\documentclass[10pt,a4paper,twoside,fancy,bg=print,twocolumn,openany,nodeprecatedcode]{dndbook}
\usepackage[utf8]{inputenc}
\usepackage{tabulary}
\usepackage[normalem]{ulem}
\usepackage{color,soul}
\usepackage{float}
\usepackage{caption}
\usepackage{wrapfig}
\usepackage{enumitem}
\usepackage[hidelinks]{hyperref}
\raggedbottom

\extrafloats{2000}
\definecolor {mytable} {HTML} {f4edd8}
\definecolor {mycomment} {HTML} {ffffff}
\definecolor {mysidebar} {HTML} {d0cfc2}
\definecolor {myreadaloud} {HTML} {d9e8f1}
\definecolor {myreadaloudborder} {HTML} {5a6173}
\definecolor {mypagenum} {HTML} {2a2d2e}
\DndSetComplexThemeColor[mytable][mycomment][mysidebar][myreadaloud][myreadaloudborder][titlered][titlegold][titlegold]

\setlist[description]{nolistsep,listparindent=3pt,topsep=6pt}

\begin{document}
\mainmatter
\chapter*{The Netherdeep}
\DndDropCapLine{L}{eaving behind the bustling streets of} Ank'Harel and the gloom of Cael Morrow, the party now navigates the dark waters of the Netherdeep. This otherworldly tomb of Alyxian's own making is similar to a demiplane or an extradimensional space connected to the Material Plane but not part of it. In other words, the Netherdeep is its own plane of existence. It is where the Apotheon's negative emotions have festered, creating labyrinthine trenches of regret, fury, yearning, and despair.


\section{Running This Chapter}
Refamiliarize yourself with the Apotheon's back story in the introduction before running this chapter. Because this chapter leads up to the final confrontation with the Apotheon, understanding his past and his frame of mind is crucial. The ending of the adventure hinges on whether the characters can help him overcome his pain and despair.

\section{Netherdeep Overview}
The Netherdeep is an underwater dungeon at the bottom of a lightless abyss. Spells or magic items that enable characters to breathe underwater are a necessity for exploring the Netherdeep.

The following sections describe the Netherdeep's recurring features, the Fragments of Suffering that the characters must acquire before they can face the Apotheon, and the main threats the characters must overcome as they make their way toward the Apotheon's prison.

\subsection{Netherdeep Features}
The following sections summarize important or recurring features of the Netherdeep.

\subsubsection{Ceilings and Floors}
Tunnels have 10-foot-high ceilings, while a chamber or grotto has no ceiling unless its description states otherwise. If an area has no ceiling, assume that its walls are about 100 feet tall, and beyond them is a dark, underwater abyss. Creatures that swim upward never reach the surface and can only return downward to the area they left.

Unless otherwise noted, all tunnels have irregular floors that are at the same elevation as the floors of the chambers they adjoin.

\subsubsection{Fragments of Suffering}
To enter the Heart of Despair, where the Apotheon is imprisoned, a creature must have absorbed at least one Fragment of Suffering—a tangible representation of an emotion felt by the Apotheon. Nine fragments exist in the Netherdeep. A Fragment of Suffering looks like an immobile mote of light about the size of a fig. The Fragments of Suffering table shows the names and locations of the fragments.

\begin{DndTable}[header=Fragments of Suffering]{XX}

Fragment Name & Area\\
Fragment of Despondence & N4a\\
Fragment of Attachment & N5a\\
Fragment of Pity & N7a\\
Fragment of Rancor & N14a\\
Fragment of Abhorrence & N16a\\
Fragment of Melancholy & N18\\
Fragment of Intransigence & N22a\\
Fragment of Deception & N24a\\
Fragment of Loathing & N25\\
\end{DndTable}

\subparagraph{Absorbing a Fragment}
As an action, a creature can touch a fragment it can see and thereby absorb the fragment into its body. If a character absorbs a fragment, give that character's player the card that corresponds to the fragment (see appendix D). The fragment disappears and becomes part of the creature's soul until the creature dies, transfers the fragment to another willing creature (see "Releasing a Fragment" below), or enters the Heart of Despair.

A creature can absorb up to three Fragments of Suffering. If a creature that carries three fragments tries to absorb a fourth, the attempt fails automatically.

\subparagraph{Releasing a Fragment}
A creature can use an action to transfer one fragment it is carrying to another creature by touching the fragment's intended recipient. When a fragment leaves a creature, that creature loses the benefit and drawback conferred by the fragment. If a character initiates a transfer, have that character's player give the corresponding card to the new recipient.

If a creature that is carrying a fragment dies or enters the Heart of Despair, the fragment immediately leaves the creature and returns to its original location (noted in the Fragments of Suffering table).

\subsubsection{Light}
Locations in the Netherdeep are either engulfed in darkness or dimly lit by ruidium deposits (see "Ruidium" below). All non-ruidium light sources in the Netherdeep shed light half as far as usual. For example, a \textit{light} spell gives off bright light in a 10-foot radius and dim light for an additional 10 feet.

\subsubsection{Ruidium}
Ruidium is plentiful throughout the Netherdeep. Brief contact with the vermilion mineral isn't harmful, but contact lasting more than a few seconds can lead to ruidium corruption (see "Ruidium Corruption" in the introduction). The save DC to resist the corruption is 20.

\subsubsection{Secret Doors}
Secret doors in the Netherdeep are made of the same ruidium-veined stone as the walls surrounding them. In any location that has a secret door, the text explains how the door can be found and opened.

A character who spends at least 10 minutes searching an area for secret doors can make a DC 20 Intelligence (Investigation) check at the end of that search. If the check succeeds and there's a secret door to be found, the character finds it. Secret doors in the Netherdeep are unlocked unless otherwise noted, and they open easily despite their great weight.

When the characters pass through a secret door, it closes behind them unless it is held or wedged open. A character who has passed through a secret door once can always find and open it again from either side of the door.

A secret door has AC 17, 50 hit points, and immunity to poison and psychic damage. Any creature within 10 feet of a secret door that deals damage to it must make a DC 15 Constitution saving throw, taking 22 (4d10) psychic damage on a failed save, or half as much damage on a successful one.

\subsubsection{Water Pressure}
The Netherdeep is a supernatural abyss, and the water pressure it exerts on creatures within it is beyond anything most mortals can survive. Creatures not native to the Netherdeep take 35 (10d6) force damage at the end of every minute they spend within the Netherdeep, unless they are wearing or wielding an item infused with ruidium or have ruidium embedded in their bodies. See appendix B for more information on ruidium items.

\subsection{Apotheon Lore}
In this chapter, the characters experience fleeting visions that shed light on Alyxian's past. The characters experience these visions as a group regardless of how close they are to one another. The visions come to them in the order given in the Apotheon Lore table.

A vision occurs whenever one of the following events takes place:


\begin{itemize}
\item The characters, as a group, finish a short or long rest in the Netherdeep.
\item A character absorbs a Fragment of Suffering.
\item A character falls unconscious in the Netherdeep.
\item Some other event in the Netherdeep triggers a vision, as noted in the text.
\end{itemize}

After the characters receive all twenty visions, they can piece together the Apotheon's whole story.


{
  \par \vspace { 9pt plus 3pt minus 3pt } \noindent
  \DndFontTableTitle{Apotheon Lore} \nopagebreak
}
\begin{dndlongtable}[c p{0.8\linewidth}]

Order & Vision\\
1st & A child is born, eerily silent, and the midwife refuses to hold him. "Look. Even he knows," she rasps to his mother as she clutches the baby, trembling. "His little soul prepares for the dread omens that will haunt him the rest of his life."\\
2nd & Two parents place a baby in front of a priest of Pelor and beg a blessing. The priest makes a gesture over the infant's head, then recoils. "There will be no help for him here. Don't be foolish. There is no cure for the mark of the red moon."\\
3rd & A baker hides a mussed, dirty boy behind the counter in his shop, slipping him a tart when no one is looking. Outside, children crow and call for the omen-touched boy, laughing, as a patron furtively mutters to the baker, "He's full of Betrayer's blood, I heard."\\
4th & A teenager sleeps fitfully under a vermilion moon, its glow washing him in faint red tones. He doesn't rest easily until a cloud drifts in front of the moon, masking its sight from his slumber.\\
5th & A young man stands atop a mountain and beholds a green Xhorhasian vista. He holds a handful of lush flowers. The vision shifts, and now he looks out over the same landscape blighted by the Calamity—the greenery is burned to the ground, corpses litter the hillside, and he grips a sword dripping blood.\\
6th & A soldier crouches over her dead commander in a plaza filled with the bodies of people and monsters. The weight of responsibility settles uncomfortably on her. "Harbinger!" she spits at Alyxian. "War-bringer! Omen! We sheltered you!"\\
7th & A guttural roar leaves the throat of the young warrior as he cleaves through the foes in front of him. His rapid breathing makes allied soldiers around him skittish. "Omen of malice," one hisses to another. "Don't get on his bad side."\\
8th & The warrior, now an adult, kneels before an altar in a ruined temple of Avandra the Change Bringer. "Fight for the freedom of others," he whispers in reverence. His hand, scarred and shaking, grips the hilt of his sword. "Luck favors the daring," he says with determination. "The Change Bringer doesn't bless cowards."\\
9th & The warrior holds his hands up to the sky beneath the sheltering, silvery glow of the moon Catha, and a drop of moonlight falls into his cupped hands. It turns into a swirl of colors, then takes the form of a small golden amulet. The warrior, awestruck, whispers a prayer as he presses the \textit{Jewel of Three Prayers} to his chest.\\
10th & Bleeding heavily, the warrior limps through a burning battlefield toward a wailing child. She screams when she sees him; he flinches, but he picks her up, then carries her toward the remnants of civilization in the distance.\\
11th & Close to sleep, huddled beneath a tree, and cloaked in a ratty blanket, the warrior runs his finger along the edge of the \textit{Jewel of Three Prayers}, finding some comfort in its familiar shape.\\
12th & During a thunderstorm, the warrior bellows, "Why? Avandra, this is not the change I seek! Tell me why I walk these steps!" The roar of the storm is all he hears in response. He rakes his hands through his hair. "Tell me!"\\
13th & The warrior kneels in front of two moonlit graves bearing the names of his parents. He places incense between them, resting it on a simple box decorated with wedding bells. The box bears the inscription "For those whose love never wavered." Around his neck, he wears the \textit{Jewel of Three Prayers}, which carries the blessings of Avandra the Change Bringer and Sehanine the Moon Weaver.\\
14th & The warrior's sword cleaves a beast in two. Its corpse falls onto the blood-soaked ground, joining many other victims. The warrior, wounded, stands above the carnage. Two people lean out of nearby houses—then four. Ten. Dozens. They run to him and cry out, some of them sobbing. "The Apotheon!" "They said you were coming!" "We knew it!"\\
15th & The warrior kneels before an elf monarch. "I've heard the most curious thing about you, Alyxian," the elf drawls. "You keep refusing the titles I offer you. The land, the riches, the servants." The warrior whispers in reply, "The power to change fate. That's all I have ever asked."\\
16th & The light is warm and dim in a tavern, where the weary warrior listens to a bard's song: "Don't kiss underneath the red moon, nor promise yourself to sweet eyes! They'll leave for another one's home too soon, or string up your heart in their lies!" He pays for his drink and leaves, despite the early hour.\\
17th & The warrior stands at the door of a ruined house and peers inside at a huddled, frightened family. They draw back when they see him, but he pulls off his cloak and drapes it over a shivering youth. "Go to the city," he rasps. "It is a bastion of hope."\\
18th & The night air is so thick with the metallic smell of blood that the warrior can sense the massacre before he sees it. He crests a hill and stares out over a city in southern Marquet, trembling with a mixture of anticipation and dread as the light of the moon Ruidus bathes the bodies in red.\\
19th & Standing in front of a mirror, the warrior considers the toll the Calamity has taken on him—the scars, the exhaustion etched into his features, all the result of the endless clash between faith and fear. He clasps the \textit{Jewel of Three Prayers}, now complete with the addition of the blessing of Corellon the Arch Heart. "Three divine mantles," he whispers. "Let us all be free of this duty."\\
20th & A warrior stands atop a temple, beneath the towering form of a one-eyed titan, as if intending to test his human-sized spear against the Betrayer God's massive weapon. "I am cloaked in moonlight," he whispers, alight with a silvery aura. "My heart is steeled by deep magic." Brilliant green suffuses the silver of the aura. The demigod faces down Gruumsh's apocalyptic spear. "And armed with the power to change fate!" Golden light flares, just before the god's spear thrusts downward and all goes dark.\\
\end{dndlongtable}

\subsection{Rivals in the Netherdeep}
In this chapter, the characters' rivals use their tier 3 stat blocks (see appendix A).

Unless they were routed or killed earlier, the rivals are also trying to reach the Apotheon. Depending on how earlier events played out, the rivals are either ahead of the characters or behind them. Regardless, you must choose when and where the characters first meet the rivals in the Netherdeep. If you're not sure where this encounter should occur, choose one of the following options:


\begin{itemize}
\item The rivals lie in wait for the characters in area N1. Their goal is to force the characters to flee while warning them never to come back.
\item The two groups meet in area N26. If the rivals are hostile toward the characters, this is their final stand—a fight to the death. If the rivals are friendly or indifferent toward the characters, they fight as hard as they can but are willing to retreat if the battle turns against them.
\end{itemize}

If the characters are trailing the rivals, you can drive home the point that "the race is on" by having the rivals leave some trace of their recent passage through an area. Determine what they left behind by rolling on the Evidence of Rivals table. Don't use this table more than once in any particular area.

\begin{DndTable}[header=Evidence of Rivals]{cX}

d4 & Evidence\\
1 & The corpses of one or more Netherdeep creatures killed by the rivals\\
2 & A broken piece of armor, twisted and bent out of shape\\
3 & A taunting message carved into a wall, telling the characters that they have fallen behind\\
4 & A lost or discarded weapon, piece of ammunition, potion bottle, or trinket\\
\end{DndTable}

\subsubsection{Race to the Bottom}
If the characters have a friendly relationship with the rival party, the Apotheon tries to ruin their friendship and put the two groups at odds. Here are some ways in which the Apotheon's interference can manifest:


\begin{itemize}
\item The characters discover the rivals tending their wounds after a skirmish with one or more Netherdeep creatures. Ayo warns the characters of a possible threat or a treasure in the next area, but the Apotheon twists locations around, making it seem as if Ayo has lied to them. (Essentially, the characters are moved to another part of the dungeon without their knowing it. If they head back the way they came, they end up in the area they previously visited and are not displaced.)
\item The Apotheon whispers to Galsariad that the characters know secrets of the Netherdeep that they refuse to share. Galsariad bitterly accuses the characters of knowing more than they've let on, or he convinces his companions that one of the characters should be captured and interrogated.
\item When the characters are close to discovering a Fragment of Suffering, the Apotheon recreates the screams of the rival party to lure the characters to another location where danger awaits.
\end{itemize}

If your goal is to turn the rivals into hated adversaries, consider using one or more of the following tactics:


\begin{itemize}
\item After the characters overcome a challenge to acquire a Fragment of Suffering, they find that the rivals have already claimed the fragment and left behind some kind of calling card.
\item When the characters arrive at a new location, they find the rivals exiting it, being chased by several creatures snapping hungrily at them. Ayo spots the characters, winks maliciously at them, and casts \textit{pass without trace} on her party, causing the monsters to lose sight of their quarry and turn their attention to the newcomers.
\item Galsariad uses his \textit{telekinesis} spell to trigger the partial collapse of a tunnel, blocking the characters' path until they spend 10 minutes clearing the rubble.
\end{itemize}

\section{Netherdeep Regions}
The Netherdeep has three interconnected regions:

Grottoes of Regret (Areas N1–N9). Characters who pass through the rift in Cael Morrow arrive in this region, which is shaped by the Apotheon's sorrows of a life unfulfilled. The Apotheon's memories and regrets manifest here in various ways.

Vents of Fury (Areas N10–N18). This region is dimly lit by fiery vents that belch the Apotheon's wrath. It is filled with violent creatures and shards of the Apotheon's persona.

Chasm of Yearning (Areas N19–N25). These twisting corridors of illusion represent the Apotheon's desire for freedom. The region's creatures and architecture try to disorient intruders.

The three regions are linked by a central hub (area N26) that contains the Heart of Despair, a cocoon of pain and misery that the Apotheon has constructed for himself. Chapter 7 describes the interior of his self-made prison in more detail. Mortal creatures can enter the Heart of Despair only by taking his pain into their souls (see "Fragments of Suffering" earlier in the chapter), thus becoming more like the Apotheon.

\section{Grottoes of Regret}
\DndQuote%
{}
{''I lived in service of others, selflessly, for I knew no other way. Now, forgotten, hanging between life and death, I can think of naught but regret for a life haunted and unlived.''}
{Alyxian the Apotheon}
Rough walls of ruidium-veined stone are interrupted by grand friezes of warriors in their phalanxes, standing proudly against monstrous hordes and towering titans. Wispy figments of memory linger before the carved walls, lost in endless, wistful dreams. Crevasses that turn into winding tunnels connect the grottoes, which the Apotheon uses as vessels for his sorrowful memories.

Walls and floors in the Grottoes of Regret are festooned with foliage as varied in color as coral—often ghostly white, but eerie colors also pierce the softer, lush landscape, interrupted by glowing crimson veins of ruidium.

\subsection{Living Memories}
When they enter certain chambers (areas N2, N3, N4, N5, and N7), the characters are transported to new sites that are not underwater—places where they can interact with characters and locations from Alyxian's past without actually changing the past. The Netherdeep is never fully isolated from any of these quasi-real locations. Ruidium formations, sea grass, and white coral appear on the outskirts of these places, and puddles of water lie on the ground—indications that the Netherdeep is not far away.

Like living dioramas, these scenarios consist of more than just the characters witnessing events of Alyxian's life that he regrets; the characters are experiencing these moments from Alyxian's perspective and are given a chance to change what he did or might have done—not to change the past, but to help the Apotheon confront the choices he made. Some players will figure out this situation quickly, but others might need additional hints. For example, figures dressed in old styles of clothing and using obscure idioms might suggest to the characters that they've been transported not only to another place, but another time.

\begin{DndSidebar}[float=!b]{Emotional Torment}
Various effects in the Netherdeep can twist and amplify the emotions of visitors, causing them to act irrationally. They might be tempted to give in to their vices, face up to their flaws, or suffer in the same ways that Alyxian has. This effect doesn't use the rules for madness; rather, the visitors are being tormented by unresolved suffering, and this distinction will make for a more rewarding outcome if the characters help Alyxian reach catharsis in the final confrontation.



\end{DndSidebar}

\subsection{Roleplaying the Apotheon}
The Apotheon's voice accompanies the characters while they explore the Grottoes of Regret. Within the grottoes, the Apotheon's tone is sorrowful, even reluctant, but helpful. He has spent many years wallowing in these torments and clings to them—they are familiar comforts at the same time that they wound him.

The Apotheon is never more than a disembodied voice in the Grottoes of Regret, but he's more willing to talk to the characters in this region than anywhere else in the Netherdeep. When the characters enter certain areas for the first time, the Apotheon pulls them out of the Netherdeep and drops them into scenes that hearken back to past events that continue to haunt him. His reason for involving the characters in these regretful episodes comes from one or more of these motivations:


\begin{itemize}
\item Alyxian is trying to see some of his feared memories through the eyes of others to gain perspective on his current imprisonment.
\item Instinctively and without malice, Alyxian is trying to elicit empathy from the party.
\item Alyxian is fixated on his own tragedy and can't help but pull others into it.
\end{itemize}

You can choose one motivation for all the areas in the Grottoes of Regret, or you can choose a different option for each area, but do your best to portray the Apotheon sympathetically in this region.

\subsection{Grottoes of Regret Locations}
The Grottoes of Regret encompass areas N1 through N9 on the Netherdeep map.

\subsubsection{N1: Antechamber}
Read or paraphrase the following when the characters first pass through the rift in Cael Morrow and enter the Netherdeep:

\begin{DndReadAloud}{}
You burst from the planar rift and are plunged into deep, cold water.

\end{DndReadAloud}

Characters who emerge from the rift are immediately subjected to the pressure of the deep (see "Water Pressure" earlier in the chapter). Describe their surroundings as follows:

\begin{DndReadAloud}{}
You are floating ten feet above the floor of a cavernous pit, the rift at your back. Two enormous statues have settled into the seafloor at awkward angles. Each statue is missing its head, and one of the statues has an outstretched hand that looks like it might grab anything that gets too close to it. Covering the statues and clinging to the walls are crystalline tendrils that shed dim, eerie red light. A ten-foot-high tunnel at floor level leads out of this sunken chamber, which has patches of sea grass growing everywhere.

\end{DndReadAloud}

Swimming in the dark water 20 feet above the characters (30 feet above the floor) are two devourers and a \textbf{slithering bloodfin} (see appendix A for their stat blocks) that act as sentries. Characters who find themselves outmatched by the monsters or unable to cope with the water pressure can retreat through the open rift and return to Cael Morrow.

\subparagraph{Statues}
The ruidium-covered statues are depictions of Alyxian the Apotheon as a noble warrior in his prime, but the missing heads make identification difficult. The heads are nowhere to be found, symbolizing Alyxian's loss of self.

\subsubsection{N2: Child's Lament}
This water-filled grotto has tunnels to the south, northeast, and northwest. These short passages lead to areas N1, N3, and N4, respectively.

The first time the characters enter this area, read:

\begin{DndReadAloud}{}
The water shimmers as you are transported out of the undersea gloom and into a child's room. Simple comforts are scattered around the room's perimeter: a teddy bear overgrown with ruidium, a cradle lined with ruidium spikes, and a rocking horse with tendrils of ruidium crawling up one side of it. Simple toys are scattered on the floor, with more of them on small wooden shelves. All of them are laced with ruidium.

Sunlight shines into the room through two small, high windows. You hear muffled voices coming from outside and an occasional dull thud.

\end{DndReadAloud}

The characters are transported to Alyxian's childhood home, a modest wooden dwelling with three rooms: a 20-foot-square room containing a kitchen and a living area, with doors leading to two 10-foot-square bedrooms. The characters appear in Alyxian's bedroom. A door off the kitchen leads outside.

\subparagraph{Alyxian's Toys}
A character who examines the rocking horse notices the name "Alyxian" carved on its flank in Common. Characters who want to keep one or more of Alyxian's toys can do so, and those toys remain with the characters after this vision ends. The first time a character picks up one of these toys, the Apotheon whispers, "Such a thing no longer holds a special place in my heart, but it might be worth something to you." These toys prove useful in area N6. Any character who handles a toy without gloves or similar protective gear might be subject to ruidium corruption, at your discretion (see "Ruidium" in the introduction).

\subparagraph{Bedroom Windows}
Alyxian's bedroom has a 10-foot-high ceiling and two windows. Each window is a 2-foot-square opening whose bottom is 7 feet above the floor. A Medium character can peer out a window easily enough by standing on the back of the rocking horse.

\subparagraph{Outdoor Scene}
Any character who looks out a window or exits the cottage sees two humans being attacked by a band of six people armed with farm implements. The two folks being accosted are representations of Alyxian's parents. His father has a kindly face and a bushy brown beard. His mother has a worry-lined, scarred face. Both have olive skin and wear simple, dirt-stained garments made of a single piece of cloth. Their attackers are shadowy, faceless facsimiles of people draped in worn-out clothes. Red eyes burn from beneath their wide-brimmed hats. These shadowy assailants use the \textbf{berserker} stat block, but their weapon attacks deal psychic damage. Their vague forms are due to the fact that Alyxian's memory of them has faded.

As soon as one or more characters see what's going on outside, Alyxian's voice whispers to them:

\begin{DndReadAloud}{}
"You have found me, my friends. Witness my earliest memory. The pain of a child born under the ill-fated sign of Ruidus. We were turned on by friends and kinfolk as war raged around us. Many people suffered from my presence, my parents most of all. My story begins with them. Their story dies with me."

With a loud crack, one of the shadowy figures strikes Alyxian's father, sending him to the ground.

\end{DndReadAloud}

Characters who exit the cottage not only witness the unfolding scene but also attract the attention of the assailants. Four of them break off to attack the characters while the remaining two continue to strike Alyxian's parents on the ground. These two assailants turn to fight the characters only when the four who attacked the characters are dead.

If the characters confront the faceless assailants, have the players roll initiative for their characters, Make one initiative roll for all the assailants.

\subparagraph{Placating the Apotheon's Regret}
To placate this regret of Alyxian's, the characters must try to rescue his parents. If the assailants haven't all been defeated by the start of their fourth turn, Alyxian's father dies from his wounds, and his mother is left for dead. Administering magical healing to the father delays his death at the hands of the assailants by 1 round.

As they attack, the assailants hiss justifications for why young Alyxian and his parents should die. Examples of what they might say:


\begin{itemize}
\item "You have brought this ruin upon us!"
\item "Monsters birth monsters!"
\item "No good ever comes from sheltering a child of Ruidus!"
\end{itemize}

If Alyxian's parents are saved, they speak to the characters as if they were speaking to Alyxian:

\begin{DndReadAloud}{}
"Get back inside, Alyx. Everything's all right. They just don't understand."

\end{DndReadAloud}

If one or more characters offer words of comfort to Alyxian's parents after rescuing them from the assailants, this interaction has a positive effect on the final battle against Alyxian (see "Emotional Healing" in chapter 7).

\subparagraph{Ending the Vision}
Once this episode has played out, the characters are transported back to the underwater grotto along with anything they are wearing or carrying, including any toys they gathered from Alyxian's bedroom. Once they are back in the water-filled chamber, Alyxian whispers the following to them in their minds:

\begin{DndReadAloud}{}
"What I understood from a young age is this: because of me, they would always suffer."

\end{DndReadAloud}

\subparagraph{Subsequent Visits}
The vision involving Alyxian's childhood home is not repeated if the characters leave this area and return later. On future visits, this grotto is dark and empty except for a few patches of sea grass.

\subsubsection{N3: Unforgotten Fallacy}
This water-filled grotto is connected to area N2 by a short tunnel. A longer tunnel in the north wall bends west and passes under area N4a on its way to area N5.

The first time the characters enter this area, read:

\begin{DndReadAloud}{}
The walls of this chamber shimmer as you enter another memory. Fragmentary recollections fill your mind: you recall Alyxian as a daring teenage boy who wanted to take part in a village hunt. The vermilion moon was shining full that night, and you remember Alyxian telling himself that he was strong enough to overcome it. You remember him donning a cloak and joining the hunt in secret.

Your awareness shifts to a clearing in a moonlit forest. Its beauty is belied by the thick scent of blood and the sound of frantic shouts.

\end{DndReadAloud}

The characters are whisked away to a forest where a hunt has just gone terribly wrong. The misfortune that has followed Alyxian his entire life has now doomed his hunting party. One of the hunt captains, a neutral \textbf{drow elite warrior} named Saqiri Yestrana, was mauled by two owlbears in the red moonlight, and another hunter caught a glimpse of the young Alyxian before he parted company with Saqiri and darted into the woods.

When the characters enter the scene, the owlbears are escaping, scattering a herd of deer in the process, and Saqiri is bleeding out. As Saqiri coughs up blood, the characters hear the following shouts of panic from the hunters:


\begin{itemize}
\item "Those beasts were going to see us through the winter! We'll be short now!"
\item "Alyxian has run off!"
\item "Saqiri's not going to make it! She needs tending!"
\end{itemize}

Alyxian whispers the following in the characters' minds after they take in the scene:

\begin{DndReadAloud}{}
"I learned the price of my selfishness early. I just wanted to—I just wanted to hunt among them, to be with them." He sighs and growls ruefully. "It doesn't matter. Saqiri died that night, and I will carry that regret until I join her."

\end{DndReadAloud}

\subparagraph{Placating the Apotheon's Regret}
To placate Alyxian's regret, the characters must hunt for meat, heal Saqiri, or both. Accomplishing one or both tasks has a positive effect on the final battle against Alyxian (see "Emotional Healing" in chapter 7).

\subparagraph{Hunting for Meat}
One or more characters can hunt the owlbears and the deer to bring back meat to the hunters' camp. At least one character must succeed on a DC 16 Wisdom (Survival) check to track the owlbears through the forest. As they explore the forest, the characters see beds of long, flowing kelp emerging from among the trees—a hint that the Netherdeep isn't far away. If one or more characters succeed on the check, they encounter the two owlbears. When an owlbear dies, its body splits open and a \textbf{swarm of sorrowfish} (see appendix A) bursts out and attacks. These sorrowfish swim through the air as if it were water. Once the swarms are defeated, the characters can take the owlbear meat—enough to sustain the village—back to the hunters' camp. If no characters succeed on the check, they don't encounter the owlbears and find only the carcasses of sickly deer that couldn't keep up with the rest of the herd. If these carcasses are returned to the camp, they provide some edible meat—enough to placate the Apotheon's regret.

\subparagraph{Healing Saqiri}
Four hunters (scouts) guard Saqiri. These hunters have shadowy forms and hazy faces, because Alyxian's memory of them has faded somewhat.

The hunters regard the characters with the same contempt they hold for Alyxian, shouting to them, "Be off with you! You're a curse on all of us!" A character can get to Saqiri's side by succeeding on a DC 14 Dexterity (Stealth) check to slip past the guards. A character can also approach the guards and convince them not to attack with a successful DC 14 Charisma (Persuasion) check. The guards attack any character who fails either check. Characters who try to force their way past the guards are also attacked.

Saqiri is unconscious and dying from her wounds. If healing restores her to 1 hit point or more, she struggles to sit up and speaks the following words to the character closest to her, as if that character were Alyxian: "Go on, boy. Take one of the carcasses to your ma. Tell her you did a good job tonight. I don't know if a curse truly follows you—but people believe it does, and that's enough for them." The characters then hear Alyxian whisper the following words: "I have watched Saqiri's death a thousand times. It never stops tearing my soul apart—but now I know what she would have said had she lived." More quietly, he adds, "It is not what I expected after so long."

\subparagraph{Ending the Vision}
Once this scene has played out, the forest fades away and the characters are returned to the underwater grotto. If they accomplished one or both tasks in the vision, a small, unlocked chest appears on the floor nearby. It contains the following rewards based on what tasks the characters completed:


\begin{itemize}
\item If the characters healed Saqiri, the chest contains fifteen tiny rubies (100 gp each).
\item If the characters brought owlbear meat back to the camp, the chest contains two 6-inch-tall, solid gold owlbear figurines (750 gp each), one on all fours and the other reared up on its hind legs. If they returned with deer meat for the camp instead, the chest contains a 3-inch-tall figurine of a deer carved out of translucent, grayish-blue chalcedony (150 gp).
\end{itemize}

\subparagraph{Subsequent Visits}
The vision of Alyxian's first hunt is not repeated if the characters leave this area and later return. On future visits, this grotto is dark and empty except for a few patches of sea grass.

\subsubsection{N4: Haunted Ashes}
This water-filled grotto has short tunnels leading to areas N2 and N5. A third tunnel leading to area N4a is hidden behind a secret door.

The first time the characters enter this area, read:

\begin{DndReadAloud}{}
The water in this chamber feels as slick as oil. A shiver moves down your spine as halos of pale red light surround you and your companions. You all sink to the bottom of the chamber and land softly as your feet touch the floor.

You are no longer underwater but standing on the edge of a small village surrounded by hills. A melee has broken out about a hundred yards away in the middle of the village. Shadowy soldiers are pillaging the settlement while villagers mount a futile effort to rebuff them. Before you can act, one of these soldiers raises his sword above his head and howls, "You ought to be kissing the ground we walk upon, you spineless cowards!" In a flash, he stabs a villager, whose death wail rings in your ears.

\end{DndReadAloud}

The characters have been transported to a sacked village in the distant past, during the time of the Calamity.

\subparagraph{Red Halos}
The halo surrounding each of the characters sheds bright light in a 10-foot radius and dim light for an additional 10 feet. These halos vanish when the characters leave this memory and are returned to the Netherdeep.

A character encircled by a red halo leaves blotches of ruidium wherever their feet touch the ground. These ruidium deposits slowly send out tendrils across the ground in the character's wake. In this vision, any successful attack a character makes with a melee or ranged weapon deals an extra 11 (2d10) psychic damage to the target.

\subparagraph{Raiders}
A group of soldiers—allies of Alyxian in the battle against Torog—have decided that since they saved this village, it is theirs to pillage. Four sword wraith warriors led by a \textbf{sword wraith commander} (see appendix A for their stat blocks) take the shadowy shape of Alyxian's crass allies. They are bullying a group of ten commoners in the middle of the village. As the characters approach, the commander, a neutral evil human in his mid-twenties named Kalagothe, looks to them as if they were Alyxian and roars, "Commander Alyxian, these people refuse to pay us tribute!"

Have the players roll initiative for their characters, then make one initiative roll for Kalagothe; the remaining sword wraiths act immediately after him in the initiative order. Roll initiative once for all the commoners, but don't waste time devising tactics for them; they cower or retreat from the soldiers on their turn.

A character can use an action to try to convince the soldiers to stand down, doing so with a successful DC 27 Charisma (Intimidation or Persuasion) check. The DC drops to 17 if Kalagothe is dead and only the sword wraith warriors need to be dealt with. Only three such attempts can be made, after which the sword wraiths are immune to such coercion.

Unless the soldiers are slain or forced to stand down, they capture and murder the commoners.

\subparagraph{Placating the Apotheon's Regret}
If the characters rescue one or more of the villagers, the Apotheon's voice tentatively whispers, "There are bright spots amid the ruin, sometimes. Beauty. Kindness." This interaction has a positive effect on the final battle against Alyxian (see "Emotional Healing" in chapter 7).

\subparagraph{Ending the Vision}
If the characters rescue at least five of the ten commoners, the Apotheon whispers to them, "You handled that better than I did. I wish I could have seen their faces happy with relief that they were saved, rather than fraught with sorrow at their losses." After Alyxian speaks these words, the vision fades and the characters are returned to the underwater grotto. As the characters reappear in the grotto, a secret door in the north wall opens (see "Secret Doors" earlier in the chapter). Beyond this door lies a short tunnel leading to a small chamber (area N4a).

\subparagraph{Subsequent Visits}
The vision of the village is not repeated if the characters leave this area and later return. On future visits, this grotto is dark and empty except for a few patches of sea grass growing near the walls.

\subsubsection{N4a: Fragment of Despondence}
\begin{DndReadAloud}{}
In the middle of this chamber stands a marble pedestal, above which floats a mote of opalescent light about the size of a fig. A bolt of crimson energy occasionally crackles across its glowing surface. Resting on the pedestal below this light is a tiny gray doll in the shape of some kind of four-legged animal.

\end{DndReadAloud}

When they enter this area, the characters hear the Apotheon's bitter words in their minds:

\begin{DndReadAloud}{}
"Yes. You let go of hope because you know nothing will ever change."

\end{DndReadAloud}

\subparagraph{Doll}
The doll on the pedestal looks like a moorbounder. It has sentimental value to Alyxian. If a character picks it up, everyone in the room hears the Apotheon whisper, "Maeska, my old friend." The moorbounder doll might prove useful to the characters in area N6.

\subparagraph{Fragment of Despondence}
The mote of light above the pedestal is the Fragment of Despondence (see "Fragments of Suffering" earlier in the chapter). When a character touches the mote, it disappears and becomes part of the character's soul. When this happens, all the characters receive a fleeting vision (see the Apotheon Lore table earlier in the chapter).

The character who is carrying the Fragment of Despondence gains the following benefit and drawback, and is aware of both:



\begin{description}
\item[Benefit.]
You are immune to the charmed condition.
\item[Drawback.]
You can't take the Help action.
\end{description}



Give the Fragment of Despondence card (see appendix D) to the player whose character is carrying the fragment.

\subparagraph{Tunnel Exits}
Two short tunnels lead away from this chamber, one toward area N4 and the other toward area N6. Each tunnel ends at a secret door (see "Secret Doors" earlier in the chapter). A creature leaving this location automatically finds either secret door.

\subsubsection{N5: Alyxian's Refuge}
This water-filled grotto has short tunnels leading to areas N4 and N6. Between them, a longer tunnel travels east, passing under area N4a on its way to area N3. A fourth tunnel leading to area N5a is hidden behind a secret door.

The first time the characters enter this area, read:

\begin{DndReadAloud}{}
As your consciousness melds with the Apotheon's, you remember a tale from the Calamity when Alyxian took shelter in the home of three women. For one night, the women fed Alyxian and gave him a bed to sleep in. Now, you stand outside their home, near which are planted three unsightly gravestones made of jagged red crystal.

In your mind, you hear Alyxian say, "They welcomed me into their home. Why can't I remember their names or faces? Please, help me remember."

\end{DndReadAloud}

The characters have been transported to another place Alyxian remembers from his troubled past. When he was a young man fated to be caught up in the Calamity, Alyxian stayed in the home of a family for one evening, and ruin followed him to their doorstep. The two women and their teenage daughter were killed when the village was sacked the day after Alyxian's departure. He learned about their deaths only several months later when another soldier told him about the village's destruction. Alyxian is desperate to remember the names of the three victims, which is why the characters have been brought here.

\subparagraph{Gravestones}
The gravestones near the ruined house are made of ruidium and bear no inscriptions. No graves lie beneath them, so no corpses are available to examine or question by using magic.

\subparagraph{Ruined House}
If the characters explore the ruins of the house, read:

\begin{DndReadAloud}{}
Corpses and smashed furniture are strewn throughout the ground floor of the house, which contains a living room and a kitchen plus a small dining room. A damaged staircase leads up to a second-floor bedroom.

\end{DndReadAloud}

A thorough search of the ground floor reveals ten corpses. They belong to the three women, five grimlocks, and two human warriors clad in bloodstained chain mail. The warriors were apparently slain trying to protect the family from the grimlocks. All these creatures look as if they died several days a day ago.

\subparagraph{Forgotten Names}
To satisfy Alyxian, the characters must learn the names of the three women: Marisa, Celeste, and Meri. Marisa was the eldest, Celeste was her wife, and Meri was their teenage daughter. Their family name is Zenthas. Marisa's name can be discovered in the kitchen, Celeste's in the living room, and Mari's in the bedroom, as described below:


\begin{description}
\item[Kitchen.]
Marisa was making a loaf of bread when the attack came. A character who spends at least 1 minute searching the kitchen can make a DC 12 Intelligence (Investigation) check. On a successful check, the character finds a folded note next to a torn bag of cinnamon spattered with blood and tied shut with a green ribbon. The note reads: "Marisa—for that little cake you like so much."

\item[Living Room.]
Celeste used the living room as a space to indulge her hobby: incense-making. Herbs, bottles, and incense sticks lie scattered across a worktable and on the floor in one corner of the room. Any character who spends at least 1 minute searching the living room can make a DC 12 Intelligence (Investigation) check. On a successful check, the character finds a book titled Scents for the Seasons lying under a chair. Written on the book's inside cover is the name "Celeste Zenthas."

\item[Bedroom.]
Meri, on the cusp of adulthood, was training to be a soldier. Any character who spends at least 1 minute searching the bedroom can make a DC 12 Intelligence (Investigation) check. On a successful check, the character finds the name "Meri" carved into the empty sheath of a shortsword.
\end{description}

\subparagraph{Ghosts of the Remembered}
When the characters learn all three names, three \textbf{ghosts} rise from the women's corpses and attack the intruders until at least one of the ghosts' names is spoken aloud. The ghosts resemble the three women, but they are not restless spirits. Rather, they are manifestations of Alyxian's rage and regret. The characters can destroy these ghosts by reducing each of them to 0 hit points.

If a character addresses a ghost by the name of the woman it resembles, that ghost's attitude toward the party changes from hostile to indifferent, and the ghost addresses each of the characters as if they were Alyxian. If a character apologizes to a ghost on Alyxian's behalf for not being around in its time of need, the ghost's attitude changes to friendly.

The following sections describe the three women and their behavior as ghosts:


\begin{description}
\item[Marisa.]
Marisa's ghost smells faintly of sugar and cinnamon. On her first turn in combat, she wails, "Our home, defiled!" If a character addresses Marisa's ghost by name and apologizes to her on Alyxian's behalf, the ghost says, "Alyxian! When they came for us, I feared they would kill you, too. It is good to hear your voice. Be well, and do not mourn. No one can harm us anymore." After speaking these words, the ghost fades away.

\item[Celeste.]
Celeste's ghost smells faintly of incense. On her first turn in combat, she snarls, "We sheltered him! We died for him!" If a character addresses Celeste's ghost by name and apologizes to her on Alyxian's behalf, the ghost's visage softens, and she says, "I would die for you, moon-cursed and moon-blessed. A thousand times, if need be." After speaking these words, the ghost fades away.

\item[Meri.]
Meri has wavy, shoulder-length hair, strong arms, and sharp eyes. The resemblance to Marisa is unmistakable. She doesn't speak unless she is addressed by name. If a character does so and then apologizes to Meri on behalf of Alyxian, she says, "I died a soldier, like my mother. What happened was neither our fault nor yours. Be well, Alyxian." After speaking these words, the ghost fades away.
\end{description}

\subparagraph{Placating the Apotheon's Regret}
Alyxian is placated if one or more ghosts are laid to rest. After speaking the ghosts' names aloud, he adds, "The Crawling King's horde destroyed them, and I was not there to help. Yet they forgave me. Was that truly all it took?"

Each ghost that vanishes on its own instead of being defeated in combat has a positive effect on the final battle against Alyxian (see "Emotional Healing" in chapter 7).

\subparagraph{Ending the Vision}
When all three ghosts are destroyed or vanish on their own, the characters are returned to the underwater grotto. A secret door opens in the northwest wall, revealing a narrow tunnel that leads to area N5a (see "Secret Doors" earlier in the chapter).

\subparagraph{Subsequent Visits}
The vision of the ruined house is not repeated if the characters leave this area and later return. On future visits, this grotto is empty except for a few patches of sea grass.

\subsubsection{N5a: Fragment of Attachment}
\begin{DndReadAloud}{}
Rising from the floor in the center of this grotto is a marble pedestal above which floats a wispy mote of light.

\end{DndReadAloud}

The first time one or more characters enter this area, all the characters hear Alyxian whisper the following words in their minds:

\begin{DndReadAloud}{}
"These regrets are all I have. The only thing that remains is to take the path I walked to get here."

\end{DndReadAloud}

\subparagraph{Fragment of Attachment}
The mote of light above the pedestal is the Fragment of Attachment (see "Fragments of Suffering" earlier in the chapter). When a character touches the mote, it disappears and becomes part of the character's soul. When this happens, all the characters receive a fleeting vision (see the Apotheon Lore table earlier in the chapter).

The character who is carrying the Fragment of Attachment gains the following benefit and drawback, and is aware of both:



\begin{description}
\item[Benefit.]
You can't be frightened while within 10 feet of an ally. If you're already frightened and move within 10 feet of an ally, the frightened condition ends on you.
\item[Drawback.]
You have disadvantage on Wisdom and Charisma saving throws while you aren't within 10 feet of an ally.
\end{description}



Give the Fragment of Attachment card (see appendix D) to the player whose character is carrying the fragment.

\subparagraph{Tunnel Exits}
Two tunnels lead from this chamber: one travels east to area N5, and the other spirals downward, passing under the eastern tunnel before opening in the wall of a wider tunnel that connects areas N10 and N11. Each tunnel ends at a secret door (see "Secret Doors" earlier in the chapter). A creature leaving this location automatically finds either secret door.

\subsubsection{N6: Loss of Innocence}
\begin{DndReadAloud}{}
This chamber is filled with hauntingly beautiful seaweed and ghostly white coral. A stone altar in the middle of the room is decorated with carvings of cyclopean skulls and grasping hands.

A barrier made of red energy covers a doorway in the north wall, beyond which you can make out another underwater chamber. Carved above the doorway in Common is the following inscription:

\textit{Apotheon, Apotheon}.

\textit{Leave your childhood yearnings behind.}

\textit{It is time to don your armor and go to war.}

\end{DndReadAloud}

\subparagraph{Altar}
A character who examines the altar and makes a successful DC 13 Intelligence (Religion) check recognizes its imagery as symbolic of the gods Gruumsh and Torog. Placing certain objects on this altar can help the characters bypass the door of red energy (see below).

\subparagraph{Door of Red Energy}
This magical barrier is impervious to damage but can be negated with a successful \textit{dispel magic} (DC 19). Any creature that touches the door must make a DC 15 Constitution saving throw. On a failed saving throw, the creature takes 35 (10d6) force damage and is pushed 10 feet away from the door. On a successful save, the creature takes half as much damage and isn't pushed.

If the characters can't dispel the door, they can place a sentimental item from Alyxian's childhood on the altar ("Leave your childhood yearnings behind"), which causes the item to disappear and dispels the door for 8 hours. Items that qualify include a toy taken from Alyxian's childhood bedroom (area N2) and the moorbounder doll found in area N4a. A character can substitute a trinket or toy from their own childhood, with the same effect. Any item that disappears from the altar can't be retrieved by any means short of a \textit{wish} spell.

In addition, a character who dons a suit of armor in this chamber or in area N7 can pass through the doorway ("It is time to don your armor and go to war"). A character who entered either chamber already wearing armor must first take off the armor and then put it on again before they can pass safely through the doorway.

\subsubsection{N7: Battle at the Betrayers' Rise}
When the characters enter this area for the first time, describe it as follows:

\begin{DndReadAloud}{}
The roof of this underwater sepulcher has three naturally formed domes, each one reaching its apex about thirty feet above the floor. Veins of ruidium cling to the walls, forming a lattice that casts red light on the clumps of sea grass growing along the room's perimeter.

\end{DndReadAloud}

The characters have only a few seconds to observe their surroundings before they are whisked away to another time and place:

\begin{DndReadAloud}{}
The underwater gloom gives way to a scene showing a parched battlefield under a bleak, ashen sky. All around you, shadowy figures clash swords, and the air is filled with battle cries and the drums of war.

Familiar-looking black towers in the distance provide clues to where and when you are: at the gate of the Betrayers' Rise during one of the fiercest battles of the Calamity. In this conflict, Alyxian fought the Crawling King's horde and scored a hard-won victory. In this version of the scene, it seems, the battle is yours to win or lose.

\end{DndReadAloud}

The characters are experiencing the Apotheon's memories of a terrifying battle. Most of the battlefield is occupied by clashing armies—Alyxian's allies versus the monstrous hordes of Torog the Crawling King. These creatures have shadowy forms and hazy faces, reflecting the fact that Alyxian's memory of them has faded somewhat, and most of them ignore the characters.

\subparagraph{Battle Rules}
As the battle rages around them, the characters are about to be faced with three waves of foes (described below) and have little choice but to defend themselves. Have the players designate one character to act as the party's leader for the duration of the battle. (It might be the character with the highest Charisma score or a character who customarily assumes a leadership role in the group.) This character remains the designated leader for the entire battle, even if they are killed or unable to act. Once the leader is determined, have the players roll initiative for their characters while you roll initiative for the hostile creatures in each wave as they enter the fray. The characters don't reroll initiative when the second and third waves appear. At no point during this battle can the characters take a short or long rest.

The characters can move anywhere they want on the battlefield, but a character who moves farther than 30 feet from the party leader or starts their turn in such a location takes 16 (3d10) psychic damage from the chaos around them, even if the leader is dead or incapacitated. A character can take this damage no more than once per turn.

\subparagraph{First Wave}
This wave consists of one \textbf{sword wraith warrior} (see appendix A) per character, up to a maximum of five. These grim knights represent the evil soldiers of the Crawling King. As the battle is joined, Alyxian chimes in:

\begin{DndReadAloud}{}
"After so long, their faces have faded into the haze of the past. I try to remember them—not my enemies, but the soldiers who fought alongside me. The least I could do was wade into the heart of the battle, where the suffering was already the worst, and take it all into me."

\end{DndReadAloud}

Alyxian's words should make it clear that the only way to put this battle to rest is to conquer wave after wave of foes.

\subparagraph{Second Wave}
After the first wave is defeated, the characters have 2 rounds to heal and prepare for the second wave: two \textbf{moorbounders} (see appendix A) accompanied by five \textbf{sword wraith warriors}. Before the characters can do anything to prevent it, the sword wraiths merge into one creature—a \textbf{slithering bloodfin} that swims through the air as easily as it normally swims through water. The howling faces of the warriors are visible on the bloodfin's flesh. The bloodfin is not what the Apotheon sought to conjure in this moment; rather, it represents Alyxian's inability to keep the Netherdeep from intruding upon his memories.

When the second wave is defeated, Alyxian's allies roar in unison, and Alyxian chimes in:

\begin{DndReadAloud}{}
"You cut them down, endlessly, a scythe in the hands of a reaper. It doesn't matter how many you kill. You water the lands with blood, and your allies die anyway."

\end{DndReadAloud}

\subparagraph{Third Wave}
After the second wave is defeated, the characters have 1 round to heal and prepare for the final wave: a \textbf{sword wraith commander} and four \textbf{sword wraith warriors} on the back of a \textbf{horizonback tortoise}. Alyxian shares his thoughts as the tortoise lumbers toward the party:

\begin{DndReadAloud}{}
"To break the spine of the Crawling King's horde, I knew I had to slay both the commander and the tortoise. I don't remember which one I killed first."

\end{DndReadAloud}

When the commander drops to 0 hit points, it lets out a terrible scream, and any remaining sword wraith warriors flee on their next turn. When the tortoise drops to 0 hit points, it falls prone and lets out one last shudder before expiring. The characters are victorious when both the commander and the tortoise are slain.

\subparagraph{Placating the Apotheon's Regret}
When all three waves of foes are defeated, Alyxian's allies cheer as the enemy is routed. A \textbf{veteran} covered in blood limps toward the characters and says, "You won the day, brave Alyxian. My wounds burn, but if you command it, I will follow you to the end." Alyxian whispers an answer that only the characters can hear: "It never ends." The wounded soldier, named Talmyth, has 1 hit point. If a character heals him, that act has a positive effect on the final battle against Alyxian (see "Emotional Healing" in chapter 7).

\subparagraph{Secret Doors}
Once the battle has ended and the scene with Talmyth has played out, the characters are returned to the underwater grotto. At the same time, two secret doors open in the north wall (see "Secret Doors" earlier in the chapter). Behind one secret door is a narrow, dimly lit tunnel that connects to area N7a; behind the other is area N8.

\subparagraph{Subsequent Visits}
The vision of the battle at the Betrayers' Rise is not repeated if the characters leave this area and later return. On future visits, this grotto is dark and empty except for the patches of sea grass growing near the walls.

\subsubsection{N7a: Fragment of Pity}
\begin{DndReadAloud}{}
Rising from the floor in the center of this grotto is a marble pedestal, above which floats a wispy mote of light. At the foot of the pedestal, half buried in the sandy floor, is a small wooden chest.

\end{DndReadAloud}

\subparagraph{Fragment of Pity}
The mote of light above the pedestal is the Fragment of Pity (see "Fragments of Suffering" earlier in the chapter). When a character touches the mote, it disappears and becomes part of the character's soul. When this happens, all the characters receive a fleeting vision (see the Apotheon Lore table earlier in the chapter).

The character who is carrying the Fragment of Pity gains the following benefit and drawback, and is aware of both:



\begin{description}
\item[Benefit.]
Each time you spend a Hit Die to regain hit points, you regain additional hit points equal to your proficiency bonus.
\item[Drawback.]
You have disadvantage on death saving throws.
\end{description}





Give the Fragment of Pity card (see appendix D) to the player whose character is carrying the fragment.

\subparagraph{Treasure}
The wooden chest is closed and unlocked. It contains four stoppered Potion of Greater Healing in red crystal vials and a fully charged \textit{wand of secrets}. If the rivals reach this area before the characters do, Dermot takes the magic items and doesn't bother closing the chest after emptying it.

\subparagraph{Tunnel Exits}
Two tunnels, one to the north and one to the south, lead away from this chamber. The south tunnel heads to area N7. The north tunnel is longer, descends almost vertically to area N16, and is dimly lit by red light emanating from that location. Each tunnel ends at a secret door (see "Secret Doors" earlier in the chapter). A creature leaving this location automatically finds either secret door.

\subsubsection{N8: Ruiner's Spear}
When the characters peer into the grotto, it appears to be devoid of inhabitants or noteworthy features. When they enter the chamber, the secret door to area N7 closes behind them unless it is being held or propped open. At the same time, a dark form rises from the floor:

\begin{DndReadAloud}{}
A shadow rises from the sandy floor and assumes a giant form. Its single, glowing red eye remains fixed on you as it rises to its full height of one hundred feet. An enormous spectral spear is gripped in its hands.

\end{DndReadAloud}

This chamber and its challenge represent the last great sacrifice of the Apotheon. The spectral colossus is an image of Gruumsh the Ruiner. This image can't be fought or harmed. Read the following as the Betrayer God brings his spear down upon the characters:

\begin{DndReadAloud}{}
The water swirls around the spear as it plunges down toward you.

\end{DndReadAloud}

Every character in the grotto and any character who enters the grotto while the manifestation of the Betrayer God is present falls prone on the floor and has their speed reduced to 0 by the force of the Betrayer God's spear, as though the character were being impaled by it. No force, magical or otherwise, can move a character from the space in which they fall prone.

One you know which characters are prone, have the players roll initiative. Any character who starts their turn in the grotto must succeed on a saving throw. Have each player describe how their character tries to resist the pressure of the Betrayer God's spear on their turn, then use the ability that most closely matches that description for the saving throw:


\begin{description}
\item[Strength.]
The character tries to use their physical prowess to push back against the god's attack.
\item[Dexterity.]
The character contorts their body to avoid the full force of the god's attack.
\item[Constitution.]
The character braces themself against the pain of the attack, trying to endure it.
\item[Intelligence.]
The character tries to refute the attack as unreal.
\item[Wisdom.]
The character tries to combat the psychic power of the Netherdeep by concentrating on their own joyful memories.
\item[Charisma.]
The character tries to use sheer force of will to counteract the pressure.
\end{description}

The saving throw DC is 15 to start and increases by 2 each time the initiative count reaches 0, making it harder for characters to resist the spear as time goes on. On a failed saving throw, a character takes 14 (4d6) psychic damage, and their mind is filled with the voices of lost loved ones telling them to stop resisting and surrender to the attack instead. On a successful saving throw, the character takes half as much damage and hears Alyxian's voice in their mind, pondering why he gave his life to save others when they wouldn't have done the same thing for him. The effect ends on any character who either is reduced to 0 hit points by the damage or succeeds on the saving throw three times. Such a character no longer feels the spear or sees the Betrayer God pinning them down with his mighty spear (though a character reduced to 0 hit points is dying and either needs to stabilize or be healed to avoid death).

Alyxian whispers in the mind of the first character to succeed on three saves, saying, "Perhaps you will be remembered. Now, what about your companions?" That character can take the Help action to speak words of encouragement to other characters still trapped by the spear, granting them advantage on the next saving throw they make to resist it.

When all the characters who were pinned by the Betrayer God's spear end the effect on themselves, the challenge is over, and the vision ends. Read:

\begin{DndReadAloud}{}
You have felt the Betrayer God's fury, as Alyxian did. The towering image of the one-eyed god is gone, and the chamber stands empty once more. Part of the north wall crumbles away, creating a jagged, open doorway between this grotto and another one.

\end{DndReadAloud}

The doorway leads to area N9.

\subsubsection{N9: Hall of the Portal}
The characters can enter this chamber by stepping through the open doorway in the south wall after surviving the attack in area N8. Hidden in the north wall is a secret door that leads to area N26 (see "Secret Doors" earlier in the chapter).

When the characters enter, read:

\begin{DndReadAloud}{}
An oval portal with glowing edges appears in the middle of the room, floating between the hands of a giant statue depicting Alyxian in armor. Suspended vertically in the water, the portal is ten feet tall and five feet wide. Through it, you can see the giant statues that occupy the first chamber of the Netherdeep that you entered.

"You have endured," whispers Alyxian's voice. "Take this opportunity to recover and prepare for what yet lies in wait."

\end{DndReadAloud}

The portal in this chamber appears only if the characters enter from area N8. The portal links to an identical one that appears simultaneously in area N1, permanently connecting these two areas. Looking through the portal is like looking through an open doorway into the other chamber. Any creature or object that passes through a portal emerges in the other location. Neither portal can be dispelled, and the portals last until Alyxian leaves the Netherdeep.

\section{Vents of Fury}
\DndQuote%
{}
{''I sacrificed everything for them! And now they tear at my flesh to slake their greed, with no thought for their savior. There is only one way to repay their ill gratitude.''}
{Alyxian the Apotheon}
The lowest chambers of the Netherdeep are called the Vents of Fury because of the steam and thermal fissures in the floor. These vents are manifestations of Alyxian's wrath, which also coalesces into monsters that inhabit this region. These creatures are marred with ugly ridges and streaks of ruidium that resemble festering scabs.

\subsection{Roleplaying the Apotheon}
In the Vents of Fury, the Apotheon's presence is a cruel and malicious hindrance. He doesn't seek understanding; he wants the characters to suffer as he has. Bitterness flavors every word he speaks.

The first time a character speaks out against the Apotheon while in the Vents of Fury, the Apotheon's voice audibly booms out, "You will not speak of me this way! Not after all we have shared!" This voice can be heard out to a range of 100 feet underwater.

The second time a character speaks out against the Apotheon in the Vents of Fury, the water around the character turns boiling hot, dealing 21 (6d6) fire damage to the character, and the Apotheon's voice roars, "Fool! You reveal the truth of your greedy, mortal heart."

If a character speaks out against the Apotheon a third time in the Vents of Fury, the Apotheon bellows, "So be it. You too shall know the suffering of betrayal." Sludgy ichor oozes out from the character's skin and coalesces, taking the form of a \textbf{cloaker} whose face resembles that of someone the character loves or thinks of fondly. This cloaker has the following changes to its stat block:


\begin{itemize}
\item The cloaker can breathe air and water.
\item It has a swimming speed equal to its flying speed.
\end{itemize}

The cloaker attacks the character, ignoring all other creatures. When either the cloaker or its target drops to 0 hit points, the cloaker disappears in a vermilion haze that fills its space and disperses after 1 minute. The area within the haze is Vision and Light.

\subsection{Vents of Fury Locations}
The Vents of Fury encompass areas N10 through N18 on the Netherdeep map.

\subsubsection{N10: Belching Gallows}
\begin{DndReadAloud}{}
This cavernous space slightly expands and contracts like the lung of a living, breathing creature. Natural fissures in the floor belch out columns of steam and bubbles, and the water here is warm. Red light also emanates from these vents, illuminating the chamber and the mouths of the tunnels that lead away from it.

\end{DndReadAloud}

Characters who have a passive Wisdom (Perception) score of 17 or higher notice something more:

\begin{DndReadAloud}{}
At the far end of the thirty-foot-high cavern, a monstrous fish with a red-streaked body chases a school of much smaller fish.

\end{DndReadAloud}

The expansion and contraction of the cavern is a natural phenomenon and can't be dispelled.

In the distance, a \textbf{slithering bloodfin} chases a \textbf{swarm of sorrowfish} (see appendix A for these creatures' stat blocks). A character who isn't carrying a light source can move around in the flooded cavern without attracting the attention of these creatures by succeeding on a DC 14 Dexterity (Stealth) check. If the check fails or the character makes no attempt to hide, the bloodfin notices and attacks the intruder. The sorrowfish are attracted to blood and attack any creature in the cavern that is wounded, including the bloodfin.

\subparagraph{Bubbling Water}
The bubbling columns of water rising from the vents are hot, but not hot enough to damage creatures that enter them. If the bloodfin is killed, the vents in the floor flare with bright red light and spew forth a flood of boiling-hot water. Any creature that starts its turn in the boiling-hot water must make a DC 15 Constitution saving throw, taking 21 (6d6) fire damage on a failed saving throw, or half as much damage on a successful one. The scalding heat lasts for 1 hour, after which the vents return to their normal illumination and the water temperature in the cavern again becomes bearably warm.

\subparagraph{Secret Door}
Hidden in the north wall is a secret door (see "Secret Doors" earlier in the chapter). On the other side of the secret door is a narrow passage that climbs steeply and ends before another secret door that opens into area N19. Either secret door is easily found from inside the connecting passage.

\subsubsection{N11: Frothing Hollow}
\begin{DndReadAloud}{}
A current threatens to pull you into this chamber, the walls of which are lined with ruidium spines. The Apotheon's ragged breathing is audible in every heave of water. Thermal vents in the floor give off swirling bubbles, dim light, and warmth.

\end{DndReadAloud}

A \textbf{light devourer} (see appendix A) swims in the dark water 50 feet above the cavern floor. It attacks any character who brings a light source into the chamber.

\subparagraph{Current}
If the characters kill the light devourer, the current intensifies for 1 hour. Until the current returns to normal, any creature that starts its turn in this chamber must succeed on a DC 14 Strength saving throw or be thrust against a wall by the current and skewered by ruidium spines. The creature takes 5 (1d10) piercing damage from the spines and must succeed on a DC 20 Charisma saving throw or gain 1 level of exhaustion. If the creature is not already suffering from ruidium corruption, it becomes corrupted when it fails the saving throw.

\subparagraph{Secret Doors}
Hidden in the south wall are two secret doors (see "Secret Doors" earlier in the chapter). One opens to reveal a narrow, spiraling tunnel that climbs gradually to area N5a. The other leads to a small side chamber (area N11a).

\subsubsection{N11a: Hidden Treasure}
\begin{DndReadAloud}{}
This small cave has a fifteen-foot-high ceiling and a sandy floor covered with tarnished gold coins. Lying among the coins are two curious items: a ceremonial dagger made of jade and a stoppered vial containing a bubbling purple liquid.

\end{DndReadAloud}

The characters can rest here without being disturbed, assuming they don't mind being underwater.

\subparagraph{Treasure}
Characters can recover the following treasure from the sandy floor:


\begin{itemize}
\item 750 gp (minted before the Calamity, these coins are worth double their monetary value to scholars or historians)
\item A jade dagger, its hilt and crossguard carved to resemble two green snakes coiled around one another (250 gp)
\item A potion that combines the beneficial effects of a Potion of Superior Healing and a \textit{potion of heroism}
\end{itemize}

\subsubsection{N12: Scalding Pit}
\begin{DndReadAloud}{}
Natural vents spew red light and jets of bubbling water into this warm, oblong chamber. Two enormous jellyfish-like creatures, faintly illuminated by the light coming from the vents, drift in the water at opposite ends of the chamber, their tentacles swaying. Each is within easy reach of one of the tunnels leading to this chamber.

\end{DndReadAloud}

The creatures floating in the chamber are two death embraces (see appendix A). They are oblivious to what happens around them and attack only if damaged or otherwise disturbed (see "Apotheon's Fury" below). Moving around a death embrace without touching it doesn't disturb it.

\subparagraph{Treasure}
In the middle of the chamber, lodged in two vents 10 feet apart, are a \textit{ruidium greatsword} and a \textit{ruidium shield} (both described in appendix B). To extract either item from its location, a character must reach into the boiling water, taking 5 (1d10) fire damage. A \textit{telekinesis} spell or similar magic can be used to dislodge these items safely.

\subparagraph{Apotheon's Fury}
If either the \textit{ruidium greatsword} or the \textit{ruidium shield} is removed from its resting place, the vents in the floor flare with bright red light and spew forth a flood of boiling-hot water, drastically raising the temperature of the water in the chamber, but only to a height of 30 feet. (The walls climb to a height of 100 feet; see "Features of the Netherdeep" earlier in the chapter.)

Any creature that starts its turn in the boiling water must make a DC 15 Constitution saving throw, taking 21 (6d6) fire damage on a failed save, or half as much damage on a successful one. The scalding heat lasts for 1 hour, after which the vents return to their normal illumination and the water temperature changes from boiling hot to bearably warm.

The death embraces, if they are still present, also take damage from the boiling water at the start of their turns. Roused by this, they attack any other creatures in the grotto, using their tentacles to grapple as many foes as they can before they float upward, away from the boiling-hot water.

\subparagraph{Secret Door}
A secret door in the north wall has a narrow, winding tunnel behind it (see "Secret Doors" earlier in the chapter). The tunnel ascends steeply and ends before another secret door, beyond which is area N21. Either secret door is easily found from inside the connecting tunnel.

\subsubsection{N13: Warriors of Wrath}
\begin{DndReadAloud}{}
Ruidium covers almost every surface of this thirty-foot-high cavern, in some places forming crystalline spikes that range from one foot to eight feet long. Life-sized, ruidium-encrusted stone statues of warriors jut from the walls, floor, and ceiling at odd angles, like jagged teeth in a diseased maw.

\end{DndReadAloud}

The statues in the cavern depict warriors who served with Alyxian. All are carved to look like they've been gravely wounded, and a few are missing limbs or their heads.

Three scuttling serpentmaws (see appendix A) hide behind the statues and ruidium spikes in the middle of the area. A character who tries to cross the cavern must succeed on a DC 19 Wisdom (Perception) check to avoid being Surprise when the serpentmaws attack.

\subparagraph{Apotheon's Fury}
When all three serpentmaws are dead, dozens of tendrils of ruidium emerge from the statues in the cavern. The tendrils persist for 1 minute, after which they recede back into the statues.

As long as the tendrils are present, any creature that enters the cavern or starts its turn there must succeed on a DC 15 Dexterity saving throw or be grappled by a tendril (escape DC 15). While grappled in this way, a creature is also restrained. No more than one tendril can grapple a single creature. Each tendril is an Ooze that has AC 15; 10 hit points; immunity to all damage except slashing; Strength, Dexterity, and Constitution scores of 20; and Intelligence, Wisdom, and Charisma scores of 1. A tendril reduced to 0 hit points disintegrates in a crimson haze, like blood diluted in water.

On initiative count 0, each creature grappled by a tendril takes 3 (1d6) piercing damage and 3 (1d6) psychic damage as it is flung against one of the cavern's ruidium spikes. The creature remains grappled after taking this damage; in addition, it must succeed on a DC 20 Charisma saving throw or gain 1 level of exhaustion. If the creature is not already suffering from ruidium corruption, it becomes corrupted when it fails the saving throw.

\subsubsection{N14: Carmine Sepulcher}
\begin{DndReadAloud}{}
Thermal vents in the floor bring light and heat into this eerie sepulcher, which has a ceiling thirty feet high. Strewn around the vents are the bones of large fish, though it's not clear what killed them.

The ruidium growing on the walls seems to have engulfed dozens of people, trapping them in red crystal sarcophagi as if they were bugs in amber. Every figure bears an expression of anguish.

\end{DndReadAloud}

The figures attached to the walls and entombed in ruidium sarcophagi are illusions. Most of them are images of unfamiliar people from Alyxian's past, but characters who examine the figures closely can see individuals they met while exploring the Grottoes of Regret, such as Alyxian's parents (area N2), Saqiri the drow hunter (area N3), the soldiers Kalagothe (area N4) and Talmyth (area N9), and young Meri (area N5). Any character who tries to chip away at the ruidium sarcophagus to free an illusory figure must succeed on a DC 15 Charisma saving throw or take 11 (2d10) psychic damage. As the attempt is made, the illusory figure trapped inside pleads for release, although its voice is indistinct through its coating of ruidium.

Each ruidium sarcophagus is a Large object that has AC 13, 27 hit points, and immunity to poison and psychic damage. Reducing a sarcophagus to 0 hit points destroys it and causes the illusory figure inside the sarcophagus to vanish.

One of the sarcophagi is empty—a hint that it is actually a secret door (see "Secret Door" below).

\subparagraph{Bones}
The floor contains the bones of three slithering bloodfins. If the secret door to area N14a is opened, the bones quietly knit themselves back together and attack any other creatures in the cavern. These skeletal creatures have glowing red eyes and use the \textbf{slithering bloodfin} stat block (see appendix A), with these changes:


\begin{itemize}
\item The skeletal bloodfins are Undead.
\item They have immunity to poison damage; resistance to piercing and slashing damage; and immunity to the charmed, deafened, frightened, paralyzed, and poisoned conditions.
\item They require no food, drink, air, or sleep.
\end{itemize}

\subparagraph{Secret Door}
In the southwest wall is a secret door made to look like one of the chamber's ruidium sarcophagi (see "Secret Doors" earlier in the chapter). Beyond the secret door is area N14a.

\subsubsection{N14a: Fragment of Rancor}
\begin{DndReadAloud}{}
A pale glow suffuses this small, dead-end chamber. The ceiling is fifteen feet high, and the floor is covered with fine, white sand. Tucked in an alcove is an altar of glittering crystal carved with small feet. A shimmering mote of light floats above it.

\end{DndReadAloud}

The crystal altar bears inscriptions in Celestial that say this shrine is dedicated to Avandra the Change Bringer. A character who doesn't speak Celestial can discern that the altar is tied to Avandra by succeeding on a DC 15 Intelligence (Religion) check. A follower of Avandra has advantage on this check.

\subparagraph{Fragment of Rancor}
The mote of light above the pedestal is the Fragment of Rancor (see "Fragments of Suffering" earlier in the chapter). When a character touches the mote, it disappears and becomes part of the character's soul. When this happens, all the characters receive a fleeting vision (see the Apotheon Lore table earlier in the chapter).

The character who is carrying the Fragment of Rancor gains the following benefit and drawback, and is aware of both:



\begin{description}
\item[Benefit.]
When you hit a creature with an attack, you can deal either an extra 2d6 psychic damage to the creature or 4d6 psychic damage to each creature within 5 feet of it. After you use this benefit, you must finish a short or long rest before you can use it again.
\item[Drawback.]
Whenever you are not unconscious and fail a saving throw, you take 2d6 psychic damage.
\end{description}





Give the Fragment of Rancor card (see appendix D) to the player whose character is carrying the fragment.

\subsubsection{N15: Spear of the Apotheon}
\begin{DndReadAloud}{}
This pit is devoid of thermal vents. The floor is strewn with humanoid skeletons clad in rusty armor as well as the bones of other creatures. Tendrils of ruidium weave along and through the battlefield like ivy. A shining spear stands upright in the center of the room, its shaft emerging from a pile of bones and rusting weapons.

\end{DndReadAloud}

The skeletons of warriors who died during the Calamity are mixed with the bones of giants, mastodons, and moorbounders. These bones are magical fabrications that remind the Apotheon of the carnage he witnessed as a warrior on the battlefield.

\subparagraph{Treasure}
Characters who search the area notice a dark recess in the northeast wall, at the back of which is a human skeleton in rusted chain mail. Tendrils of ruidium coil around the remains. The skeleton has a \textit{ruidium shortsword} (see appendix B) clenched in its bony hand.

\subparagraph{Taking the Spear}
The Spear of Warning sticking out of the pile of bones is a replica of Alyxian's spear. It shows no signs of rust, decay, or ruidium corruption, unlike the other weapons found throughout the chamber. It is imbued with the properties of a \textit{weapon of warning}. If the characters have already confronted Alyxian in the Heart of Despair (see chapter 7), the spear is just another piece of treasure waiting to be claimed, and there are no repercussions for someone who takes it.

If a character pulls the spear from the pile of bones before the characters confront Alyxian in the Heart of Despair, all the characters receive a fleeting vision (see the Apotheon Lore table earlier in the chapter). After the vision concludes, two beings suddenly appear in the chamber:

\begin{DndReadAloud}{}
A loud growl reverberates throughout the cavern. As the sound fades, a soft, melodic voice says, "You shouldn't have done that." The voice comes from a translucent, pale blue figure of a child who appears to have materialized out of nowhere. Floating next to him is a slender sculpture of a man carved out of ruidium and bearing a solemn expression.

"I am Theo Nathope," says the ghostly boy. He points to the ruidium golem. "This is Alyxian the Hunter. It has come to destroy you, but it won't do anything as long as you're with me. We must go now. Do you want to save the Apotheon and end this nightmare?"

\end{DndReadAloud}

\subparagraph{Theo}
Theo Nathope bears a striking resemblance to Alyxian as a young boy. (One or more players might realize that "Theo Nathope" is an anagram of "The Apotheon.") This apparition can't be harmed, banished, or dispelled. If the characters accept his offer to guide them, Theo leads them through area N16, shows them a secret door to area N17, and follows them into that location. There he hopes the characters will learn secrets that will help them redeem Alyxian. See the "Roleplaying Theo" sidebar for more information on how to portray the spectral child.

Theo glides through the water as an apparition should, moving at a speed of 30 feet. He can pass through solid objects and end his turn inside one, and he has truesight out to a range of 60 feet.

\begin{DndSidebar}[float=!b]{Roleplaying Theo}
Everything about Theo hints that there is more to him than what appears. His words are carefully chosen and intelligent. He often pauses before he speaks, and he tends to give vague or cryptic responses to questions. Although he might seem unhelpful, his coyness serves a purpose: he wants the characters to \textit{want} to save Alyxian, not simply to follow instructions.



Theo's intent is to assess the moral fiber of the characters. If he thinks they can save Alyxian by healing his centuries-old trauma, rather than letting him go free without a second thought, Theo tries to lead them to the tools that will help them do so. He tries to bond with characters who are empathetic and gentle, but he also takes a keen interest in discerning and clever characters, and he will ask them questions both personal and philosophical. Only after he bonds with a character does he show his more childlike side, sometimes revealing an innocent sense of wonder, a desire to play games and laugh, or a need for fondness and affection.



\end{DndSidebar}

\subparagraph{Alyxian the Hunter}
This ruidium golem resembles Alyxian the Apotheon. It remains as rigid as a statue until the characters attack it or until Theo vanishes from the Netherdeep (which happens in area N18). Either event causes the golem to animate and attack the characters.

Alyxian the Hunter can track the characters unerringly and knows the locations of all the secret doors in the Netherdeep. The golem can't leave the Netherdeep, however, and it is instantly destroyed if the spear from this area is taken into the Heart of Despair.

\textbf{Alyxian the Hunter} uses the \textbf{stone golem} stat block, with these changes:


\begin{itemize}
\item The golem is Medium. (Its Hit Dice and hit points remain the same.)
\item The golem has a swimming speed of 30 feet.
\item As an action, the golem can conjure a magic dart made of ruidium and hurl it at another creature it can see within 120 feet of itself. The dart strikes the target unerringly and vanishes on contact. The target must make a DC 17 Charisma saving throw. On a failed saving throw, the target takes 24 (7d6) psychic damage and gains 1 level of exhaustion. In addition, if the target is not already suffering from ruidium corruption, it becomes corrupted when it fails the saving throw. On a successful saving throw, the target takes half as much damage and suffers no other effects.
\end{itemize}

When Alyxian the Hunter drops to 0 hit points, the golem turns into an immobile vermilion haze that can't be harmed or dispelled. If the spear from this area has been taken into the Heart of Despair, the haze disperses as the golem is destroyed. Otherwise, the haze coalesces and solidifies into a new ruidium body 1 minute later, whereupon Alyxian the Hunter regains all its hit points and resumes the hunt.

\subsubsection{N16: Curse of Ruidus}
\begin{DndReadAloud}{}
Floating twenty feet above the floor in the middle of this chamber, rotating slowly on a slightly tilted vertical axis, is a ruidium sphere pockmarked with crater-like indentations. The moonlike sphere is not quite fifty feet in diameter. Two giant sharks, their mouths and skin encrusted with ruidium, swim around the moon's equator in opposite directions. Thermal vents in the floor heat the grotto and bathe the lower hemisphere of the moon in red light.

\end{DndReadAloud}

Two corrupted giant sharks (see appendix A) orbit the red moon. The sharks don't attack until certain conditions are met (described below) or they are attacked.

If Theo is with the characters, he crinkles his nose and gives the moon a disapproving look, then says, "Ruidus was the one enemy Alyxian could never face, yet the one he hated the most." Theo warns the characters against attacking the sharks and tells them not to "wake the red moon" by touching it, moving too quickly through the grotto, or casting a spell.

\subparagraph{Secret Doors}
Hidden in the southwest wall are two secret doors (see "Secret Doors" earlier in the chapter). Opening either door doesn't awaken the red moon.

\subparagraph{Sinister Red Moon}
The moon is immune to all damage and can't be affected by magic, nor can it be moved. If a character in the chamber makes an attack, casts a spell, touches the moon, or moves 15 feet or more during a single turn, dozens of large bloodshot eyes open on the face of the moon and gaze in all directions. Three things happen at once when this occurs:


\begin{itemize}
\item The secret door to area N16a opens.
\item If the sharks in this area are still alive, they attack the character who caused the moon's eyes to open.
\item Each character inside the chamber when the eyes open must succeed on a DC 15 Wisdom saving throw or be cursed by the moon's corrupted gaze. Determine the effect of each character's curse by rolling on the Curses of Ruidus table. The effect on a character can be ended by any magic that removes a curse.
\end{itemize}

Once the moon's eyes come open, they remain open, watching creatures that enter and leave the area without harming them.

If the moon disappears because of events that transpire in area N25 and the sharks are still alive, they attack any creatures that enter this chamber.

\begin{DndTable}[header=Curses of Ruidus]{cX}

d6 & Curse\\
1–2 & \textbf{Agent of Doom}. In your mind, a voice that sounds like Alyxian's says, "Those you love must die for your greatness, and believe me, I feel your pain." While within 30 feet of you, your allies have disadvantage on their saving throws until this curse on you ends.\\
3–4 & \textbf{Harbinger of Misadventure}. In your mind, a voice that sounds like Alyxian's says, "Your victories will bring ruin to those you love, as mine have done." While within 30 feet of you, your allies have disadvantage on their attack rolls until this curse on you ends.\\
5–6 & \textbf{Vortex of Malice}. In your mind, a voice that sounds like Alyxian's says, "Embrace the malice that swirls deep within your heart, as it does within mine." You can't take the Help action, nor can you use a spell, magic item, or class feature to restore hit points or grant temporary hit points to another creature, until this curse on you ends.\\
\end{DndTable}

\subsubsection{N16a: Fragment of Abhorrence}
This small chamber is hidden behind a secret door. If the characters explore this area while Theo is with them, he waits for them outside.

\begin{DndReadAloud}{}
In the middle of this small chamber is a half-buried stone statue of Alyxian tilted back at a slight angle. Bits of the statue have broken off, and tendrils of ruidium rise from the damaged areas. A dull red key with a crescent moon-shaped handle is sticking out from a hole in the statue's chest, and a bright mote of light is cupped in the statue's outstretched hands.

\end{DndReadAloud}

The key is made of rusty iron and is safe to touch. It opens the door in area N23. The characters can remove the key by carefully jiggling it free. In area N23, the key glows brightly enough to shed dim light in a 5-foot radius.

\subparagraph{Fragment of Abhorrence}
Cupped in the statue's hands is the Fragment of Abhorrence (see "Fragments of Suffering" earlier in the chapter). When a character touches the mote of light, it disappears and becomes part of the character's soul. When this happens, all the characters receive a fleeting vision (see the Apotheon Lore table earlier in the chapter).

The character who is carrying the Fragment of Abhorrence gains the following benefit and drawback, and is aware of both:



\begin{description}
\item[Benefit.]
Once on each of your turns when you hit a creature with a weapon attack roll, you can force the target to move up to 10 feet away from you in a direction of your choice. A creature that can't be frightened is immune to this effect.
\item[Drawback.]
If you start your turn frightened, you take 2d6 psychic damage.
\end{description}





Give the Fragment of Abhorrence card (see appendix D) to the player whose character is carrying the fragment.

\subsubsection{N17: Heart of Truth}
This cavern is hidden behind two secret doors (see "Secret Doors" earlier in the chapter). If Theo has guided the characters to this location, he draws their attention to the northernmost secret door, then passes through it and waits for them on the other side.

When the characters enter this area, read the following boxed text:

\begin{DndReadAloud}{}
Thermal vents in the floor heat and illuminate this thirty-foot-high, oblong cavern. The ruidium-veined walls expand and contract in regular pulses like the beating of a heart. A ruidium archway in the far wall leads to a darker chamber.

\end{DndReadAloud}

The expansion and contraction of the cavern is a natural phenomenon and can't be dispelled.

\subparagraph{Ruidium Archway}
The ruidium archway that stands between this chamber and the adjoining grotto (area N18) is 10 feet tall and 8 feet wide. The first time a character passes through the archway, all the characters receive a fleeting vision (see the Apotheon Lore table earlier in the chapter).

\subparagraph{Theo}
If Theo is with the characters, he shares some thoughts on how they can help Alyxian:

\begin{DndReadAloud}{}
"Alyxian's suffering has made him callous and selfish," Theo says. "I know he can be reminded of his goodness."

A human-sounding howl echoes throughout the grotto. Darkness engulfs the cavern, swallowing you in an inky embrace, flooding your mind with visions. Of Ank'Harel. Of Wildemount. Of a godlike wrath that ravages the world until only corpses and smoking ruins remain. The visions abate, and you are back where you started. "That is what awaits you in the Heart of Despair," says the ghostly child. "Well, kind of."

\end{DndReadAloud}

If asked what he means by "kind of," Theo solemnly beckons for the characters to follow him as he passes through the ruidium archway into area N18.

\subsubsection{N18: Fragment of Melancholy}
\begin{DndReadAloud}{}
This pit is devoid of thermal vents, and the water here is cold and dark. A few tendrils of ruidium cling to the walls around the archway. Within the darkness, you see a small, pale white light hovering six feet above the ground in a crevice to the north.

\end{DndReadAloud}

When a character touches the mote of light (see "Fragment of Melancholy" below), three chuuls hidden under the sandy floor in the middle of the grotto rise and attack. Theo shrieks and cowers from the chuuls, which do not attack him. The chuuls pursue prey into area N17 but won't go any farther.

\subparagraph{Fragment of Melancholy}
The mote of light in the northern crevice is the Fragment of Melancholy (see "Fragments of Suffering" earlier in the chapter). When a character touches the mote, it disappears and becomes part of the character's soul. When this happens, all the characters receive a fleeting vision (see the Apotheon Lore table earlier in the chapter).

The character who is carrying the Fragment of Melancholy gains the following benefit and drawback, and is aware of both:



\begin{description}
\item[Benefit.]
When you fail a Wisdom or Charisma saving throw, you can reroll it with advantage, potentially turning a failure into a success. After you use this benefit, you must finish a long rest before you can use it again.
\item[Drawback.]
If you use this fragment's benefit and it doesn't turn a failed saving throw into a success, you are incapacitated until the end of your next turn, overcome with despair.
\end{description}





Give the Fragment of Melancholy card (see appendix D) to the player whose character is carrying the fragment.

\subparagraph{Theo}
If Theo is with the characters, he shares the following information with them before they leave the grotto:

\begin{DndReadAloud}{}
"You have seen and heard the creatures in here. Their rage is Alyxian's rage. Their sorrow is his own. Their yearning to escape is his yearning to escape. And when he is angry, Alyxian curses and scorns the gods. But he always returns to prayer, as if hoping for an answer to why the gods abandoned him.

"When Alyxian's rage is spent, temporarily, the Netherdeep grows dark and quiet. In those moments, I hear gentle laughter echoing through the grottoes and I feel a lightness... dare I say a budding joy. I think in such moments he finds peace.

"A piece of the Apotheon remains uncontaminated by regret, fury, yearning, and despair, submerged beneath all his anger. When you see him in his present form, appeal to this shard of his psyche. Please, you must not simply free him. Do whatever it takes to truly save him."

\end{DndReadAloud}

When the characters are prepared to move on, Theo adds:

\begin{DndReadAloud}{}
"Please don't judge him too harshly. There is goodness in him still, buried deep under so many centuries of anguish."

\end{DndReadAloud}

The first chance he gets after sharing this information, Theo disappears—at a moment when no one is looking at him—and never returns.

\subparagraph{Alyxian the Hunter}
If Alyxian the Hunter is stalking the characters, they find the golem waiting for them at the northern end of the tunnel outside area N17. See area N15 for more information about the golem.

\section{Chasm of Yearning}
\DndQuote%
{}
{''Those I saved live on. Their joy, their struggles, their ambitions, filter into my home in the depths like sunlight through the waves. How I long to feel that again—the fluttering of another living soul.''}
{Alyxian the Apotheon}
This region of the Netherdeep reflects the furtive tunneling of an imprisoned mind yearning for freedom. Yet not all is as it seems. The Chasm of Yearning is a region filled with falsehoods and mirages of things too good to be true. If lit by magic, the walls of this region are a deep, soothing blue, except where interrupted by illusions or ruidium growths.

The Apotheon's longing for freedom and friendship is felt most strongly here, and the creatures that reside in this region of the Netherdeep appear withered, starved, and sometimes translucent.

\subsection{Roleplaying the Apotheon}
Within the Chasm of Yearning, the Apotheon's presence manifests as raw desire: the desire to be seen and understood by others, the desire for rest, and the desire to be among heroes who know the meaning of sacrifice. In this region, Alyxian can manifest forms that represent his loneliness and yearning, appearing as a human made of sand, a will-o'-wisp, or an air bubble (see "Weird Events" below).

\subsection{Weird Events}
Whenever the characters leave one location and move to another one in the Chasm of Yearning, roll on the Weird Events table to determine what occurs. An event can occur more than once.

\begin{DndTable}[header=Weird Events]{cX}

d10 & Event\\
1 & The characters suddenly find themselves in another area in the Chasm of Yearning, along with all their equipment, but with no memory of how they got there. You choose the new location, excluding areas N24 and N25.\\
2 & A 1-foot-diameter air bubble drifts toward the party, moving gently while avoiding creatures in its path. If a creature touches it, the bubble pops and releases a crash of thunder that can be heard out to a range of 60 feet. Any creature within 10 feet of the popping bubble takes 5 (1d10) thunder damage. The bubble follows one randomly determined character and disappears when that character leaves the Chasm of Yearning or when you roll on this table again.\\
3 & A \textbf{will-o'-wisp} becomes visible and tries to lure the characters to area N25, then turns invisible once it arrives there. If attacked, the will-o'-wisp retaliates. Otherwise, it is indifferent toward the characters and harmless.\\
4–5 & Sand on the floor assumes a vaguely humanlike form in an unoccupied space as close to the characters as possible. This sandy shape follows one randomly determined character. It is harmless and can't be damaged or dispelled, but it loses its form and reverts to ordinary sand when the character it is following leaves the Chasm of Yearning or when you roll on this table again.\\
6–10 & The characters experience a vision from Alyxian's past (see the Apotheon Lore table earlier in the chapter).\\
\end{DndTable}

\subsection{Chasm of Yearning Locations}
The Chasm of Yearning encompasses areas N19 through N25 on the Netherdeep map.

\subsubsection{N19: Cavern of Many Hands}
\begin{DndReadAloud}{}
This narrow cavern is fifty feet long, with a ten-foot-high ceiling and an exit at the far end. The floor and walls are covered in a five-foot-thick, rippling carpet of thin, pale arms ending in hands that have long, delicate fingers. At first glance, these limbs might be taken for seaweed. Whispering voices begin to echo in your mind, saying things like "Join us," "Be with us," and "Never be alone again."

\end{DndReadAloud}

The arms and hands that grow out of the floor and walls reach for passersby, trying to touch them. This effect is a manifestation of the Apotheon's longing to feel the touch of another living creature. These limbs are not aggressive—in fact, they are chillingly gentle as they try to hold the characters' hands, caress their arms, and cradle their heads.

Any character who crosses the cavern notices, with a successful DC 13 Wisdom (Perception) check, that several of the hands are holding items: a box, a key, a flask of red liquid, and a rolled-up scroll. These items are as follows:


\begin{itemize}
\item The locked box is decorated with wedding bells and inscribed with the words "For those whose love never wavered." The box contains a black pearl (500 gp), a violet garnet (100 gp), and a pair of delicate platinum rings (250 gp each).
\item The key is made of tarnished silver and unlocks the box. (Without this key, a character can use \textit{thieves' tools} to try to pick the box's lock, doing so with a successful DC 14 Dexterity check.)
\item The flask contains a Potion of Superior Healing.
\item The rolled-up scroll is a \textit{spell scroll} of \textit{Otiluke's resilient sphere} tied with a green ribbon.
\end{itemize}

\subparagraph{Secret Door}
Hidden in the south wall is a secret door that requires a successful DC 25 Wisdom (Perception) check to locate (see "Secret Doors" earlier in the chapter). This door is hard to spot because of the spindly arms and grasping hands that grow on and around it. Behind the secret door is a narrow, winding tunnel that steeply ascends to area N10.

\subsubsection{N20: Cavern of Many Eyes}
\begin{DndReadAloud}{}
Weblike strands of bright ruidium cling to the walls of this grotto, spreading out from dozens of large, bloodshot eyes that follow your every move. These eyes are embedded in the walls between ten and twenty feet off the floor.

Five patches of seaweed sprout from the sandy floor. Inside each patch is an oval mirror four feet tall and two feet wide. The seaweed not only holds the mirrors in place but also frames them. You see displayed in each mirror what appears to be an important moment in the life of the Apotheon.

\end{DndReadAloud}

\subparagraph{Eyes}
The bloodshot eyes embedded in the walls track creatures as they move through the room. Their stony eyelids snap shut if necessary to protect the eyes against incoming attacks, blinding light, or harmful magic. Each eye is 1 foot in diameter and resembles a bloodshot human eyeball in appearance and texture. The eyes do nothing but watch and stare.

\subparagraph{Mirrors}
A \textit{detect magic} spell reveals an aura of divination magic around each mirror. The five mirrors show particular moments in Alyxian's past, each one playing a short, silent scene that repeats every 30 seconds. If the characters haven't yet seen all the visions in the Apotheon Lore table presented earlier in the chapter, you can replace any or all of the scenes in the mirrors with entries from that table. Otherwise, the scenes are as follows:


\begin{description}
\item[Mirror 1:]
Born in the Light of Ruidus. A baby boy lifts his arms into the air as he cries. A red light washes over him. Instead of toys lying around him, a spear and a dagger rest near where he lays his head.
\item[Mirror 2:]
A Selfless Soul. A young Alyxian, ten years old at most, holds out a piece of bread to a sick, elderly figure, who reaches for it with a warm smile.
\item[Mirror 3:]
Beginnings of a Hero. This solemn scene depicts Alyxian as a young man, holding a sword and looking at it with a pained expression. The ground around his feet is strewn with rubble and bodies.
\item[Mirror 4:]
Defending the Weak. Alyxian, as an adult, faces a gloomstalker that bears down on him while two children huddle behind him.
\item[Mirror 5:]
Beacon of Hope. Alyxian, his demeanor exuding determination, stands atop a raised mound, thrusting his spear in the air as he delivers a rallying cry. The soldiers below respond by raising their weapons in salute.
\end{description}

A character who gazes into a mirror while standing directly in front of it sees themself instead of Alyxian, and that mirror's image now depicts an important or defining moment in the character's past. Ask the player to describe the scene that their character sees reflected in the mirror's surface. If the player does a particularly good job describing a scene from their character's past, award \textit{inspiration} to the character; no character can receive inspiration in this way more than once. If the player doesn't come up with a scene before deciding that the character looks away, the character takes 22 (4d10) psychic damage and is incapacitated for 1 minute.

Each mirror has AC 13, 8 hit points, and immunity to poison and psychic damage. When a mirror is destroyed or cut loose from the seaweed that holds it in place, it ceases to display images and becomes nonmagical.

\subparagraph{Leaving the Grotto}
Any character who tries to leave the grotto without first seeing themself in at least one mirror hears a garbled voice in their mind say, "Don't leave yet! Let us see! Show us a moment from your past—one that shaped your heroic journey." If met with belligerence or silence, the voice replies, "Oh, have you no worth, then? You, too, shall be forgotten." If met with confusion or gentleness, the voice replies, "Please, friend, let us adore you. Let us catch a glimpse of the path you had to walk to become the hero we see before us." A character can freely ignore the voice, and once the character leaves this area, the voice no longer speaks to them (even if the character returns to this area later).

\subsubsection{N21: Cavern of Many Peaceful Ends}
When the characters first peer into this area, but before they enter it, describe it as follows:

\begin{DndReadAloud}{}
Weblike strands of bright ruidium cling to the walls of this otherwise empty underwater grotto.

\end{DndReadAloud}

When one or more characters enter this chamber for the first time, add:

\begin{DndReadAloud}{}
A dim red light appears above your head and takes the form of a crown with jagged spikes.

\end{DndReadAloud}

A crown that sheds dim light in a 5-foot radius appears above the head of each character who has just entered the chamber for the first time. When a crown appears, the character sporting it must make a DC 19 Wisdom saving throw. On a failed save, the character is incapacitated. While incapacitated in this way, the character is also charmed and does nothing except make their way to area N22 by the most direct route—most likely the north tunnel (unless that path is blocked). A character who can't be charmed can still be incapacitated by the crown, but isn't compelled to go to area N22.

While charmed by the crown, a character experiences a vision of a wonderful future in which the character believes they have achieved their goals, won happiness, and found paradise. In effect, the crown shows the character the sort of happy ending that Alyxian wishes he had experienced. Ask the player to describe their character's idea of perfect happiness. Perhaps the character falls in love with Ank'Harel and becomes a person of influence in the city, or perhaps the character imagines returning to Jigow and raising a family. If a player is having trouble with this request, use one or more of the following questions as prompts:


\begin{itemize}
\item What would make your character happy?
\item Where does your character imagine spending the best years of their life?
\item Whom does your character imagine spending the rest of their life with?
\end{itemize}

When a character bearing a crown enters area N22 or is targeted by any magic that ends a curse, the crown disappears and its effect on that character ends.

\subparagraph{Secret Door}
Hidden in the west wall is a secret door (see "Secret Doors" earlier in the chapter), beyond which is a narrow, winding tunnel that steeply ascends to area N12. The secret door outside area N12 is found automatically from inside the tunnel.

\subsubsection{N22: Cavern of Many Graves}
Each character who peers into this area sees the same thing:

\begin{DndReadAloud}{}
Jutting out of the sandy floor of this grotto are a dozen seaweed-choked gravestones. Most of the stones are engraved with the names of people you know or knew, but one of them has your name on it. Below your name on the stone is a shallow indentation with an unusual yet familiar shape.

\end{DndReadAloud}

A \textit{detect magic} spell reveals an aura of illusion magic around each gravestone, and this magic is the reason why each character sees different inscriptions on the stones. The gravestone that bears a character's name is the same one for each character. It doesn't take an ability check to recognize that the shallow indentation beneath the name perfectly matches the size and shape of the Jewel of Three Prayers (Exalted) in its Exalted State. If the jewel is pressed into the indentation, the secret door to area N22a swings open.

\subparagraph{Gravedigging}
If the characters start digging up graves, they find a number of faceless corpses buried under the sand in shallow graves equal to the number of characters in the party. These corpses have nothing in their possession. Removing a corpse from its grave causes all the corpses in the area to animate and attack. Each corpse uses the \textbf{revenant} stat block and transforms into a 1-ounce pebble of ruidium when reduced to 0 hit points.

\subparagraph{Secret Door}
A secret door is hidden in the southwest wall (see "Secret Doors" earlier in the chapter). Unlike most of the secret doors in the Netherdeep, this one is locked. It can be opened only by pressing the Jewel of Three Prayers (Exalted) into the indentation (as described above) or by the use of a \textit{knock} spell or similar magic.

\subsubsection{N22a: Fragment of Intransigence}
\begin{DndReadAloud}{}
Floating about this grotto are deteriorated bedroom furnishings: a bed frame with no mattress, a chair, a table, a dresser missing its drawers, and a wooden chest, all within easy reach. The chest's lid is emblazoned with the symbol of Corellon: a four-pointed star flanked by a pair of crescents. A soft white glow spills out of the chest's many cracks.

\end{DndReadAloud}

The chamber's furnishings fall apart if handled roughly. The chest is unlocked and contains a tiny mote of light—the Fragment of Intransigence (see "Fragments of Suffering" earlier in the chapter).

\subparagraph{Fragment of Intransigence}
When a character touches the mote of light, it disappears and becomes part of the character's soul. When this happens, all the characters receive a fleeting vision (see the Apotheon Lore table earlier in the chapter).

The character who is carrying the Fragment of Intransigence gains the following benefit and drawback, and is aware of both:



\begin{description}
\item[Benefit.]
Another creature can't force you to move anywhere you don't want to go.
\item[Drawback.]
You can't take the Disengage or Dodge action.
\end{description}





Give the Fragment of Intransigence card (see appendix D) to the player whose character is carrying the fragment.

\subsubsection{N23: Grotto of Many Futures}
The following description assumes that the characters arrive from area N20 or N22:

\begin{DndReadAloud}{}
The walls of this grotto have been meticulously carved with images of shops and town houses. The features of the buildings change from time to time, as though an invisible force were swapping out one facade for another or replacing one architectural flourish with another. The one feature that seems never to change is a wooden door set into the far wall at floor level.

In the middle of the grotto floats a giant jellyfish-like creature with sixty-foot-long tentacles. Swarms of tiny, ugly fish flit between its translucent tentacles.

\end{DndReadAloud}

This fake city street is a manifestation of Alyxian's longing for a home he can love, a fact that a character can intuit with a successful DC 15 Wisdom (Insight) check.

In the center of the cavern is a \textbf{death embrace} that has a \textbf{swarm of sorrowfish} darting around its tentacles (see appendix A for these creatures' stat blocks). These creatures ignore the characters, attacking only in self-defense or when something becomes entangled in the death embrace's tentacles, such as what can happen when the east door is touched (see below).

\subparagraph{East Door}
The door in the east wall resembles a sturdy cellar door with a built-in lock. The first time the door is touched, it gives off a flash of crimson light. Any creature within 10 feet of the door when this happens must succeed on a DC 18 Strength saving throw or be pushed 20 feet away from the door—and into the death embrace's tentacles, if that creature is still alive and present.

The key to the door can be found in area N16a. A character can use an action to try to pick the lock with \textit{thieves' tools}, doing so with a successful DC 22 Dexterity check. A \textit{knock} spell or similar magic also opens the door.

After the door is opened, the characters can see that the water doesn't pass beyond its threshold. On the other side of the door is a flight of dry steps leading down to a red-lit cavern (area N24).

\subsubsection{N24: Cavern of Many Waterfalls}
The following boxed text assumes that the characters are coming from area N23:

\begin{DndReadAloud}{}
This thirty-foot-high cavern is not submerged. The air smells stale, coral-like growths made of ruidium festoon the walls, and the floor is covered with a soft layer of dry sand. Sparkling waterfalls trickle down the walls in a few places, but the water disappears in a glittery froth as it strikes the floor.

At the east end of the cavern, opposite you, is a floor-to-ceiling opening with a waterfall pouring down in front of it. Standing in front of this waterfall is a blue-skinned figure with glowing red eyes and veins of ruidium on her face and arms. Folded behind her back is a pair of feathery white wings.

\end{DndReadAloud}

The blue-skinned angel is a deva named \textbf{Perigee}. She is a servant of the Moon Weaver who fought with Alyxian in his final battle. Perigee was destroyed but later re-formed at the Moon Weaver's side. For centuries thereafter, she searched for Alyxian and eventually found her way to the Netherdeep, where she became mired in despair and corrupted by ruidium. In the Celestial tongue, she accuses the characters of disturbing the Apotheon's dreams and attacks them.

A character can use an action to try to convince Perigee that the party wants to save Alyxian. Allow the player to roleplay this interaction, then have the character make a DC 16 Charisma (Persuasion) check, with advantage if you think the player roleplayed the interaction well. On a failed check, Perigee isn't swayed. On a successful check, Perigee does nothing on her next turn as she considers the character's words (as long as she's not attacked or threatened in the meantime). At the end of that turn, she coolly remarks that anyone who can't defeat her has little hope of besting Alyxian in the Heart of Despair. On her next turn, she resumes her attack. Each character can attempt the check only once.

\subparagraph{Perigee's Defeat}
If Perigee is defeated, she disappears with a sigh. No longer bound to the Netherdeep, the deva returns to the Moon Weaver's side with all traces of ruidium gone. When Perigee disappears, a secret door in the north wall swings open (see "Secret Door" below).

\subparagraph{Placating the Apotheon}
If one or more characters tried to convince Perigee of the party's good intentions, their attempt to sway Perigee has a positive effect on the final battle against Alyxian (see "Emotional Healing" in chapter 7), regardless of the success of the attempt.

\subparagraph{Secret Door}
The secret door in the north wall (see "Secret Doors" earlier in the chapter) opens to reveal a short, dry tunnel leading to area N24a. The tunnel is dimly lit by silvery moonlight emanating from that chamber, and characters can hear the sound of the waterfalls coming from that direction.

\subparagraph{Waterfalls}
The waterfalls are illusory and can't be dispelled. Nothing that touches a waterfall or passes through it gets any wetter than it already is. The waterfall in the passageway to N25 doesn't block line of sight and is harmless; characters can move through it as easily as they would walk through a curtain.

\subsubsection{N24a: Fragment of Deception}
This area, like area N24, is not submerged.

\begin{DndReadAloud}{}
The air smells fresh here. Waterfalls tumble down the walls of this moonlit grotto, the water disappearing when it touches the soft soil that covers the floor. The area is filled with jungle foliage, and water droplets roll off the leaves onto fragrant tropical flowers.

A beam of silvery moonlight shines down on a curved alabaster bench facing the west wall. Carved into the west wall is a frieze depicting a young girl with wavy strands of light for limbs, accompanied by a pack of wolves. The girl's sculpted hands protrude from the wall and are held out in a cupped position.

\end{DndReadAloud}

If the characters examine the bench, add:

\begin{DndReadAloud}{}
An inscription carved into the bench reads, in Common, "Sit here and reflect on the Moon Weaver."

\end{DndReadAloud}

The waterfalls, flora, and moonlight are illusions that can't be dispelled. If a character sits on the alabaster bench and uses a mirror or other reflective surface to deflect the moonlight so that it shines on the frieze of the Moon Weaver, a mote of pale white moonlight appears in the sculpture's cupped hands. This mote of light is the Fragment of Deception (see "Fragments of Suffering" earlier in the chapter).

\subparagraph{Fragment of Deception}
When a character touches the mote of light in the Moon Weaver's cupped hands, it disappears and becomes part of the character's soul. When this happens, all the characters receive a fleeting vision (see the Apotheon Lore table earlier in the chapter).

The character who is carrying the Fragment of Deception gains the following benefit and drawback, and is aware of both:



\begin{description}
\item[Benefit.]
When you take damage, you can use your reaction to turn invisible and teleport up to 60 feet to an unoccupied space you can see. You remain invisible until the start of your next turn or until you make a damage roll or cast a spell.
\item[Drawback.]
You have disadvantage on Wisdom checks.
\end{description}





Give the Fragment of Deception card (see appendix D) to the player whose character is carrying the fragment.

\subsubsection{N25: Fragment of Loathing}
Characters who pass through the waterfall at the east end of area N24 find themselves submerged in water once again:

\begin{DndReadAloud}{}
This enormous, submerged grotto has tall clumps of white sea grass sprouting from the sandy floor. Looming above you is a red moon pockmarked by craters. You sense great fury emanating from it.

Spaced around the perimeter of the moon, resting among the sea grass, are three stone altars carved with the symbols of Avandra the Change Bringer, Corellon the Arch Heart, and Sehanine the Moon Weaver. Each altar has an inscription carved into its top and a ruidium spear leaning against it.

\end{DndReadAloud}

This area represents the culmination of Alyxian's dreams: a world in which Ruidus can be destroyed. The moon, which hovers 20 feet above the floor, is called the Effigy of Ruidus. It is an immobile, Gargantuan object that has AC 17, 200 hit points, and immunity to poison and psychic damage. Reducing the Effigy of Ruidus to 0 hit points causes it to explode into ruidium dust and causes the moon in area N16 to crumble to dust as well.

When a creature moves within 20 feet of the Effigy of Ruidus, or when the effigy takes damage, read:

\begin{DndReadAloud}{}
Colors begin to swirl around the moon's surface, forming three layers of shimmering energy: an inner layer of bright vermilion, a middle layer of rich maroon, and an outer layer of crackling purple.

\end{DndReadAloud}

When the shield appears around the Effigy of Ruidus, have all creatures in the area roll initiative. On initiative count 20 and again on initiative count 10, the effigy makes two Mind Warp attacks (see below), targeting creatures in the area at random.


\subparagraph{Mind Warp}
\textsl{Melee or Ranged Spell Attack:} +9 to hit, reach 5 ft. or range 120 ft., one creature. Hit: 14 (4d6) psychic damage.





\subparagraph{Shield Layers}
Each shield layer around the Effigy of Ruidus has its own property:


\begin{description}
\item[Purple Layer.]
This layer crackles with magical energy that deals 14 (4d6) lightning damage to any creature that starts its turn within 10 feet of the Effigy of Ruidus.
\item[Maroon Layer.]
This layer releases a pulse of necrotic energy every time the Effigy of Ruidus takes damage. The pulse deals 7 (2d6) necrotic damage to creatures in the chamber.
\item[Vermilion Layer.]
This layer increases the moon's AC from 17 to 22.
\end{description}

\subparagraph{Spears and Altars}
The ruidium spears that lean against the altars can be used as nonmagical weapons. Their primary purpose is to destroy the shield layers that protect the Effigy of Ruidus.

A character who wears gloves or gauntlets can handle any of the ruidium spears safely. Without such protection, a character gains 1 level of exhaustion at the start of each of their turns for each ruidium spear they are holding; in addition, a character not already suffering from ruidium corruption becomes corrupted when it gains 1 or more levels of exhaustion in this way.

To destroy a shield layer, a character must use an action to imagine the layer gone while grasping the spear, use the spear to damage the effigy, or hold the spear while casting a spell that deals damage to the effigy. The inscriptions on the altars are written in Celestial. The spear that leans against an altar is the one referred to in that altar's inscription:


\begin{description}
\item[Altar of Avandra (Northwest of the Effigy).]
The inscription on this altar reads, "Grasp this spear and banish the purple shield."
\item[Altar of Corellon (Southwest of the Effigy).]
The inscription on this altar reads, "Grasp this spear and use it to dispel the maroon shield."
\item[Altar of Sehanine (East of the Effigy).]
The inscription on this altar reads, "Grasp this spear and imagine the red moon without its vermilion shield."
\end{description}

\subparagraph{Fragment of Loathing}
When the effigy of Ruidus is destroyed, the Fragment of Loathing (see "Fragments of Suffering" earlier in the chapter), in the form of a tiny mote of white light, appears in its place and floats down until it's a few feet above the floor. When a character touches the mote of light, it disappears and becomes part of the character's soul. When this happens, all the characters receive a fleeting vision (see the Apotheon Lore table earlier in the chapter).

The character who is carrying the Fragment of Loathing gains the following benefit and drawback, and is aware of both:



\begin{description}
\item[Benefit.]
When a creature damages you with a weapon attack or a spell, you can focus your hatred on that creature. Until the end of your next turn, you have advantage on attack rolls you make against the creature. You can focus your hatred on only one creature at a time.
\item[Drawback.]
You have disadvantage on Charisma checks.
\end{description}





Give the Fragment of Loathing card (see appendix D) to the player whose character is carrying the fragment.

\subsubsection{N26: Sanctuary of Despair}
\begin{DndReadAloud}{}
This immense cavern contains a spherical black mass resembling a giant organ or cyst that nearly fills the space. Measuring over one hundred feet in diameter, it has veins of ruidium crisscrossing its surface, which is covered with spiky protuberances that retract and elongate as the sphere pulsates.

\end{DndReadAloud}

This cavern is an appropriate location for the characters to encounter their rivals (see "Rivals in the Netherdeep" earlier in the chapter). There are three ways this scene can play out:


\begin{description}
\item[Friendly Rivals.]
If the rivals have become the characters' friends or allies, they agree to help the characters in the final battle against Alyxian, assuming they've all acquired enough Fragments of Suffering to enter the Heart of Despair together. The characters might have to choose which of the rivals to take into this final confrontation and which ones to leave behind, based on who has a Fragment of Suffering and who doesn't. Friendly rivals are willing to transfer one or more of their Fragments of Suffering to characters who don't have one.
\item[Indifferent Rivals.]
The rivals and the characters are both trying to reach the heart first. The two groups have deep-seated animosity and bonds of love or camaraderie in equal measure. The rivals might be willing to transfer Fragments of Suffering they've acquired to characters whom they view as better equipped to deal with Alyxian. Dermot and Maggie are the most helpful, while Irvan and Galsariad are the least helpful. Ultimately, Ayo makes the decisions.
\item[Hostile Rivals.]
Ruidium corruption has made the rivals callous and selfish. They still need one or more Fragments of Suffering to enter the heart as a group, and they demand that the characters surrender any fragments they've obtained. If the characters refuse, the rivals attack them. If one or more characters enter the Heart of Despair during this confrontation, as many rivals pursue them as are able to do so.
\end{description}

\subparagraph{Entering the Heart}
The surface of the Heart of Despair feels like mud to any creature carrying one or more of Fragments of Suffering. Such a creature can push its way through the outer surface and enter the heart's air-filled interior. To other creatures, the Heart of Despair feels like a crusty and impenetrable solid object.

If a character brings the spear from area N15 into the Heart of Despair, the golem-like creature called Alyxian the Hunter is destroyed instantly and permanently. If the characters have avoided Alyxian the Hunter so far, have it catch up to the party in this area, so that one or more characters can witness its destruction when the spear is brought into the Heart of Despair.

When one or more characters enter the Heart of Despair, read:

\begin{DndReadAloud}{}
The Fragments of Suffering shield you from harm as you are drawn into the fleshy, undulating Heart of Despair. The fragments then burn within you for a moment before they are torn from your soul.

\end{DndReadAloud}

The Fragments of Suffering disappear from the characters as they enter the Heart of Despair, which they can't leave until they confront the Apotheon (as described in the next chapter).

\chapter*{The Heart of Despair}
\DndDropCapLine{T}{he Heart of Despair is a prison of the} Apotheon's own making. Nothing traps him in the Netherdeep except his belief that he doesn't deserve to leave. He has called the characters to his location so that they can deem him worthy of returning to the world. By the time they're ready to enter the Heart of Despair, the characters should have learned enough about the Apotheon to pass judgment on his future. Their options are as follows:



\begin{itemize}
\item Allow the Apotheon to go free without first healing his emotional wounds (and in doing so, cause his corruption to wreak havoc throughout Exandria)
\item Deem the Apotheon beyond redemption and destroy him
\item Heal the Apotheon's emotional wounds and bring him peace before granting him his freedom
\end{itemize}

\section{Running This Chapter}
Read this chapter thoroughly before running it, paying particular attention to the stat blocks for Alyxian's forms and the "Lair Actions" section, which describes actions the Apotheon can take in any of his forms.

\section{Within the Heart}
Characters who enter the Heart of Despair find themselves standing in the middle of the area, just south of a statue of Alyxian (area H1). Describe the environment as follows:

\begin{DndReadAloud}{}
You pass through clammy darkness and tumble onto the glassy surface of a lake, which feels as solid under your feet as any floor. Rising out of the water are three rocky isles, upon which rest crumbling structures that look like ancient prayer sites.

Waterfalls tumble from the clear sky, giving off mist that blankets the area in shimmering silver. Directly in front of you, a life-sized statue of Alyxian armed with a spear and a shield stands on a white stone pedestal that rises inches above the water's surface. The statue's expression is serene, much like the water.

\end{DndReadAloud}

The interior of the Heart of Despair is a cylindrical space 200 feet high and over 200 feet in diameter. The water here is 20 feet deep.

Every visitor to the Heart and each of Alyxian's forms gains the benefit of a \textit{water walk} spell upon arrival; this effect lasts for 1 hour. A creature can end the effect on itself at any time (no action required), enabling it to sink below the water's surface. Once the effect ends on a creature, it can't be reactivated.

\subsection{Dealing with the Rivals}
Rivals who follow the characters into the Heart of Despair behave in one of two ways, depending on their attitude toward the characters:


\begin{description}
\item[Friendly or Indifferent Rivals.]
The rivals allow the characters to deal with Alyxian in whatever manner they see fit, fighting alongside them if need be.
\item[Hostile Rivals.]
The rivals attack the characters.
\end{description}

\subsection{Locations in the Heart}
The following locations are keyed to the Heart of Despair map.

\subsubsection{H1: Statue of Alyxian}
This painted stone statue depicts Alyxian as a young man. See "Alyxian Speaks" below for more information about the statue.

\subsubsection{H2–H4: Prayer Sites}
The ruins on the three islands are prayer sites dedicated to Avandra the Change Bringer (area H2), Corellon the Arch Heart (area H3), and Sehanine the Moon Weaver (area H4). The Apotheon draws on the power of these ruins when he takes his lair actions, as described in the "Lair Actions" section.

\subsection{Alyxian Speaks}
Before the characters begin exploring the Heart of Despair, read:

\begin{DndReadAloud}{}
The Apotheon's voice emanates from the statue and says, "After so long, we can finally speak to one another, and I am overjoyed. You have seen what I have done. What I have become. You, better than anyone, have the clarity of sight to judge me. May I leave?"

\end{DndReadAloud}

The Apotheon's inviting words should give the characters an uneasy, hesitant feeling. Alyxian believes they are the only beings alive who understand him well enough to judge him worthy of returning to the world. Whatever their decision, it must be unanimous (see "Deciding Alyxian's Fate" below).

The statue contains Alyxian's spirit. It is an object with AC 18, 50 hit points, and immunity to poison and psychic damage. If the characters destroy the statue, jump ahead to the "Battle for the Soul of the Apotheon" section.

\subsubsection{Deciding Alyxian's Fate}
If the characters grant the Apotheon his freedom, he gasps in surprise. This is the outcome he wants, but his own self-loathing did not prepare him for this answer, since deep down he knows he's in no condition to leave. He asks the characters if they are sure of their decision, whereupon any character who succeeds on a DC 16 Wisdom (Insight) check can sense the doubt that caused the Apotheon to question their decision. If the characters choose to free him at this point, proceed with the "Worst Ending: Unleashing Devastation" section.

Alyxian is willing to chat amicably with the characters and answer any questions they have—up to a point. Freedom is within his sight after eons of suffering. He is impatient, restless, and determined to end this chapter of his existence. He keeps the conversation going as long as he believes he's convincing the characters that he's a good person. If the characters don't come to a decision or start to question his assertions, frustration starts to creep into his tone, then anger.

If the characters ask about certain topics, you can read or paraphrase Alyxian's replies. Whenever he gives an answer, a character can make a DC 21 Wisdom (Insight) check. On a failed check, the character views him as sincere, sorrowful, and grateful. On a successful check, the character can tell he's deeply conflicted. On a check that succeeds by 5 or more, the character can also tell that he is growing impatient for the characters to judge him worthy. The following are the topics and Alyxian's replies:


\begin{description}
\item[Alyxian's Personality.]
"I have been cursed by fate. By some cruel joke of the cosmos. If I am to suffer, was it not right that I should allow my suffering to spare others from the same fate?"
\item[Alyxian's Past.]
"I was born as the Calamity raged around me. I grew up knowing fear, and I watched the Betrayer Gods turn my home into a barren wasteland. My only regrets are that I could not save others... and that my sacrifices were not enough to free me from my fate."
\item[Alyxian's Future.]
"I have dreamed of my future for so long, but now that you ask, I find my dreams have abandoned me. I want to walk by myself along the streets of a city where no one has cause to fear or loathe me."
\item[Alyxian's Regret.]
"You have felt my remorse and witnessed all the memories that I have prayed to be rid of. I am grateful for your actions in the retelling of those memories. I have watched them so many times, I wept to see that they could have turned out happily."
\item[Alyxian's Fury.]
"I am filled with hatred. For the god who condemned me to this hell, and for the people who tear at my essence to satisfy their own greed. Does that scare you?" (The god Alyxian refers to is Gruumsh. The people Alyxian refers to are the factions who want to take ruidium from the Netherdeep and use it for their own ends.)
\item[Alyxian's Yearning.]
"I have been imprisoned for so long, I dreamed of many things. Those dreams gave me hope, but—but now you are here! They needn't be fantasies any longer."
\item[The Heart of Despair.]
"It defies belief that this beautiful place is the source of so much suffering. Let us not speak of this place any longer. I simply wish to be free of it."
\item[Ruidium.]
"Do not speak to me of that wretched stuff. I would sooner slay anyone who takes it from my domain than see it used again. Let us be rid of this evil place and deal with ruidium later."
\end{description}

\subsubsection{Upsetting the Apotheon}
If a character points out that Alyxian seems angry during their conversation, he hastily retorts, "I am impatient. The world above sings to me, and you withhold judgment from me. What is your decision? Am I free to go? Choose!"

If any option other than unconditional release from his prison is suggested, if the characters take too long to decide, or if the characters doubt him openly, the Apotheon turns hostile. Read:

\begin{DndReadAloud}{}
The Heart of Despair rumbles. The water churns as if in anger, and the sounds of breaking glass and cracking stone pierce the space as the Apotheon speaks. "I have prayed for freedom. I have begged for release. I have screamed to an uncaring world and the impassive heavens for salvation. Yet even here, with you before me, I am denied! No more waiting. I will claim everything I have waited these long centuries for, and you will be the ones the world forgot."

\end{DndReadAloud}

Proceed with the next section.

\section{Battle for the Soul of the Apotheon}
Angering the Apotheon or denying him his release triggers a three-stage battle that begins as follows:

\begin{DndReadAloud}{}
The statue of Alyxian melts into a mass of pulsing black ichor. It fills your nose with a putrid scent, like a combination of iron and rotted meat. Thunder booms as crimson storm clouds blot out the sky. Bright red light flares from the ruined prayer sites, which are crawling with tendrils of ruidium that glow with malevolent power. Silence hangs over the lake for a moment.

Then the silence breaks. A milky-white monster with dozens of arms, too many blinking eyes, and a thick hide festooned with blades bursts out of the inky ooze. It rears back, then crashes down onto the lake as it howls in vengeance.

\end{DndReadAloud}

The battle with the Apotheon begins.

\subsection{How to Run the Encounter}
The following sections describe how the battle against the Apotheon is expected to play out.

\subsubsection{Multiple Stages}
The Apotheon has multiple forms that the characters encounter in sequence, each with its own stat block. Three of these forms—\textbf{Alyxian the Tormented}, \textbf{Alyxian the Callous}, and \textbf{Alyxian the Dispossessed}—are hostile. The Apotheon has a fourth form—\textbf{Alyxian the Absolved}—that poses no threat to the characters and appears only if the Apotheon is redeemed. The sequence is as follows:

Stage 1. The battle begins with the appearance of Alyxian the Tormented. When this form is reduced to 0 hit points, it disappears and is replaced by Alyxian the Callous.

Stage 2. When Alyxian the Callous is reduced to 0 hit points, this form of the Apotheon disappears and is replaced by Alyxian the Dispossessed.

Stage 3. When Alyxian the Dispossessed is reduced to 0 hit points, the Apotheon is well and truly dead. If Alyxian the Dispossessed is redeemed instead of killed, he transforms into Alyxian the Absolved. Either way, the battle ends.

\subsubsection{Emotional Healing}
While in the form of Alyxian the Tormented, Alyxian the Callous, or Alyxian the Dispossessed, the Apotheon taunts the characters, but an underlying sadness and regret hint that something humane still lives within him—a fragment of his original psyche that is desperate to curb his destructive tendencies yet unable to break free of the torment that drives his actions. By reaching out with words of empathy and compassion, the characters can help Alyxian achieve a kind of catharsis, leading to his redemption.

Have Alyxian speak to the characters often, to convey that this is not just a drawn-out combat encounter but an opportunity to rid Alyxian of his corruption. During the battle, the characters can improve Alyxian's emotional state by speaking to him with supportive words and succeeding on one or more Charisma (Persuasion) checks. With each successful check, they weaken his current form and push him toward assuming a less monstrous form. See the various stages of the battle, described below, for how to resolve these checks.

\subparagraph{Placating the Apotheon}
Count the number of times the characters' actions in chapter 6 placated the Apotheon. Then share this number with the players. This number indicates how many times the characters can gain advantage on their Charisma (Persuasion) checks to improve Alyxian's emotional state.

\subsection{Lair Actions}
Rising from the lake are fallen temples to Sehanine the Moon Weaver, Avandra the Change Bringer, and Corellon the Arch Heart, imagined into being by the despairing and lonely Apotheon. These structures stand atop 20-foot-tall rocky isles. Regardless of his current form, Alyxian can draw on the power of these profane ruins to take lair actions. Each lair action is associated with a particular site; when a lair action is used, the corresponding ruin glows with bright crimson light that lingers until initiative count 20 of the next round.

On initiative count 20 (losing initiative ties), Alyxian can take one of the following lair actions; he can't take the same lair action two rounds in a row:


\begin{description}
\item[Avandra's Grasp.]
Alyxian uses the power of Avandra's ruined temple (area H2) to create a watery tendril that rises out of the lake. One creature Alyxian can see that's not submerged in the lake must succeed on a DC 15 Dexterity saving throw or be restrained by the tendril. The tendril is a Large object that has AC 10, 15 hit points, and immunity to poison and psychic damage. The tendril disappears when it is reduced to 0 hit points or the next time Alyxian takes a lair action.
\item[Corellon's Charm.]
Alyxian causes a magical wisp of light to emerge from Corellon's ruined temple (area H3). The wisp enters the body of one creature Alyxian can see that isn't inside Corellon's ruined temple or atop the rocky island that supports it. The creature must succeed on a DC 15 Wisdom saving throw or be charmed by Alyxian until initiative count 20 of the next round.
\item[Sehanine's Torment.]
Alyxian uses the power of Sehanine's ruined temple (area H4) to create a watery form in the space of one creature that is either standing on the lake or submerged in it. The watery form assumes the vague likeness of the creature whose space it occupies and makes one melee weapon attack (+7 to hit) against it. On a hit, the attack deals 8 (1d8 + 4) bludgeoning damage. Whether it hits or misses, the watery form disappears after making its attack.
\end{description}

\subparagraph{Prayers to the Gods}
A character can use an action to pray to Avandra, Corellon, or Sehanine in the deity's prayer site. The character must succeed on a DC 15 Wisdom (Religion) check, and the check has advantage if the character is a worshiper of that god. On a successful check, the profane ruins atop the island melt away, and Alyxian can no longer use the lair action associated with that prayer site. In addition, Alyxian takes 14 (4d6) psychic damage, regardless of his current form.

\subsection{Stage 1: Alyxian the Tormented}
Alyxian's first form is gigantic and monstrous, for it represents the most corrupted version of the Apotheon. All the pain that torments his soul is made manifest upon his flesh to such an extent that it eclipses the gentle soul within. In this form, he uses the accompanying Alyxian the Tormented stat block.

\subsubsection{Emotional State}
Devoid of virtually every attribute that once made him a hero, Alyxian the Tormented is nihilistic, despondent, and wrathful. He attaches the worst possible interpretation to what the characters say, believes them to be his enemies, rebuffs their attempts to help, and mocks them at every turn.

The following bits of dialogue are examples of what he might say in this emotional state:


\begin{itemize}
\item "You can't understand. No one can understand the pain of spending years beyond counting in darkness."
\item "No one remembers. No one cares. Even you!"
\item "Even you will be forgotten. Do you trust them? Your friends? Do not trust. Do not grow close."
\end{itemize}

\subsubsection{Comforting Alyxian the Tormented}
As an action, a character can speak to Alyxian the Tormented with empathy and compassion, after which the character can make a DC 15 Charisma (Persuasion) check. On a failed check, the character's words do not resonate with Alyxian. On a successful check, Alyxian the Tormented loses 20 hit points as he tries to shed his monstrous form, causing chunks of his body to slough off. As gruesome as this sight might seem, Alyxian's emotional state appears to improve as a result. The following bits of dialogue are examples of what he might say in response to the characters' helpful words:


\begin{itemize}
\item "How can I forgive them for this suffering?"
\item "Of course, you can say such things! You have had family, friends, allies, and each other."
\item "This monstrous form... you must destroy it. I can feel its hatred tearing at me. Please! I can feel my spirit inside straining to escape!"
\end{itemize}

\subsection{Stage 2: Alyxian the Callous}
When Alyxian the Tormented is reduced to 0 hit points, read or paraphrase the following text as his second form sprouts from the husk of his first one:

\begin{DndReadAloud}{}
The corpse splits open like a cocoon, and a radiant figure with angelic wings and wreathed in a halo of spears emerges from the husk. With a gleaming spear in his hand, he looks down on you as if in judgment. In a cold voice he says, "Your words mean nothing. I am the Apotheon, born of the blessing of three gods. Thousands died by my hand. You will be no different."

\end{DndReadAloud}

To spare himself the pain that comes with feeling, Alyxian has hidden his emotions beneath the expressionless mask of a killer. In this form, he uses the Alyxian the Callous stat block.

\subsubsection{Emotional State}
Alyxian the Callous is cold and distant, disregarding words he doesn't care to hear. He sees himself as a perfect, "holy" being and believes that the world should be condemned for its crimes against him.

The following bits of dialogue are examples of what he might say in this emotional state:


\begin{itemize}
\item "What crime have I committed other than craving freedom? Why do you not condemn those who left me to rot?"
\item "It is as I thought: you are imperfect judges. You have not known my suffering. No one can know it; I am my only judge."
\item "I have lived a thousand lifetimes and felt a million sorrows. I am the judge this world deserves—and mine will be the final decree."
\end{itemize}

\subsubsection{Comforting Alyxian the Callous}
As an action, a character can speak to Alyxian the Callous with empathy and compassion, after which the character can make a DC 17 Charisma (Persuasion) check. On a failed check, the character's words do not resonate with Alyxian. On a successful check, Alyxian the Callous loses 25 hit points as his current form withers. As painful as this effect might seem, Alyxian's emotional state appears to improve as a result.

The following bits of dialogue are examples of what he might say in response to the characters' helpful words:


\begin{itemize}
\item "You do have a way with words, puny one."
\item "If the world will not pay for the hurt it caused me, then what? What happens?"
\item "How can this be righteousness? How is any of this allowed by fate, by gods that once fought by my side?"
\item "If there is no justice in this world, then the least you can do is grant me an honorable defeat. Come forth! Let us waste no more time on words."
\end{itemize}

\subsection{Stage 3: Alyxian the Dispossessed}
When Alyxian the Callous is reduced to 0 hit points, read or paraphrase the following text as his third form erupts from the remains of his second one:

\begin{DndReadAloud}{}
The Apotheon's angelic form falls, striking the ground in a golden starburst. Its gleaming feathers crumble to ash as Alyxian rises yet again. Before you, hunched and weary, is an old, haggard man. His gray hair falls over his face in loose strands, and he snarls as he readies his spear and his shield.

\end{DndReadAloud}

His heroic countenance shattered, Alyxian's radiant form sloughs off, revealing an old man. In this form, he uses the Alyxian the Dispossessed stat block.

\subsubsection{Emotional State}
Alyxian the Dispossessed is tired. The weariness that has built up in him over the centuries has begun to overcome his anger and bitterness. Now he needs to be guided back into the light.

The following bits of dialogue are examples of what he might say in this emotional state:


\begin{itemize}
\item "This world has moved beyond me. There is no place in it for a bitter man such as I."
\item "How could the gods ever forgive me for what I have done? I'm too old for redemption."
\item "Your naivete shows itself. The omens of Ruidus haunt me now as they did many moons ago."
\end{itemize}

\subsubsection{Comforting Alyxian the Dispossessed}
As an action, a character can speak to Alyxian the Dispossessed with empathy and compassion, after which the character can make a DC 19 Charisma (Persuasion) check. If a character can think of another way to remind Alyxian of the good man he could still be, have the character make a DC 19 ability check using whichever ability and skill you think is appropriate. On a failed check, the character's words do not resonate with Alyxian. On a successful check, years of negative emotions seem to drain away from him. After three successful checks, Alyxian the Dispossessed achieves catharsis and transforms into Alyxian the Absolved (see "Alyxian Redeemed" below).

The following bits of dialogue are examples of what he might say as Alyxian the Absolved:


\begin{itemize}
\item "I have survived for this long. Surely it was for some purpose."
\item "I have heard echoes of the most beautiful things from the surface. Laughter. Music. Joy. I want to see them. To know them again."
\item "What is left of the world that I once knew?"
\end{itemize}

\subsection{Ending the Battle}
The climactic battle ends with Alyxian victorious, defeated, or redeemed. These outcomes are discussed below.

\subsubsection{Alyxian Victorious}
If Alyxian defeats the characters, he disappears with a cackle. His victory triggers the worst possible outcome, as described in the "Worst Ending: Unleashing Devastation" section later in the chapter.

When Alyxian disappears, all creatures inside in the Heart of Desire—alive or dead—are transported to unoccupied spaces in area N26 (see chapter 6), which is now an empty, flooded, spherical chamber.

\subsubsection{Alyxian Defeated}
If Alyxian the Dispossessed dies, any hope of the Apotheon's redemption dies with him. Read:

\begin{DndReadAloud}{}
Alyxian lets out a long, rattling sigh as ichor trickles from his wounds. He coughs, and a gout of the fluid spurts from his mouth. Ruidium veins creep up his neck, and he smiles ruefully. He drops to his knees, reduced to a shriveled, gaunt figure, and lets out a last whimper that is undeniably human.

"I have been so cold." The light leaves his eyes. "So cold and so alone." With that, he falls into the lake and sinks into the water.

\end{DndReadAloud}

After Alyxian dies, all creatures inside the Heart of Desire—alive or dead—are transported to unoccupied spaces in area N26, which is now an empty, flooded, spherical chamber. Proceed with the "Neutral Ending: Alyxian at Rest" section.

\subsubsection{Alyxian Redeemed}
Upon assuming the form of Alyxian the Absolved (see the accompanying stat block), the Apotheon no longer has any desire to fight the characters. Read:

\begin{DndReadAloud}{}
The Apotheon looks toward the sky. He takes in a long breath, then releases it, shuddering as tears run down his face. A smile of relief comes over his countenance.

"My soul feels so light," he whispers. He locks eyes with each of you, one by one, and says, "Please. Let us leave this cursed place."

\end{DndReadAloud}

If the characters unanimously grant permission for Alyxian the Absolved to leave his self-made prison, proceed with the "Best Ending: A World That Remembered" section. If they kill Alyxian the Absolved, the story takes the direction described in "Neutral Ending: Alyxian at Rest."

\section{Back into the Light}
The ending of this adventure tries to resolve any major questions that remain by providing explanations for what happened to the Apotheon, what happened to the rivals, and what happened to the factions. Other issues, including the fate of the characters and any nonplayer characters they grew attached to, are left for you and your players to resolve.

\subsection{Possible Endings}
How this story ends hinges on the outcome of the battle against Alyxian.

\subsubsection{Neutral Ending: Alyxian at Rest}
\textit{Alyxian is defeated instead of being redeemed.}

This ending applies if the characters, refusing to trust the Apotheon enough to allow him his freedom, destroy him during the final confrontation and thereby deprive the Netherdeep of its vital essence.

\subparagraph{Ruidium's End}
Upon Alyxian's death, all ruidium on Exandria melts away, leaving no trace of the substance anywhere in the world. Creatures suffering from ruidium corruption are instantly rid of it. Magic items that contain ruidium transform as described in appendix B.

\subparagraph{Escape from the Netherdeep}
Without the power of the Apotheon to sustain it, the Netherdeep begins to collapse. Creatures have 30 minutes to escape from the Netherdeep as it crumbles around them; those who fail to get out are lost. If the characters have a clear path back to area N1 (see chapter 6), escaping in this amount of time is doable. If the characters pass through one or more areas that contain hostile creatures, these creatures (which might include the rival adventurers) do their best to waylay the characters.

\subparagraph{Parting Vision}
Characters who escape from the Netherdeep receive a vision of the Apotheon. The characters see him with his back turned to them as he receives the killing blow from Gruumsh. Then the characters see themselves, one after the other, striking the same blow. The Apotheon's blood seeps into the sand and disappears from sight. Time passes, and sandstorms rage, then clear, revealing that flowers have begun sprouting from the blood-drenched sand. Then buildings begin to emerge, forming the city of Ank'Harel—which is now at peace thanks to their actions.

\subsubsection{Worst Ending: Unleashing Devastation}
\textit{Alyxian is granted freedom without regaining compassion.}

If the characters take pity on the Apotheon and grant him his freedom without healing his emotional wounds, read:

\begin{DndReadAloud}{}
The statue of Alyxian explodes, spraying debris and black ichor. The prayer sites, too, explode into rubble as the Apotheon siphons off their last vestiges of power. His voice resounds from everywhere as he screams, "Never again will they forget my name!" A titanic mass of oily, slick shadow gushes past you and tears a hole in the edge of the Heart of Despair.

A wave of darkness comes over you. You are left gasping for breath, but that is the least of your worries.

Ruin is coming to Exandria.

\end{DndReadAloud}

After escaping from the Heart of Despair, the Apotheon races through the Netherdeep and erupts in Ank'Harel. Darkness fills the sky above the city, as if it were a moonless, starless night—and then Ruidus's light blazes abruptly, filling the inky sky with crimson ire.

After ruminating over his regret, fury, yearning, and despair for an age, the Apotheon condemns the world that allowed him to be imprisoned after a lifetime of selflessness. His vengeance is all-consuming, his pain unending. For the next year and a day, Ruidus remains a full moon, shining with unnatural brightness. During this time, Alyxian can appear anywhere in the world in the form of an oily black cloud that can reduce cities to rubble in a matter of days.

\subparagraph{Escape from the Netherdeep}
Without the power of the Apotheon to sustain it, the Netherdeep begins to collapse (see "Neutral Ending: Alyxian at Rest" above). When the characters emerge from the Netherdeep, they find Ank'Harel ablaze, its people screaming and fleeing for safety. Alyxian's voice, distorted and hateful, booms out, "Soon, all shall remember!"

To continue a campaign in the midst of this devastation, you can devise a way to stop this wrathful avatar from razing all of Exandria, which could be the characters receiving blessings from other Prime Deities to compete against his power; raising an Exandrian army to face him and the occult worshipers who flock to him; or, for a party that is desperate or daring, seeking help from the Betrayer Gods (see "What of the Deities?" below).

\subsubsection{Best Ending: A World That Remembered}
\textit{Alyxian is shown kindness and reminded of love.}

When the Apotheon's corruption has been neutralized, the characters can speak to Alyxian's true self. They can remind him of the great deeds he has done and could still do, as well as express gratitude for his sacrifice on behalf of the world at large. As Alyxian sheds the last of his bitterness, pain, and fury, he becomes a humble man once more. In this form, he is truly ready to return to the world above.

\begin{DndSidebar}[float=!b]{Roleplaying Alyxian the Absolved}
If the characters save the Apotheon's tortured soul, Alyxian the Dispossessed fades away and is replaced by Alyxian the Absolved. Though he is not as he was before the horrors of the Netherdeep, this form of Alyxian has been relieved of his burdens. He has found freedom from years of anguish and regret. Alyxian the Absolved is neither an aberrant monster nor a divine avenger; he is, as he always was, a man seeking serenity.



Alyxian the Absolved can help the characters escape from the Netherdeep as it collapses, and he might even come to the party's aid in future adventures—though he would prefer to live out his remaining mortal days in peace. See "Best Ending: A World That Remembered" for more information.



\end{DndSidebar}

\subparagraph{Ruidium's End}
The Apotheon who caused the spread of ruidium is dead, replaced by a version of the Apotheon who is at peace with himself and the world. The destruction of Alyxian's corrupted forms causes all ruidium throughout Exandria to melt away, leaving no trace of the substance anywhere in the world. Creatures suffering from ruidium corruption are instantly rid of it. Magic items that contain ruidium transform as described in appendix B.

\subparagraph{Escape from the Netherdeep}
The Netherdeep reacts violently when the characters heal the scars on the Apotheon's soul. Creatures have 30 minutes to escape from the Netherdeep as it crumbles around them; those who fail to get out are lost. If the characters have a clear path back to area N1 (see chapter 6) and can breathe water, escaping in this length of time is a simple task. If the characters pass through one or more areas that contain hostile creatures, these creatures do their best to waylay the characters. Alyxian the Absolved casts \textit{water breathing} and fights by the characters' side in any such conflict. If the rivals are among these hostile creatures, they flee rather than face Alyxian in battle.

\subparagraph{Finale in the Light}
No longer hampered by his fury, his memories, and the empty silence of the Netherdeep, Alyxian gapes in wonderment at Ank'Harel. The busy city that stretches out before him is filled with pleasant chatter, birdsong, and sunlight. He flinches, overwhelmed, and feels out of place in this new world. The characters can either help him settle into his new life or leave him to his own devices. Either way, in time, he finds peace and makes new friends.

At your discretion, one of the Prime Deities most closely tied to Alyxian (Avandra, Corellon, or Sehanine) might reward each character with a blessing of your choice (see "\textit{Supernatural Gifts}" in the \textit{Dungeon Master's Guide}).

\subsection{What of the Factions?}
The nature of the characters' ongoing relationships with the factions of Ank'Harel depends on what becomes of Alyxian, what happens to ruidium, and what the characters do after emerging from the Netherdeep.

Characters who represented their faction vigorously and who provided resources (ruidium or otherwise) for their faction might be asked to assume positions of leadership. If either the story's neutral ending or best ending comes to pass, all ruidium is destroyed—an outcome that delights the Cobalt Soul, vexes the Allegiance of Allsight, and infuriates the Consortium of the Vermilion Dream.

\subsection{What of the Rivals?}
The characters' rivals have experienced their share of triumphs, tragedies, and trials. The following sections describe likely ending scenarios for the rivals, depending on their attitude toward the characters.

\subsubsection{Friendly Rivals}
If the rivals are friendly, their actions from here on depend on what happens to Alyxian:


\begin{description}
\item[Alyxian Is Killed.]
The rivals help the heroes escape from the Netherdeep. Sometime later, they invite the characters to a tavern in Ank'Harel to share their memories of the Netherdeep, then try to put it all behind them.
\item[Alyxian Is Unleashed.]
The rivals become staunch allies of the characters. The rivals agree to go wherever the party sends them to help stem the disaster.
\item[Alyxian Is Redeemed.]
The rivals, recognizing and accepting that the characters have prevailed, honor the heroes during a night of celebration together. They might become fast friends of the characters in time.
\end{description}

\subsubsection{Indifferent Rivals}
If the rivals are indifferent, their actions from here on depend on what happens to Alyxian:


\begin{description}
\item[Alyxian Is Killed.]
The rivals report what happened to their faction leaders, casting themselves as neither better nor worse than the characters. They take their leave in search of a place to build a new reputation for themselves.
\item[Alyxian Is Unleashed.]
The rivals blame the characters for the cataclysm now facing the world but realize that their only hope of survival is to band together against their common enemy.
\item[Alyxian Is Redeemed.]
The rivals thank the characters for engaging in the "race" with them, which didn't really seem to have a winner. Ayo remarks, "Next time there's something you want, we'll be there for you—maybe just to be a thorn in your side." Her parting words seem warm enough.
\end{description}

\subsubsection{Hostile Rivals}
If the rivals are hostile, their actions from here on depend on what happens to Alyxian:


\begin{description}
\item[Alyxian Is Killed.]
The rivals report to their faction that the characters have destroyed their access to ruidium and killed a legendary hero. They sabotage the characters' relationship with their own faction as much as they can and angle to become viewed as the more accomplished adventurers.
\item[Alyxian Is Unleashed.]
The rivals blame the characters for the devastation Alyxian causes. They complicate the characters' dealings with political leaders and nations, taking every opportunity to portray the characters as selfish, untrustworthy meddlers.
\item[Alyxian Is Redeemed.]
The rivals, realizing that their further efforts here would be to no avail, set out for distant lands to find a new way to make their mark in the world.
\end{description}

\subsubsection{Heroes Unsung}
If the rivals died in the Netherdeep, the characters are haunted by dreams of them for a few weeks. In these dreams, they might see Maggie's solemn, steady gaze as she speaks, but not be able to hear what she is trying to tell them. They might see Ayo erupting in an outburst of anger, Galsariad fruitlessly trying to turn back time to change his party's fate, Dermot mourning his friends, and Irvan wondering if he ever did something great.

Eventually, the dreams subside, but on their final night, Ayo steps up to each character in turn and offers a salute. "We didn't do too badly, huh?" she says quietly. "Remember us. Maybe with a monument back in Jigow? That would be nice."

\subsection{What of the Deities?}
This adventure has chronicled Alyxian's relationship with three Prime Deities: Avandra, Corellon, and Sehanine. By the end of the adventure, the characters are heroes in their own right—and they have had opportunities to act in the name of the same deities Alyxian once heralded. If Alyxian threatens to bring doom to the world, these gods might send celestial envoys to entreat directly with the characters:


\begin{description}
\item[Avandra the Change Bringer.]
Avandra might reach out to paladins, druids, clerics, wizards, the adventurous, and the bold. She encourages her faithful to grasp their fate, embrace change for the better, and rise against the threat posed by Alyxian. If the characters need her aid, her solars and planetars might deliver useful information about places they hope to visit, or Avandra might provide the occasional swing of fate in the characters' favor.
\item[Corellon the Arch Heart.]
Corellon might reach out to elves, spellcasters, the artistic, and the clever. Although Corellon offers direct support less often than other deities, their servants or the strange visions Corellon provides often coax adventurers into searching for forgotten secrets and spells, empowering these heroes through the fruits of discovery and inspiration. Arcane tools of this sort can be useful in overcoming any future threat.
\item[Sehanine the Moon Weaver.]
Sehanine might reach out to bards, rogues, druids, clerics, and the mischievous. Her childlike form appears to such individuals in dreams. If the characters need her aid, she might send the deva Perigee to offer them support or a place of refuge.
\end{description}

\subparagraph{Betrayer Gods}
The Betrayer Gods might tempt adventurers (either the same group who freed Alyxian or a new band of heroes) into forming an alliance against Alyxian if the Apotheon can't be subdued any other way. The following are boons the Betrayer Gods might offer to a party—for a price:

\textbf{Asmodeus} might provide a legion of devils.

\textbf{Gruumsh and Bane} might provide armies of marauders and warriors.

\textbf{Lolth} might offer to sow dissension among the mortals that have begun to worship Alyxian, fostering chaos and rebellion to stem that tide.

\textbf{Tiamat} might offer to send an ancient chromatic dragon—or five—to aid in an assault or a siege at a critical moment.

\textbf{Torog} might send a horde of demons and purple worms to assail the fortifications of the characters' enemies.

\textbf{Vecna} might offer previously forbidden arcane power, gifting the characters with powerful magic items that will eventually corrupt them.

\textbf{Zehir} might send assassins to their aid, or he might furnish the characters with a poison that weakens anyone imbued with divinity.

Although any of these alliances are potentially beneficial in the short term, receiving any of these benefits will require the characters to partake in actions or rituals that strengthen the Betrayer Gods and their hold over Exandria. Then, once Alyxian is dealt with, their cults' domination of the world can begin.

\subsection{What of Our Heroes?}
When all is said and done, the characters are credited as the ones who brought peace to a tormented soul or blamed for unleashing that soul upon an unsuspecting world.

If the characters achieve some measure of renown, their reputation precedes them if they travel to faraway lands. They are approached by eager folk with questions about what lies below Ank'Harel. "Tell us of Cael Morrow," they'll say. As time passes, the characters might become known more broadly as "the Redeemers," "the Prayer Bearers," "the Hope Bringers," "the Soul Healers," or any one of a hundred other names that bring to mind the mythic heroes who plunged into the Netherdeep and experienced both triumph and tragedy.

\appendix
\chapter{Creatures}

\pagestyle{empty}

\setlist[description]{nolistsep,listparindent=3pt,topsep=6pt,font={\bfseries\sffamily\small}}

\begin{DndMonster}[float*=b,width=\textwidth + 8pt]{Alyxian the Absolved}
\begin{multicols}{2}
 \DndMonsterType{Medium humanoid (human), lawful good}
\DndMonsterBasics[
armor-class = {15 (\textit{leather armor})},
hit-points = {\DndDice{19d8 + 76}},
speed = {30 ft.}
]
\DndMonsterAbilityScores[
 str = 19,
 dex = 15,
 con = 18,
 int = 13,
 wis = 14,
 cha = 19
]
\DndMonsterDetails[
saving-throws = {Str +9, Con +9, Wis +7},
skills = {Athletics +9, Insight +7, Perception +7, Persuasion +9},
senses = {darkvision 120 ft., passive perception 17},
languages = {Common, Elvish},
challenge = 14,
]
\DndMonsterAction{Divinely Blessed}
Alyxian can't be Surprise and can't be changed into another form against his will.

\DndMonsterAction{Legendary Resistance (2/Day)}
If Alyxian fails a saving throw, he can choose to succeed instead.

\DndMonsterAction{Spellcasting}
Alyxian casts one of the following spells, requiring no material components and using Charisma as the spellcasting modifier (spell save DC 17):
\begin{DndMonsterSpells}
  \DndInnateSpellLevel{guidance, light}
  \DndInnateSpellLevel[1e]{lesser restoration, remove curse, water breathing}
\end{DndMonsterSpells}
\DndMonsterSection{Actions}
\DndMonsterAction{Multiattack}
Alyxian makes two Spear attacks.

\DndMonsterAction{Spear}
\textsl{Melee or Ranged Weapon Attack:} +9 to hit, reach 5 ft. or range 20/60 ft., one target. \textsl{Hit:} 7 (1d6 + 4) piercing damage, or 8 (1d8 + 4) piercing damage when used with two hands to make a melee attack, plus 9 (2d8) radiant damage.

\DndMonsterAction{Force Wave (1/Day)}
Alyxian strikes the ground, creating a burst of energy that ripples outward. Each creature he chooses within a 30-foot-radius sphere centered on himself must succeed on a DC 17 Constitution saving throw or take 35 (10d6) force damage and be knocked prone. A creature that succeeds on the saving throw takes half as much damage and isn't knocked prone.

\DndMonsterSection{Reactions}
\DndMonsterAction{Parry}
Alyxian adds 3 to his AC against one attack roll that would hit him. To do so, Alyxian must see the attacker and be wielding a melee weapon.

\DndMonsterSection{Legendary Actions}
Alyxian the Absolved can take 3 legendary actions, choosing from the options below. Only one legendary action can be used at a time and only at the end of another creature's turn. Alyxian the Absolved regains spent legendary actions at the start of its turn.

\begin{DndMonsterLegendaryActions}
\DndMonsterLegendaryAction{Arch Heart's Strike}{Alyxian makes one Spear attack, and all damage from this attack is radiant.
}
\DndMonsterLegendaryAction{Change Bringer's Dance}{Alyxian moves up to his speed. This movement doesn't provoke opportunity attacks.
}
\DndMonsterLegendaryAction{Moon Weaver's Veil (Costs 2 Actions)}{Alyxian becomes invisible until the end of his next turn. Any equipment he is wearing or carrying becomes invisible with him.
}
\end{DndMonsterLegendaryActions}

\end{multicols}
\end{DndMonster}



\begin{DndMonster}[float*=b,width=\textwidth + 8pt]{Alyxian the Callous}
\begin{multicols}{2}
\DndMonsterType{Large celestial, lawful evil}
\DndMonsterBasics[
armor-class = {17 (half plate armor)},
hit-points = {\DndDice{15d10 + 75}},
speed = {30 ft., fly 90 ft.}
]
\DndMonsterAbilityScores[
 str = 23,
 dex = 18,
 con = 21,
 int = 15,
 wis = 16,
 cha = 20
]
\DndMonsterDetails[
saving-throws = {Str +10, Con +9, Wis +7},
skills = {Deception +9, Insight +7, Perception +7},
damage-immunities = {radiant},
senses = {truesight 120 ft., passive perception 17},
languages = {Celestial, Common, Elvish, telepathy 120 ft.},
challenge = 12,
]
\DndMonsterAction{Apotheonic Rejuvenation}
When Alyxian the Callous drops to 0 hit points, his body dies and sheds its wings, and he rises to his feet in his third form, Alyxian the Dispossessed. His initiative count doesn't change.

\DndMonsterAction{Divinely Blessed}
Alyxian can't be Surprise and can't be changed into another form against his will.

\DndMonsterAction{Legendary Resistance (2/Day)}
If Alyxian fails a saving throw, he can choose to succeed instead.

\DndMonsterSection{Actions}
\DndMonsterAction{Multiattack}
Alyxian makes two Radiant Spear attacks. He can replace one of these attacks with Blinding Teleport.

\DndMonsterAction{Radiant Spear}
\textsl{Melee or Ranged Weapon Attack:} +10 to hit, reach 5 ft. or range 120 ft., one target. \textsl{Hit:} 15 (2d8 + 6) radiant damage.

\DndMonsterAction{Blinding Teleport}
Alyxian teleports, along with any equipment he is wearing or carrying, to an unoccupied space he can see within 120 feet of himself. Each creature within 5 feet of his new location must succeed on a DC 17 Constitution saving throw or be blinded until the end of its next turn.

\DndMonsterAction{Heavenly Destruction (1/Day)}
Alyxian releases divine energy in a 60-foot cone. Each creature of his choice in that area must make a DC 17 Constitution saving throw. On a failed saving throw, the creature takes 31 (7d8) radiant damage and is knocked prone. On a successful save, a creature takes half as much damage and isn't knocked prone.

\DndMonsterSection{Legendary Actions}
Alyxian the Callous can take 3 legendary actions, choosing from the options below. Only one legendary action can be used at a time and only at the end of another creature's turn. Alyxian the Callous regains spent legendary actions at the start of its turn.

\begin{DndMonsterLegendaryActions}
\DndMonsterLegendaryAction{Attack}{Alyxian makes one Radiant Spear attack.
}
\DndMonsterLegendaryAction{Celestial Chains}{Alyxian targets one creature he can see within 120 feet of himself. The target must succeed on a DC 17 Strength saving throw or be restrained by magical chains for 1 minute. While restrained in this way, the target can't leave the space by any means. A creature can repeat the saving throw at the end of each of its turns, ending the effect on itself on a success.
}
\DndMonsterLegendaryAction{Flare (Costs 2 Actions)}{Alyxian targets one creature he can see within 30 feet of himself. Light flares around the target, dealing 17 (5d6) radiant damage to it and creatures within 10 feet of it.
}
\end{DndMonsterLegendaryActions}
\DndMonsterSection{Lair Actions}
On initiative count 20 (losing initiative ties), Alyxian can take one of the following lair actions; he can't take the same lair action two rounds in a row:


\begin{description}
\item[Avandra's Grasp.]
Alyxian uses the power of Avandra's ruined temple (area H2) to create a watery tendril that rises out of the lake. One creature Alyxian can see that's not submerged in the lake must succeed on a DC 15 Dexterity saving throw or be restrained by the tendril. The tendril is a Large object that has AC 10, 15 hit points, and immunity to poison and psychic damage. The tendril disappears when it is reduced to 0 hit points or the next time Alyxian takes a lair action.

\item[Corellon's Charm.]
Alyxian causes a magical wisp of light to emerge from Corellon's ruined temple (area H3). The wisp enters the body of one creature Alyxian can see that isn't inside Corellon's ruined temple or atop the rocky island that supports it. The creature must succeed on a DC 15 Wisdom saving throw or be charmed by Alyxian until initiative count 20 of the next round.

\item[Sehanine's Torment.]
Alyxian uses the power of Sehanine's ruined temple (area H4) to create a watery form in the space of one creature that is either standing on the lake or submerged in it. The watery form assumes the vague likeness of the creature whose space it occupies and makes one melee weapon attack (+7 to hit) against it. On a hit, the attack deals 8 (1d8 + 4) bludgeoning damage. Whether it hits or misses, the watery form disappears after making its attack.

\end{description}

\DndMonsterSection{Prayers to the Gods}
A character can use an action to pray to Avandra, Corellon, or Sehanine in the deity's prayer site. The character must succeed on a DC 15 Wisdom (Religion) check, and the check has advantage if the character is a worshiper of that god. On a successful check, the profane ruins atop the island melt away, and Alyxian can no longer use the lair action associated with that prayer site. In addition, Alyxian takes 14 (4d6) psychic damage, regardless of his current form.

\end{multicols}
\end{DndMonster}



\begin{DndMonster}[float*=b,width=\textwidth + 8pt]{Alyxian the Dispossessed}
\begin{multicols}{2}
\DndMonsterType{Medium humanoid (human), lawful evil}
\DndMonsterBasics[
armor-class = {15 (\textit{leather armor})},
hit-points = {\DndDice{18d8 + 72}},
speed = {30 ft.}
]
\DndMonsterAbilityScores[
 str = 19,
 dex = 15,
 con = 18,
 int = 13,
 wis = 14,
 cha = 20
]
\DndMonsterDetails[
saving-throws = {Str +9, Con +9, Wis +7},
skills = {Insight +7, Perception +7, Persuasion +10},
senses = {darkvision 120 ft., passive perception 17},
languages = {Celestial, Common, Elvish, telepathy 120 ft.},
challenge = 13,
]
\DndMonsterAction{Divinely Blessed}
Alyxian can't be Surprise and can't be changed into another form against his will.

\DndMonsterAction{Legendary Resistance (2/Day)}
If Alyxian fails a saving throw, he can choose to succeed instead.

\DndMonsterAction{Special Equipment}
Alyxian carries a magic dagger that he uses to make Stoneheart Dagger attacks. In the hands of creatures other than Alyxian, the dagger is nonmagical and has no special properties.

\DndMonsterSection{Actions}
\DndMonsterAction{Multiattack}
Alyxian makes two Spear attacks.

\DndMonsterAction{Spear}
\textsl{Melee or Ranged Weapon Attack:} +9 to hit, reach 5 ft. or range 20/60 ft., one target. \textsl{Hit:} 7 (1d6 + 4) piercing damage, or 8 (1d8 + 4) piercing damage when used with two hands to make a melee attack, plus 9 (2d8) radiant damage.

\DndMonsterAction{Stoneheart Dagger}
\textsl{Melee or Ranged Weapon Attack:} +9 to hit, reach 5 ft. or range 20/60 ft., one target. \textsl{Hit:} 11 (3d4 + 4) force damage, and if the target is a creature, it must succeed on a DC 17 Constitution saving throw or become petrified as the dagger lodges itself in the target's body. While the dagger is lodged in the target, magic can't end the petrified condition on it and Alyxian can't make Stoneheart Dagger attacks. A creature within reach of the petrified target can use an action to try to remove the dagger, doing so with a successful DC 17 Strength check.

\DndMonsterSection{Bonus Actions}
\DndMonsterAction{Summon Dagger}
Alyxian's Stoneheart Dagger teleports into his free hand, provided he has one.

\DndMonsterSection{Reactions}
\DndMonsterAction{Parry}
Alyxian adds 3 to his AC against one attack roll that would hit him. To do so, Alyxian must see the attacker and be wielding a melee weapon.

\DndMonsterSection{Legendary Actions}
Alyxian the Dispossessed can take 3 legendary actions, choosing from the options below. Only one legendary action can be used at a time and only at the end of another creature's turn. Alyxian the Dispossessed regains spent legendary actions at the start of its turn.

\begin{DndMonsterLegendaryActions}
\DndMonsterLegendaryAction{Attack}{Alyxian makes one Spear or Stoneheart Dagger attack.
}
\DndMonsterLegendaryAction{Warrior's Stride}{Alyxian moves up to his speed. This movement doesn't provoke opportunity attacks.
}
\DndMonsterLegendaryAction{Ruidium Shard (Costs 2 Actions)}{A gem-sized shard of ruidium appears at a point Alyxian can see within 120 feet of himself and explodes. Each creature, other than Alyxian, in a 10-foot-radius sphere centered on that point must succeed on a DC 18 Charisma saving throw or gain 1 level of exhaustion. If a creature isn't suffering from ruidium corruption, it becomes corrupted when it fails the saving throw.
}
\end{DndMonsterLegendaryActions}
\DndMonsterSection{Lair Actions}
On initiative count 20 (losing initiative ties), Alyxian can take one of the following lair actions; he can't take the same lair action two rounds in a row:


\begin{description}
\item[Avandra's Grasp.]
Alyxian uses the power of Avandra's ruined temple (area H2) to create a watery tendril that rises out of the lake. One creature Alyxian can see that's not submerged in the lake must succeed on a DC 15 Dexterity saving throw or be restrained by the tendril. The tendril is a Large object that has AC 10, 15 hit points, and immunity to poison and psychic damage. The tendril disappears when it is reduced to 0 hit points or the next time Alyxian takes a lair action.

\item[Corellon's Charm.]
Alyxian causes a magical wisp of light to emerge from Corellon's ruined temple (area H3). The wisp enters the body of one creature Alyxian can see that isn't inside Corellon's ruined temple or atop the rocky island that supports it. The creature must succeed on a DC 15 Wisdom saving throw or be charmed by Alyxian until initiative count 20 of the next round.

\item[Sehanine's Torment.]
Alyxian uses the power of Sehanine's ruined temple (area H4) to create a watery form in the space of one creature that is either standing on the lake or submerged in it. The watery form assumes the vague likeness of the creature whose space it occupies and makes one melee weapon attack (+7 to hit) against it. On a hit, the attack deals 8 (1d8 + 4) bludgeoning damage. Whether it hits or misses, the watery form disappears after making its attack.

\end{description}

\DndMonsterSection{Prayers to the Gods}
A character can use an action to pray to Avandra, Corellon, or Sehanine in the deity's prayer site. The character must succeed on a DC 15 Wisdom (Religion) check, and the check has advantage if the character is a worshiper of that god. On a successful check, the profane ruins atop the island melt away, and Alyxian can no longer use the lair action associated with that prayer site. In addition, Alyxian takes 14 (4d6) psychic damage, regardless of his current form.

\end{multicols}
\end{DndMonster}



\begin{DndMonster}[float=t]{Alyxian the Hunter}
\DndMonsterType{Medium construct, unaligned}
\DndMonsterBasics[
armor-class = {17 (natural armor)},
hit-points = {\DndDice{17d10 + 85}},
speed = {30 ft., swim 30 ft.}
]
\DndMonsterAbilityScores[
 str = 22,
 dex = 9,
 con = 20,
 int = 3,
 wis = 11,
 cha = 1
]
\DndMonsterDetails[
damage-immunities = {poison, psychic; bludgeoning, piercing, slashing from nonmagical attacks that aren't adamantine},
senses = {darkvision 120 ft., passive perception 10},
languages = {understands the languages of its creator but can't speak},
challenge = 10,
]
\DndMonsterAction{Immutable Form}
The golem is immune to any spell or effect that would alter its form.

\DndMonsterAction{Magic Resistance}
The golem has advantage on saving throws against spells and other magical effects.

\DndMonsterAction{Magic Weapons}
The golem's weapon attacks are magical.

\DndMonsterSection{Actions}
\DndMonsterAction{Multiattack}
The golem makes two slam attacks.

\DndMonsterAction{Ruidium Dart}
The golem can conjure a magic dart made of ruidium and hurl it at another creature it can see within 120 feet of itself. The dart strikes the target unerringly and vanishes on contact. The target must make a DC 17 Charisma saving throw. On a failed saving throw, the target takes 24 (7d6) psychic damage and gains 1 level of exhaustion. In addition, if the target is not already suffering from ruidium corruption, it becomes corrupted when it fails the saving throw. On a successful saving throw, the target takes half as much damage and suffers no other effects.

\DndMonsterAction{Slam}
\textsl{Melee Weapon Attack:} +10 to hit, reach 5 ft., one target. \textsl{Hit:} 19 (3d8 + 6) bludgeoning damage.

\DndMonsterAction{Slow (Recharge 5\textendash 6)}
The golem targets one or more creatures it can see within 10 feet of it. Each target must make a DC 17 Wisdom saving throw against this magic. On a failed save, a target can't use reactions, its speed is halved, and it can't make more than one attack on its turn. In addition, the target can take either an action or a bonus action on its turn, not both. These effects last for 1 minute. A target can repeat the saving throw at the end of each of its turns, ending the effect on itself on a success.

\end{DndMonster}



\begin{DndMonster}[float*=b,width=\textwidth + 8pt]{Alyxian the Tormented}
\begin{multicols}{2}
\DndMonsterType{Huge aberration, lawful evil}
\DndMonsterBasics[
armor-class = {16 (natural armor)},
hit-points = {\DndDice{12d12 + 72}},
speed = {40 ft., swim 40 ft.}
]
\DndMonsterAbilityScores[
 str = 21,
 dex = 18,
 con = 22,
 int = 15,
 wis = 16,
 cha = 20
]
\DndMonsterDetails[
saving-throws = {Str +9, Con +10, Wis +7},
skills = {Deception +9, Insight +7, Perception +7},
damage-immunities = {necrotic},
senses = {blindsight 120 ft., passive perception 17},
languages = {Celestial, Common, Elvish, telepathy 120 ft.},
challenge = 11,
]
\DndMonsterAction{Apotheonic Rejuvenation}
When Alyxian drops to 0 hit points, his body dies and cracks open like a cocoon. Alyxian instantly emerges from the cocoon in his second form, Alyxian the Callous, in the space where his previous form died. His initiative count doesn't change.

\DndMonsterAction{Divinely Blessed}
Alyxian can't be Surprise and can't be changed into another form against his will.

\DndMonsterAction{Legendary Resistance (2/Day)}
If Alyxian fails a saving throw, he can choose to succeed instead.

\DndMonsterSection{Actions}
\DndMonsterAction{Multiattack}
Alyxian makes two Grasping Hand attacks.

\DndMonsterAction{Grasping Hand}
\textsl{Melee Weapon Attack:} +9 to hit, reach 5 ft., one target. \textsl{Hit:} 15 (3d6 + 5) force damage, and if the target is a Medium or smaller creature, it is grappled (escape DC 15). Alyxian has six hands, each of which can grapple one target.

\DndMonsterAction{Void Eyes (Recharge 5\textendash 6)}
Alyxian targets up to two creatures he can see within 120 feet of himself. Each target must make a DC 17 Constitution saving throw, taking 28 (8d6) necrotic damage on a failed save, or half as much damage on a successful one.

\DndMonsterSection{Legendary Actions}
Alyxian the Tormented can take 3 legendary actions, choosing from the options below. Only one legendary action can be used at a time and only at the end of another creature's turn. Alyxian the Tormented regains spent legendary actions at the start of its turn.

\begin{DndMonsterLegendaryActions}
\DndMonsterLegendaryAction{Attack}{Alyxian makes one Grasping Hand attack.
}
\DndMonsterLegendaryAction{Weapons of Yore (Costs 2 Actions)}{Ancient weapons embedded in Alyxian's hide detach from Alyxian and fly about in a 10-foot-radius sphere centered on a point Alyxian can see within 60 feet of himself. A creature takes 18 (4d8) slashing damage when it enters that area for the first time on a turn or starts its turn there. The effect lasts until the end of Alyxian's next turn, when the weapons return to Alyxian and embed themselves in his hide once again (causing him no harm).
}
\end{DndMonsterLegendaryActions}
\DndMonsterSection{Lair Actions}
On initiative count 20 (losing initiative ties), Alyxian can take one of the following lair actions; he can't take the same lair action two rounds in a row:


\begin{description}
\item[Avandra's Grasp.]
Alyxian uses the power of Avandra's ruined temple (area H2) to create a watery tendril that rises out of the lake. One creature Alyxian can see that's not submerged in the lake must succeed on a DC 15 Dexterity saving throw or be restrained by the tendril. The tendril is a Large object that has AC 10, 15 hit points, and immunity to poison and psychic damage. The tendril disappears when it is reduced to 0 hit points or the next time Alyxian takes a lair action.

\item[Corellon's Charm.]
Alyxian causes a magical wisp of light to emerge from Corellon's ruined temple (area H3). The wisp enters the body of one creature Alyxian can see that isn't inside Corellon's ruined temple or atop the rocky island that supports it. The creature must succeed on a DC 15 Wisdom saving throw or be charmed by Alyxian until initiative count 20 of the next round.

\item[Sehanine's Torment.]
Alyxian uses the power of Sehanine's ruined temple (area H4) to create a watery form in the space of one creature that is either standing on the lake or submerged in it. The watery form assumes the vague likeness of the creature whose space it occupies and makes one melee weapon attack (+7 to hit) against it. On a hit, the attack deals 8 (1d8 + 4) bludgeoning damage. Whether it hits or misses, the watery form disappears after making its attack.

\end{description}

\DndMonsterSection{Prayers to the Gods}
A character can use an action to pray to Avandra, Corellon, or Sehanine in the deity's prayer site. The character must succeed on a DC 15 Wisdom (Religion) check, and the check has advantage if the character is a worshiper of that god. On a successful check, the profane ruins atop the island melt away, and Alyxian can no longer use the lair action associated with that prayer site. In addition, Alyxian takes 14 (4d6) psychic damage, regardless of his current form.

\end{multicols}
\end{DndMonster}


\begin{DndMonster}[float=t]{Chuul}
\DndMonsterType{Large aberration, chaotic evil}
\DndMonsterBasics[
armor-class = {16 (natural armor)},
hit-points = {\DndDice{11d10 + 33}},
speed = {30 ft., swim 30 ft.}
]
\DndMonsterAbilityScores[
 str = 19,
 dex = 10,
 con = 16,
 int = 5,
 wis = 11,
 cha = 5
]
\DndMonsterDetails[
skills = {Perception +4},
damage-immunities = {poison},
senses = {darkvision 60 ft., passive perception 14},
languages = {understands Deep Speech but can't speak},
challenge = 4,
]
\DndMonsterAction{Amphibious}
The chuul can breathe air and water.

\DndMonsterAction{Sense Magic}
The chuul senses magic within 120 feet of it at will. This trait otherwise works like the \textit{detect magic} spell but isn't itself magical.

\DndMonsterSection{Actions}
\DndMonsterAction{Multiattack}
The chuul makes two pincer attacks. If the chuul is grappling a creature, the chuul can also use its tentacles once.

\DndMonsterAction{Pincer}
\textsl{Melee Weapon Attack:} +6 to hit, reach 10 ft., one target. \textsl{Hit:} 11 (2d6 + 4) bludgeoning damage. The target is grappled (escape DC 14) if it is a Large or smaller creature and the chuul doesn't have two other creatures grappled.

\DndMonsterAction{Tentacles}
One creature grappled by the chuul must succeed on a DC 13 Constitution saving throw or be poisoned for 1 minute. Until this poison ends, the target is paralyzed. The target can repeat the saving throw at the end of each of its turns, ending the effect on itself on a success.

\end{DndMonster}



\begin{DndMonster}[float=t]{Cloaker}
\DndMonsterType{Large aberration, chaotic neutral}
\DndMonsterBasics[
armor-class = {14 (natural armor)},
hit-points = {\DndDice{12d10 + 12}},
speed = {10 ft., fly 40 ft.}
]
\DndMonsterAbilityScores[
 str = 17,
 dex = 15,
 con = 12,
 int = 13,
 wis = 12,
 cha = 14
]
\DndMonsterDetails[
skills = {Stealth +5},
senses = {darkvision 60 ft., passive perception 11},
languages = {Deep Speech, Undercommon},
challenge = 8,
]
\DndMonsterAction{Damage Transfer}
While attached to a creature, the cloaker takes only half the damage dealt to it (rounded down). and that creature takes the other half.

\DndMonsterAction{False Appearance}
While the cloaker remains motionless without its underside exposed, it is indistinguishable from a dark leather cloak.

\DndMonsterAction{Light Sensitivity}
While in bright light, the cloaker has disadvantage on attack rolls and Wisdom (Perception) checks that rely on sight.

\DndMonsterSection{Actions}
\DndMonsterAction{Multiattack}
The cloaker makes two attacks: one with its bite and one with its tail.

\DndMonsterAction{Bite}
\textsl{Melee Weapon Attack:} +6 to hit, reach 5 ft., one creature. \textsl{Hit:} 10 (2d6 + 3) piercing damage, and if the target is Large or smaller, the cloaker attaches to it. If the cloaker has advantage against the target, the cloaker attaches to the target's head, and the target is blinded and unable to breathe while the cloaker is attached. While attached, the cloaker can make this attack only against the target and has advantage on the attack roll. The cloaker can detach itself by spending 5 feet of its movement. A creature, including the target, can take its action to detach the cloaker by succeeding on a DC 16 Strength check.

\DndMonsterAction{Tail}
\textsl{Melee Weapon Attack:} +6 to hit, reach 10 ft., one creature. \textsl{Hit:} 7 (1d8 + 3) slashing damage.

\DndMonsterAction{Moan}
Each creature within 60 feet of the cloaker that can hear its moan and that isn't an aberration must succeed on a DC 13 Wisdom saving throw or become frightened until the end of the cloaker's next turn. If a creature's saving throw is successful, the creature is immune to the cloaker's moan for the next 24 hours.

\DndMonsterAction{Phantasms (Recharges after a Short or Long Rest)}
The cloaker magically creates three illusory duplicates of itself if it isn't in bright light. The duplicates move with it and mimic its actions, shifting position so as to make it impossible to track which cloaker is the real one. If the cloaker is ever in an area of bright light, the duplicates disappear.

Whenever any creature targets the cloaker with an attack or a harmful spell while a duplicate remains, that creature rolls randomly to determine whether it targets the cloaker or one of the duplicates. A creature is unaffected by this magical effect if it can't see or if it relies on senses other than sight.

A duplicate has the cloaker's AC and uses its saving throws. If an attack hits a duplicate, or if a duplicate fails a saving throw against an effect that deals damage, the duplicate disappears.

\end{DndMonster}



\begin{DndMonster}[float=t]{Corrupted Giant Shark}
\DndMonsterType{Huge aberration, unaligned}
\DndMonsterBasics[
armor-class = {13 (natural armor)},
hit-points = {\DndDice{11d12 + 55}},
speed = {0 ft., swim 50 ft.}
]
\DndMonsterAbilityScores[
 str = 23,
 dex = 11,
 con = 21,
 int = 1,
 wis = 10,
 cha = 5
]
\DndMonsterDetails[
skills = {Perception +4},
senses = {blindsight 60 ft., passive perception 14},
challenge = 9,
]
\DndMonsterAction{Psychic Maelstrom}
Any creature that starts its turn within 15 feet of the shark must succeed on a DC 17 Wisdom saving throw or take 11 (2d10) psychic damage.

\DndMonsterAction{Regeneration}
The shark regains 10 hit points at the start of its turn. If the shark takes radiant damage or suffers a critical hit, this trait doesn't function at the start of its next turn. The shark dies only if it starts its turn with 0 hit points and doesn't regenerate.

\DndMonsterAction{Water Breathing}
The shark can breathe only underwater.

\DndMonsterSection{Actions}
\DndMonsterAction{Bite}
\textsl{Melee Weapon Attack:} +10 to hit (with advantage if the target is a creature missing any hit points), reach 5 ft., one target. \textsl{Hit:} 22 (3d10 + 6) piercing damage, and if the target is a creature, it must succeed on a DC 17 Charisma saving throw or gain 1 level of exhaustion.

\end{DndMonster}



\begin{DndMonster}[float=t]{Death Embrace}
\DndMonsterType{Huge aberration, chaotic evil}
\DndMonsterBasics[
armor-class = {13 (natural armor)},
hit-points = {\DndDice{14d12 + 56}},
speed = {0 ft., swim 20 ft.}
]
\DndMonsterAbilityScores[
 str = 23,
 dex = 15,
 con = 19,
 int = 6,
 wis = 9,
 cha = 4
]
\DndMonsterDetails[
saving-throws = {Str +10, Wis +3},
senses = {blindsight 60 ft. (blind beyond this radius), passive perception 9},
challenge = 11,
]
\DndMonsterAction{Petrifying Tendrils}
Any creature that starts its turn in the death embrace's space must make a DC 16 Constitution saving throw. On a failed saving throw, the creature is restrained. A creature restrained in this way must repeat the saving throw at the end of its next turn, becoming petrified on a failed saving throw or ending the effect on a successful one. The petrified condition lasts until the effect is ended by a greater restoration spell or similar magic.

\DndMonsterAction{Reel}
At the start of each of its turns, the death embrace can pull each creature it is grappling up to 20 feet toward it (no action required).

\DndMonsterAction{Water Breathing}
The death embrace can breathe only underwater.

\DndMonsterSection{Actions}
\DndMonsterAction{Multiattack}
The death embrace makes two Tentacle attacks.

\DndMonsterAction{Tentacle}
\textsl{Melee Weapon Attack:} +10 to hit, reach 60 ft., one target. \textsl{Hit:} 13 (2d6 + 6) piercing damage plus 11 (2d10) psychic damage. If the target is Large or smaller, it is also grappled (escape DC 16). The death embrace has six tentacles, each of which can grapple one target.

\DndMonsterSection{Reactions}
\DndMonsterAction{Body Shield}
When the death embrace is hit by an attack, one creature the death embrace is grappling (chosen by the death embrace) takes the damage instead.

\end{DndMonster}



\begin{DndMonster}[float=t]{Drow Elite Warrior}
\DndMonsterType{Medium humanoid (elf), neutral evil}
\DndMonsterBasics[
armor-class = {18 (studded leather armor)},
hit-points = {\DndDice{11d8 + 22}},
speed = {30 ft.}
]
\DndMonsterAbilityScores[
 str = 13,
 dex = 18,
 con = 14,
 int = 11,
 wis = 13,
 cha = 12
]
\DndMonsterDetails[
saving-throws = {Dex +7, Con +5, Wis +4},
skills = {Perception +4, Stealth +10},
senses = {darkvision 120 ft., passive perception 14},
languages = {Elvish, Undercommon},
challenge = 5,
]
\DndMonsterAction{Fey Ancestry}
The drow has advantage on saving throws against being charmed, and magic can't put the drow to sleep.

\DndMonsterAction{Sunlight Sensitivity}
While in sunlight, the drow has disadvantage on attack rolls, as well as on Wisdom (Perception) checks that rely on sight.

\DndMonsterAction{Innate Spellcasting}
The drow's spellcasting ability is Charisma (spell save DC 12). It can innately cast the following spells, requiring no material components:
\begin{DndMonsterSpells}
  \DndInnateSpellLevel{dancing lights}
  \DndInnateSpellLevel[1e]{darkness, faerie fire, levitate}
\end{DndMonsterSpells}
\DndMonsterSection{Actions}
\DndMonsterAction{Multiattack}
The drow makes two shortsword attacks.

\DndMonsterAction{Shortsword}
\textsl{Melee Weapon Attack:} +7 to hit, reach 5 ft., one target. \textsl{Hit:} 7 (1d6 + 4) piercing damage plus 10 (3d6) poison damage.

\DndMonsterAction{Hand Crossbow}
\textsl{Ranged Weapon Attack:} +7 to hit, range 30/120 ft., one target. \textsl{Hit:} 7 (1d6 + 4) piercing damage, and the target must succeed on a DC 13 Constitution saving throw or be poisoned for 1 hour. If the saving throw fails by 5 or more, the target is also unconscious while poisoned in this way. The target wakes up if it takes damage or if another creature takes an action to shake it awake.

\DndMonsterSection{Reactions}
\DndMonsterAction{Parry}
The drow adds 3 to its AC against one melee attack that would hit it. To do so, the drow must see the attacker and be wielding a melee weapon.

\end{DndMonster}



\begin{DndMonster}[float=t]{Ghost}
\DndMonsterType{Medium undead, any alignment}
\DndMonsterBasics[
armor-class = {11},
hit-points = {\DndDice{10d8}},
speed = {0 ft., fly 40 ft. \textit{(hover)}}
]
\DndMonsterAbilityScores[
 str = 7,
 dex = 13,
 con = 10,
 int = 10,
 wis = 12,
 cha = 17
]
\DndMonsterDetails[
damage-resistances = {acid, fire, lightning, thunder; bludgeoning, piercing, slashing from nonmagical attacks},
damage-immunities = {cold, necrotic, poison},
senses = {darkvision 60 ft., passive perception 11},
languages = {any languages it knew in life},
challenge = 4,
]
\DndMonsterAction{Ethereal Sight}
The ghost can see 60 feet into the Ethereal Plane when it is on the Material Plane, and vice versa.

\DndMonsterAction{Incorporeal Movement}
The ghost can move through other creatures and objects as if they were difficult terrain. It takes 5 (1d10) force damage if it ends its turn inside an object.

\DndMonsterSection{Actions}
\DndMonsterAction{Withering Touch}
\textsl{Melee Weapon Attack:} +5 to hit, reach 5 ft., one target. \textsl{Hit:} 17 (4d6 + 3) necrotic damage.

\DndMonsterAction{Etherealness}
The ghost enters the Ethereal Plane from the Material Plane, or vice versa. It is visible on the Material Plane while it is in the Border Ethereal, and vice versa, yet it can't affect or be affected by anything on the other plane.

\DndMonsterAction{Horrifying Visage}
Each non-undead creature within 60 feet of the ghost that can see it must succeed on a DC 13 Wisdom saving throw or be frightened for 1 minute. If the save fails by 5 or more, the target also ages 1d4 × 10 years. A frightened target can repeat the saving throw at the end of each of its turns, ending the frightened condition on itself on a success. If a target's saving throw is successful or the effect ends for it, the target is immune to this ghost's Horrifying Visage for the next 24 hours. The aging effect can be reversed with a  \textit{greater restoration} spell, but only within 24 hours of it occurring.

\DndMonsterAction{Possession (Recharge 6)}
One humanoid that the ghost can see within 5 feet of it must succeed on a DC 13 Charisma saving throw or be possessed by the ghost; the ghost then disappears, and the target is incapacitated and loses control of its body. The ghost now controls the body but doesn't deprive the target of awareness. The ghost can't be targeted by any attack, spell, or other effect, except ones that turn undead, and it retains its alignment, Intelligence, Wisdom, Charisma, and immunity to being charmed and frightened. It otherwise uses the possessed target's statistics, but doesn't gain access to the target's knowledge, class features, or proficiencies.

The possession lasts until the body drops to 0 hit points, the ghost ends it as a bonus action, or the ghost is turned or forced out by an effect like the \textit{dispel evil and good} spell. When the possession ends, the ghost reappears in an unoccupied space within 5 feet of the body. The target is immune to this ghost's Possession for 24 hours after succeeding on the saving throw or after the possession ends.

\end{DndMonster}



\begin{DndMonster}[float=t]{Horizonback Tortoise}
\DndMonsterType{Gargantuan monstrosity, unaligned}
\DndMonsterBasics[
armor-class = {17 (natural armor) (22 while in its shell)},
hit-points = {\DndDice{13d20 + 91}},
speed = {20 ft.}
]
\DndMonsterAbilityScores[
 str = 28,
 dex = 3,
 con = 25,
 int = 4,
 wis = 10,
 cha = 5
]
\DndMonsterDetails[
saving-throws = {Str +12, Con +10},
damage-immunities = {poison},
senses = {darkvision 60 ft, passive perception 10},
languages = {understands Goblin but can't speak},
challenge = 8,
]
\DndMonsterAction{Amphibious}
The tortoise can breathe air and water.

\DndMonsterAction{Massive Frame}
The tortoise can carry up to 20,000 pounds of weight atop its shell, but moves at half speed if the weight exceeds 10,000 pounds. Medium or smaller creatures can move underneath the tortoise while it's not prone.

Any creature under the tortoise when it falls prone is grappled (escape DC 18). Until the grapple ends, the creature is prone and restrained.

\DndMonsterSection{Actions}
\DndMonsterAction{Bite}
\textsl{Melee Weapon Attack:} +12 to hit, reach 10 ft., one target. \textsl{Hit:} 28 (3d12 + 9) bludgeoning damage.

\DndMonsterAction{Shell Defense (Recharge 4\textendash 6)}
The tortoise withdraws into its shell, falls prone, and gains a +5 bonus to AC. While the tortoise is in its shell, its speed is 0 and can't increase. The tortoise can emerge from its shell as an action, whereupon it is no longer prone.

\end{DndMonster}



\begin{DndMonster}[float=t]{Light Devourer}
\DndMonsterType{Medium aberration, unaligned}
\DndMonsterBasics[
armor-class = {14},
hit-points = {\DndDice{14d8 + 56}},
speed = {0 ft., swim 30 ft.}
]
\DndMonsterAbilityScores[
 str = 18,
 dex = 18,
 con = 18,
 int = 2,
 wis = 10,
 cha = 3
]
\DndMonsterDetails[
skills = {Perception +3, Stealth +10},
damage-immunities = {radiant},
senses = {darkvision 120 ft., passive perception 13},
challenge = 6,
]
\DndMonsterAction{Light Absorption}
If the light devourer spends any part of its turn in an area of bright light or is subjected to radiant damage, it radiates dim light in a 10-foot radius for 1 hour.

\DndMonsterAction{Water Breathing}
The light devourer can breathe only underwater.

\DndMonsterSection{Actions}
\DndMonsterAction{Multiattack}
The light devourer makes two Bite attacks.

\DndMonsterAction{Bite}
\textsl{Melee Weapon Attack:} +7 to hit, reach 5 ft., one target. \textsl{Hit:} 21 (5d6 + 4) piercing damage, or radiant damage if the light devourer is radiating light (see Light Absorption).

\DndMonsterSection{Bonus Actions}
\DndMonsterAction{Radiant Discharge}
If the light devourer is radiating light (see Light Absorption), it releases the light stored in its body in a 20-foot-radius sphere centered on itself, then ceases to radiate light. Each creature in the sphere must make a DC 15 Constitution saving throw, taking 24 (7d6) radiant damage on a failed saving throw, or half as much damage on a successful one.

\end{DndMonster}


\begin{DndMonster}[float=t]{Moorbounder}
\DndMonsterType{Large beast, unaligned}
\DndMonsterBasics[
armor-class = {13 (natural armor)},
hit-points = {\DndDice{4d10 + 8}},
speed = {70 ft.}
]
\DndMonsterAbilityScores[
 str = 18,
 dex = 14,
 con = 14,
 int = 2,
 wis = 13,
 cha = 5
]
\DndMonsterDetails[
senses = {darkvision 60 ft., passive perception 11},
challenge = 1,
]
\DndMonsterAction{Standing Leap}
The moorbounder's long jump is up to 40 feet and its high jump is up to 20 feet, with or without a running start.

\DndMonsterSection{Actions}
\DndMonsterAction{Claws}
\textsl{Melee Weapon Attack:} +6 to hit, reach 5 ft., one target. \textsl{Hit:} 14 (4d4 + 4) slashing damage.

\end{DndMonster}



\begin{DndMonster}[float=t]{Owlbear}
  \DndMonsterType{Large monstrosity (Brute), unaligned}
  \DndMonsterBasics[
  armor-class = {13 (natural armor)},
  hit-points = {\DndDice{11d10 + 22}},
  speed = {40 ft.}
  ]
  \DndMonsterAbilityScores[
   str = 21,
   dex = 12,
   con = 15,
   int = 5,
   wis = 12,
   cha = 6
  ]
  \DndMonsterDetails[
  skills = {Perception +3},
  condition-immunities = {frightened},
  senses = {darkvision 60 ft., passive perception 13},
  challenge = 3,
  ]
  \DndMonsterAction{Deadly Leap}
  The owlbear's long jump is up to 30 feet and their high jump is 15 feet, with or without a running start. If the owlbear leaps at least 20 feet toward a target and then hits them with a Bite or Claw attack, the target must succeed on a DC 15 Strength saving throw or be knocked prone.
  
  \DndMonsterAction{Keen Sight and Smell}
  The owlbear has advantage on Wisdom (Perception) checks that rely on sight or smell.
  
  \DndMonsterSection{Actions}
  \DndMonsterAction{Multiattack}
  The owlbear makes one Bite and one Claw attack.
  
  \DndMonsterAction{Bite}
  \textsl{Melee Weapon Attack:} 7 to hit, reach 10 ft., one target. \textsl{Hit:} 11 (1d12 + 5) piercing damage, and the owlbear can move the target 5 feet in any direction.
  
  \DndMonsterAction{Claw}
  \textsl{Melee Weapon Attack:} 7 to hit, reach 5 ft., one target. \textsl{Hit:} 12 (2d6 + 5) slashing damage.
  
  \DndMonsterAction{Bear Hug (Recharge 4\textendash 6)}
  The owlbear attempts to grab and crush a creature they can see within 5 feet of them. The target must make a DC 15 Dexterity saving throw. On a failed save, the target takes 22 (4d10) bludgeoning damage and is grappled (escape DC 15). On a successful save, the target takes half as much damage and is not grappled. Until this grapple ends, the target is restrained and takes 5 (1d10) bludgeoning damage at the start of each of their turns. The grapple ends if the owlbear uses Bear Hug on another target or makes a Claw attack against another target.
  
  \DndMonsterSection{Reactions}
  \DndMonsterAction{Hulking Rush}
  When the owlbear takes damage, they can move up to half their speed without provoking opportunity attacks.
  
  \end{DndMonster}
  

\clearpage

\begin{DndMonster}[float=t]{Perigee}
\DndMonsterType{Medium celestial, lawful good}
\DndMonsterBasics[
armor-class = {17 (natural armor)},
hit-points = {\DndDice{16d8 + 64}},
speed = {30 ft., fly 90 ft.}
]
\DndMonsterAbilityScores[
 str = 18,
 dex = 18,
 con = 18,
 int = 17,
 wis = 20,
 cha = 20
]
\DndMonsterDetails[
saving-throws = {Wis +9, Cha +9},
skills = {Insight +9, Perception +9},
damage-resistances = {radiant; bludgeoning, piercing, slashing from nonmagical attacks},
senses = {darkvision 120 ft., passive perception 19},
languages = {all, telepathy 120 ft.},
challenge = 10,
]
\DndMonsterAction{Angelic Weapons}
Perigee's weapon attacks are magical. When Perigee hits with any weapon, the weapon deals an extra 4d8 radiant damage (included in the attack).

\DndMonsterAction{Magic Resistance}
Perigee has advantage on saving throws against spells and other magical effects.

\DndMonsterAction{Innate Spellcasting}
Perigee's spellcasting ability is Charisma (spell save DC 17). Perigee can innately cast the following spells, requiring only verbal components:
\begin{DndMonsterSpells}
  \DndInnateSpellLevel{detect evil and good}
  \DndInnateSpellLevel[1e]{commune, raise dead}
\end{DndMonsterSpells}
\DndMonsterSection{Actions}
\DndMonsterAction{Multiattack}
Perigee makes two melee attacks.

\DndMonsterAction{Mace}
\textsl{Melee Weapon Attack:} +8 to hit, reach 5 ft., one target. \textsl{Hit:} 7 (1d6 + 4) bludgeoning damage plus 18 (4d8) radiant damage.

\DndMonsterAction{Healing Touch (3/Day)}
Perigee touches another creature. The target magically regains 20 (4d8 + 2) hit points and is freed from any curse, disease, poison, blindness, or deafness.

\DndMonsterAction{Change Shape}
Perigee magically polymorphs into a humanoid or beast that has a challenge rating equal to or less than its own, or back into its true form. It reverts to its true form if it dies. Any equipment it is wearing or carrying is absorbed or borne by the new form (Perigee's choice).

In a new form, Perigee retains its game statistics and ability to speak, but its AC, movement modes, Strength, Dexterity, and special senses are replaced by those of the new form, and it gains any statistics and capabilities (except class features, legendary actions, and lair actions) that the new form has but that it lacks.

\end{DndMonster}



\begin{DndMonster}[float=t]{Revenant}
\DndMonsterType{Medium undead, neutral}
\DndMonsterBasics[
armor-class = {13 (\textit{leather armor})},
hit-points = {\DndDice{16d8 + 64}},
speed = {30 ft.}
]
\DndMonsterAbilityScores[
 str = 18,
 dex = 14,
 con = 18,
 int = 13,
 wis = 16,
 cha = 18
]
\DndMonsterDetails[
saving-throws = {Str +7, Con +7, Wis +6, Cha +7},
damage-resistances = {necrotic, psychic},
damage-immunities = {poison},
senses = {darkvision 60 ft., passive perception 13},
languages = {the languages it knew in life},
challenge = 5,
]
\DndMonsterAction{Regeneration}
The revenant regains 10 hit points at the start of its turn. If the revenant takes fire or radiant damage, this trait doesn't function at the start of the revenant's next turn. The revenant's body is destroyed only if it starts its turn with 0 hit points and doesn't regenerate.

\DndMonsterAction{Rejuvenation}
When the revenant's body is destroyed, its soul lingers. After 24 hours, the soul inhabits and animates another humanoid corpse on the same plane of existence and regains all its hit points. While the soul is bodiless, a \textit{wish} spell can be used to force the soul to go to the afterlife and not return.

\DndMonsterAction{Turn Immunity}
The revenant is immune to effects that turn undead.

\DndMonsterAction{Vengeful Tracker}
The revenant knows the distance to and direction of any creature against which it seeks revenge, even if the creature and the revenant are on different planes of existence. If the creature being tracked by the revenant dies, the revenant knows.

\DndMonsterSection{Actions}
\DndMonsterAction{Multiattack}
The revenant makes two fist attacks.

\DndMonsterAction{Fist}
\textsl{Melee Weapon Attack:} +7 to hit, reach 5 ft., one target. \textsl{Hit:} 11 (2d6 + 4) bludgeoning damage. If the target is a creature against which the revenant has sworn vengeance, the target takes an extra 14 (4d6) bludgeoning damage. Instead of dealing damage, the revenant can grapple the target (escape DC 14) provided the target is Large or smaller.

\DndMonsterAction{Vengeful Glare}
The revenant targets one creature it can see within 30 feet of it and against which it has sworn vengeance. The target must make a DC 15 Wisdom saving throw. On a failure, the target is paralyzed until the revenant deals damage to it, or until the end of the revenant's next turn. When the paralysis ends, the target is frightened of the revenant for 1 minute. The frightened target can repeat the saving throw at the end of each of its turns, with disadvantage if it can see the revenant, ending the frightened condition on itself on a success.

\end{DndMonster}



\begin{DndMonster}[float=t]{Scout}
\DndMonsterType{Medium humanoid (any race), any alignment}
\DndMonsterBasics[
armor-class = {13 (\textit{leather armor})},
hit-points = {\DndDice{3d8 + 3}},
speed = {30 ft.}
]
\DndMonsterAbilityScores[
 str = 11,
 dex = 14,
 con = 12,
 int = 11,
 wis = 13,
 cha = 11
]
\DndMonsterDetails[
skills = {Nature +4, Perception +5, Stealth +6, Survival +5},
languages = {any one language (usually Common)},
challenge = 1/2,
]
\DndMonsterAction{Keen Hearing and Sight}
The scout has advantage on Wisdom (Perception) checks that rely on hearing or sight.

\DndMonsterSection{Actions}
\DndMonsterAction{Multiattack}
The scout makes two melee attacks or two ranged attacks.

\DndMonsterAction{Shortsword}
\textsl{Melee Weapon Attack:} +4 to hit, reach 5 ft., one target. \textsl{Hit:} 5 (1d6 + 2) piercing damage.

\DndMonsterAction{Longbow}
\textsl{Ranged Weapon Attack:} +4 to hit, ranged 150/600 ft., one target. \textsl{Hit:} 6 (1d8 + 2) piercing damage.

\end{DndMonster}



\begin{DndMonster}[float=t]{Scuttling Serpentmaw}
\DndMonsterType{Small aberration, unaligned}
\DndMonsterBasics[
armor-class = {17 (natural armor)},
hit-points = {\DndDice{10d6 + 30}},
speed = {30 ft.}
]
\DndMonsterAbilityScores[
 str = 16,
 dex = 14,
 con = 16,
 int = 3,
 wis = 11,
 cha = 3
]
\DndMonsterDetails[
skills = {Stealth +6},
senses = {blindsight 60 ft. (blind beyond this radius), passive perception 10},
challenge = 4,
]
\DndMonsterAction{Amphibious}
The serpentmaw can breathe air and water.

\DndMonsterAction{Pack Tactics}
The serpentmaw has advantage on an attack roll against a creature if at least one of the serpentmaw's allies is within 5 feet of the creature and the ally isn't incapacitated.

\DndMonsterSection{Actions}
\DndMonsterAction{Multiattack}
The serpentmaw makes one Bite attack and two Claw attacks.

\DndMonsterAction{Bite}
\textsl{Melee Weapon Attack:} +5 to hit, reach 15 ft., one target. \textsl{Hit:} 10 (2d6 + 3) piercing damage, or 17 (4d6 + 3) piercing damage if the serpentmaw had advantage on the attack roll.

\DndMonsterAction{Claw}
\textsl{Melee Weapon Attack:} +5 to hit, reach 5 ft., one target. \textsl{Hit:} 7 (1d8 + 3) slashing damage.

\end{DndMonster}



\begin{DndMonster}[float=t]{Skeletal Bloodfin}
\DndMonsterType{Large undead, unaligned}
\DndMonsterBasics[
armor-class = {16 (natural armor)},
hit-points = {\DndDice{11d10 + 33}},
speed = {5 ft., swim 50 ft.}
]
\DndMonsterAbilityScores[
 str = 20,
 dex = 14,
 con = 16,
 int = 2,
 wis = 10,
 cha = 3
]
\DndMonsterDetails[
skills = {Perception +4},
damage-immunities = {poison},
damage-resistances = {piercing, slashing},
condition-immunities = {charmed, deafened, frightened, paralyzed, poisoned},
senses = {blindsight 100 ft. (blind beyond this radius), passive perception 14},
challenge = 9,
]

\DndMonsterAction{Death Burst}
When the bloodfin dies, it explodes in a cloud of toxic blood. Each creature in a 10-foot-radius sphere centered on the exploding bloodfin, including any creature swallowed by the bloodfin, must succeed on a DC 15 Constitution saving throw or take 10 (3d6) poison damage. A creature inside the bloodfin when it explodes falls prone in the space formerly occupied by the bloodfin and is no longer blinded or restrained by it.

\DndMonsterAction{Unusual Nature}
The skeletal bloodfin requires no food, drink, air, or sleep.

\DndMonsterSection{Actions}
\DndMonsterAction{Multiattack}
The bloodfin makes one Bite attack and one Tail attack.

\DndMonsterAction{Bite}
\textsl{Melee Weapon Attack:} +9 to hit (with advantage if the target is a creature missing any hit points), reach 5 ft., one target. \textsl{Hit:} 14 (2d8 + 5) piercing damage, and if the target is a creature, it is grappled (escape DC 15). Until this grapple ends, the target is restrained, and the bloodfin can't bite another target.

\DndMonsterAction{Tail}
\textsl{Melee Weapon Attack:} +9 to hit (with advantage if the target is a creature missing any hit points), reach 10 ft., one target. \textsl{Hit:} 19 (4d6 + 5) bludgeoning damage.

\DndMonsterSection{Bonus Actions}
\DndMonsterAction{Swallow}
\textsl{Melee Weapon Attack:} +9 to hit, reach 5 ft., one Medium or smaller creature the bloodfin is grappling. \textsl{Hit:} The bloodfin swallows the target. The swallowed creature is no longer grappled but is blinded and restrained. It has total cover against attacks and other effects outside the bloodfin, it takes 14 (4d6) necrotic damage at the start of each of its turns, and the bloodfin regains hit points equal to the necrotic damage dealt. A bloodfin can have only one creature swallowed at a time.

If the bloodfin takes 30 damage or more on a single turn from the swallowed creature, the bloodfin must succeed on a DC 15 Constitution saving throw at the end of that turn or regurgitate the creature, which falls prone in a space within 5 feet of the bloodfin and is no longer blinded or restrained by it.

\end{DndMonster}



\begin{DndMonster}[float=t]{Slithering Bloodfin}
\DndMonsterType{Large aberration, unaligned}
\DndMonsterBasics[
armor-class = {16 (natural armor)},
hit-points = {\DndDice{11d10 + 33}},
speed = {5 ft., swim 50 ft.}
]
\DndMonsterAbilityScores[
 str = 20,
 dex = 14,
 con = 16,
 int = 2,
 wis = 10,
 cha = 3
]
\DndMonsterDetails[
skills = {Perception +4},
damage-immunities = {poison},
senses = {blindsight 100 ft. (blind beyond this radius), passive perception 14},
challenge = 9,
]
\DndMonsterAction{Death Burst}
When the bloodfin dies, it explodes in a cloud of toxic blood. Each creature in a 10-foot-radius sphere centered on the exploding bloodfin, including any creature swallowed by the bloodfin, must succeed on a DC 15 Constitution saving throw or take 10 (3d6) poison damage. A creature inside the bloodfin when it explodes falls prone in the space formerly occupied by the bloodfin and is no longer blinded or restrained by it.

\DndMonsterAction{Water Breathing}
The bloodfin can breathe only underwater.

\DndMonsterSection{Actions}
\DndMonsterAction{Multiattack}
The bloodfin makes one Bite attack and one Tail attack.

\DndMonsterAction{Bite}
\textsl{Melee Weapon Attack:} +9 to hit (with advantage if the target is a creature missing any hit points), reach 5 ft., one target. \textsl{Hit:} 14 (2d8 + 5) piercing damage, and if the target is a creature, it is grappled (escape DC 15). Until this grapple ends, the target is restrained, and the bloodfin can't bite another target.

\DndMonsterAction{Tail}
\textsl{Melee Weapon Attack:} +9 to hit (with advantage if the target is a creature missing any hit points), reach 10 ft., one target. \textsl{Hit:} 19 (4d6 + 5) bludgeoning damage.

\DndMonsterSection{Bonus Actions}
\DndMonsterAction{Swallow}
\textsl{Melee Weapon Attack:} +9 to hit, reach 5 ft., one Medium or smaller creature the bloodfin is grappling. \textsl{Hit:} The bloodfin swallows the target. The swallowed creature is no longer grappled but is blinded and restrained. It has total cover against attacks and other effects outside the bloodfin, it takes 14 (4d6) necrotic damage at the start of each of its turns, and the bloodfin regains hit points equal to the necrotic damage dealt. A bloodfin can have only one creature swallowed at a time.

If the bloodfin takes 30 damage or more on a single turn from the swallowed creature, the bloodfin must succeed on a DC 15 Constitution saving throw at the end of that turn or regurgitate the creature, which falls prone in a space within 5 feet of the bloodfin and is no longer blinded or restrained by it.

\end{DndMonster}



\begin{DndMonster}[float=t]{Swarm of Sorrowfish}
\DndMonsterType{Medium aberration, unaligned}
\DndMonsterBasics[
armor-class = {14},
hit-points = {\DndDice{16d8 + 32}},
speed = {0 ft., swim 30 ft.}
]
\DndMonsterAbilityScores[
 str = 16,
 dex = 19,
 con = 14,
 int = 1,
 wis = 11,
 cha = 3
]
\DndMonsterDetails[
damage-resistances = {bludgeoning, piercing, slashing},
damage-immunities = {psychic},
senses = {blindsight 60 ft. (blind beyond this radius), passive perception 10},
challenge = 6,
]
\DndMonsterAction{Swarm}
The swarm can occupy another creature's space and vice versa, and the swarm can move through an opening as narrow as 1 foot wide. The swarm can't regain hit points or gain temporary hit points.

\DndMonsterAction{Virulent Sorrow}
Each time the swarm takes damage, any creature within 5 feet of it must make a DC 13 Wisdom saving throw. On a failed saving throw, the creature suffers the following effects until the end of its next turn: it has disadvantage on its attack rolls, it can't take reactions, and its speed is halved.

\DndMonsterAction{Water Breathing}
The swarm can breathe only underwater.

\DndMonsterSection{Actions}
\DndMonsterAction{Swarm of Bites}
\textsl{Melee Weapon Attack:} +8 to hit, reach 0 ft., one target in the swarm's space. \textsl{Hit:} 21 (6d6) piercing damage, or 10 (3d6) piercing damage if the swarm has half its hit points or fewer.

\DndMonsterAction{Desolate Drain (Recharge 5\textendash 6)}
Each creature in the swarm's space must make a DC 13 Wisdom saving throw. On a failed saving throw, a creature takes 24 (7d6) psychic damage and is stunned until the end of its next turn. The effect ends on a creature if the swarm moves out of the creature's space or if another creature within 5 feet of the swarm uses an action to make a DC 14 Strength check, pulling the affected creature out of the swarm on a successful check. On a successful saving throw, a creature takes half as much damage and suffers no other effects.

\end{DndMonster}



\begin{DndMonster}[float=t]{Sword Wraith Commander}
\DndMonsterType{Medium undead, lawful evil}
\DndMonsterBasics[
armor-class = {18 (\textit{breastplate})},
hit-points = {\DndDice{15d8 + 60}},
speed = {30 ft.}
]
\DndMonsterAbilityScores[
 str = 18,
 dex = 14,
 con = 18,
 int = 11,
 wis = 12,
 cha = 14
]
\DndMonsterDetails[
skills = {Perception +4},
damage-resistances = {necrotic; bludgeoning, piercing, slashing from nonmagical attacks},
damage-immunities = {poison},
senses = {darkvision 60 ft., passive perception 14},
languages = {the languages it knew in life},
challenge = 8,
]
\DndMonsterAction{Turning Defiance}
The commander and any other sword wraiths within 30 feet of it have advantage on saving throws against effects that turn Undead.

\DndMonsterAction{Unusual Nature}
The commander doesn't require air, food, drink, or sleep.

\DndMonsterSection{Actions}
\DndMonsterAction{Multiattack}
The commander makes two Longsword or Longbow attacks.

\DndMonsterAction{Longsword}
\textsl{Melee Weapon Attack:} +7 to hit, reach 5 ft., one target. \textsl{Hit:} 8 (1d8 + 4) slashing damage, or 9 (1d10 + 4) slashing damage if used with two hands.

\DndMonsterAction{Longbow}
\textsl{Ranged Weapon Attack:} +5 to hit, range 150/600 ft., one target. \textsl{Hit:} 6 (1d8 + 2) piercing damage.

\DndMonsterAction{Call to Honor (1/Day)}
If the commander has taken damage during this combat, it gives itself advantage on attack rolls until the end of its next turn, and 1d4 + 1 sword wraith warriors appear in unoccupied spaces within 30 feet of it. The warriors last until they drop to 0 hit points, and they take their turns immediately after the commander's turn on the same initiative count.

\DndMonsterSection{Bonus Actions}
\DndMonsterAction{Martial Fury}
The commander makes one Longsword or Longbow attack, which deals an extra 9 (2d8) necrotic damage on a hit, and attack rolls against it have advantage until the start of its next turn.

\end{DndMonster}



\begin{DndMonster}[float=t]{Sword Wraith Warrior}
\DndMonsterType{Medium undead, lawful evil}
\DndMonsterBasics[
armor-class = {16 (\textit{chain shirt})},
hit-points = {\DndDice{6d8 + 18}},
speed = {30 ft.}
]
\DndMonsterAbilityScores[
 str = 18,
 dex = 12,
 con = 17,
 int = 6,
 wis = 9,
 cha = 10
]
\DndMonsterDetails[
damage-resistances = {necrotic; bludgeoning, piercing, slashing from nonmagical attacks},
damage-immunities = {poison},
senses = {darkvision 60 ft., passive perception 9},
languages = {the languages it knew in life},
challenge = 3,
]
\DndMonsterAction{Unusual Nature}
The warrior doesn't require air, food, drink, or sleep.

\DndMonsterSection{Actions}
\DndMonsterAction{Battleaxe}
\textsl{Melee Weapon Attack:} +6 to hit, reach 5 ft., one target. \textsl{Hit:} 8 (1d8 + 4) slashing damage, or 9 (1d10 + 4) slashing damage if used with two hands.

\DndMonsterAction{Longbow}
\textsl{Ranged Weapon Attack:} +3 to hit, range 150/600 ft., one target. \textsl{Hit:} 5 (1d8 + 1) piercing damage.

\DndMonsterSection{Bonus Actions}
\DndMonsterAction{Martial Fury}
The warrior makes one Battleaxe or Longbow attack, and attack rolls against it have advantage until the start of its next turn.

\end{DndMonster}



\begin{DndMonster}[float=t]{Veteran}
\DndMonsterType{Medium humanoid (any race), any alignment}
\DndMonsterBasics[
armor-class = {17 (\textit{splint armor})},
hit-points = {\DndDice{9d8 + 18}},
speed = {30 ft.}
]
\DndMonsterAbilityScores[
 str = 16,
 dex = 13,
 con = 14,
 int = 10,
 wis = 11,
 cha = 10
]
\DndMonsterDetails[
skills = {Athletics +5, Perception +2},
languages = {any one language (usually Common)},
challenge = 3,
]
\DndMonsterSection{Actions}
\DndMonsterAction{Multiattack}
The veteran makes two longsword attacks. If it has a shortsword drawn, it can also make a shortsword attack.

\DndMonsterAction{Longsword}
\textsl{Melee Weapon Attack:} +5 to hit, reach 5 ft., one target. \textsl{Hit:} 7 (1d8 + 3) slashing damage, or 8 (1d10 + 3) slashing damage if used with two hands.

\DndMonsterAction{Shortsword}
\textsl{Melee Weapon Attack:} +5 to hit, reach 5 ft., one target. \textsl{Hit:} 6 (1d6 + 3) piercing damage.

\DndMonsterAction{Heavy Crossbow}
\textsl{Ranged Weapon Attack:} +3 to hit, range 100/400 ft., one target. \textsl{Hit:} 6 (1d10 + 1) piercing damage.

\end{DndMonster}



\begin{DndMonster}[float=t]{Will-o'-Wisp}
\DndMonsterType{Tiny undead, chaotic evil}
\DndMonsterBasics[
armor-class = {19},
hit-points = {\DndDice{9d4}},
speed = {0 ft., fly 50 ft. \textit{(hover)}}
]
\DndMonsterAbilityScores[
 str = 1,
 dex = 28,
 con = 10,
 int = 13,
 wis = 14,
 cha = 11
]
\DndMonsterDetails[
damage-resistances = {acid, cold, fire, necrotic, thunder; bludgeoning, piercing, slashing from nonmagical attacks},
damage-immunities = {lightning, poison},
senses = {darkvision 120 ft., passive perception 12},
languages = {the languages it knew in life},
challenge = 2,
]
\DndMonsterAction{Consume Life}
As a bonus action, the will-o'-wisp can target one creature it can see within 5 feet of it that has 0 hit points and is still alive. The target must succeed on a DC 10 Constitution saving throw against this magic or die. If the target dies, the will-o'-wisp regains 10 (3d6) hit points.

\DndMonsterAction{Ephemeral}
The will-o'-wisp can't wear or carry anything.

\DndMonsterAction{Incorporeal Movement}
The will-o'-wisp can move through other creatures and objects as if they were difficult terrain. It takes 5 (1d10) force damage if it ends its turn inside an object.

\DndMonsterAction{Variable Illumination}
The will-o'-wisp sheds bright light in a 5 to 20-foot radius and dim light for an additional number of ft. equal to the chosen radius. The will-o'-wisp can alter the radius as a bonus action.

\DndMonsterSection{Actions}
\DndMonsterAction{Shock}
\textsl{Melee Spell Attack:} +4 to hit, reach 5 ft., one creature. \textsl{Hit:} 9 (2d8) lightning damage.

\DndMonsterAction{Invisibility}
The will-o'-wisp and its light magically become invisible until it attacks or uses its Consume Life, or until its concentration ends (as if concentration on a spell).

\end{DndMonster}


\end{document}
